% ---
% file: docs/latex/Fisica/formulario_fisica.tex
% description: Compendio exhaustivo de fórmulas de física desde nivel básico hasta nivel avanzado/universitario.
% type: doc/manual
% version: 1.0.0
% date: 2026-08-24
% covers:
%   - src/data/symbols.py
% relations:
%   - docs/mates/formulario_mates.md
%   - docs/ARCHITECTURE.md
% keywords:
%   - compendio-fisica
%   - formulas-fisica
%   - mecanica-clasica
%   - termodinamica
%   - electromagnetismo
%   - mecanica-cuantica
% ---

\documentclass[11pt,a4paper]{article}
\usepackage[utf8]{inputenc}
\usepackage[T1]{fontenc}
\usepackage[spanish,es-nodecimaldot,es-noquoting]{babel}
\usepackage{amsmath,amssymb,amsfonts}
\usepackage{esint}
\usepackage{array,tabularx}

% Configuración de márgenes y espaciado estándar
\setlength{\textwidth}{16.5cm}
\setlength{\textheight}{24cm}
\setlength{\oddsidemargin}{-0.25cm}
\setlength{\evensidemargin}{-0.25cm}
\setlength{\topmargin}{-1cm}
\setlength{\parindent}{0pt}
\setlength{\parskip}{5pt}
\setlength{\emergencystretch}{3em}

\begin{document}

\section*{Compendio I: Mecánica Clásica, de Fluidos y Analítica}
\addcontentsline{toc}{section}{Compendio I: Mecánica Clásica, de Fluidos y Analítica}

\textit{(Nivel Básico a Universitario)}

\section*{1. Cinemática y Dinámica Newtoniana}
\addcontentsline{toc}{section}{1. Cinemática y Dinámica Newtoniana}

\subsection*{Nivel Básico}
\addcontentsline{toc}{subsection}{Nivel Básico}

\subsubsection*{1.1 Movimiento Rectilíneo Uniforme (MRU)}
\addcontentsline{toc}{subsubsection}{1.1 Movimiento Rectilíneo Uniforme (MRU)}

\begin{itemize}
  \item \textbf{Posición en función del tiempo:}
  \[
    x(t) = x_0 + v \cdot t
  \]
  \item \textbf{Velocidad media:}
  \[
    v_m = \frac{\Delta x}{\Delta t} = \frac{x_f - x_i}{t_f - t_i}
  \]
\end{itemize}

\subsubsection*{1.2 Movimiento Rectilíneo Uniformemente Acelerado (MRUA)}
\addcontentsline{toc}{subsubsection}{1.2 Movimiento Rectilíneo Uniformemente Acelerado (MRUA)}

\begin{itemize}
  \item \textbf{Velocidad:}
  \[
    v(t) = v_0 + a \cdot t
  \]
  \item \textbf{Posición:}
  \[
    x(t) = x_0 + v_0 t + \frac{1}{2} a t^2
  \]
  \item \textbf{Ecuación independiente del tiempo (Torricelli):}
  \[
    v_f^2 = v_0^2 + 2a (x_f - x_0)
  \]
  \item \textbf{Desplazamiento medio:}
  \[
    \Delta x = \left( \frac{v_0 + v_f}{2} \right) t
  \]
\end{itemize}

\subsubsection*{1.3 Leyes del Movimiento de Newton}
\addcontentsline{toc}{subsubsection}{1.3 Leyes del Movimiento de Newton}

\begin{itemize}
  \item \textbf{Primera Ley (Inercia):}
  \[
    \sum \vec{F} = \vec{0} \iff \vec{v} = \text{constante}
  \]
  \item \textbf{Segunda Ley (Ley Fundamental de la Dinámica):}
  \[
    \sum \vec{F} = m \vec{a} = \frac{d\vec{p}}{dt}
  \]
  \item \textbf{Tercera Ley (Acción y Reacción):}
  \[
    \vec{F}_{AB} = -\vec{F}_{BA}
  \]
\end{itemize}

\subsubsection*{1.4 Fuerzas de Fricción / Rozamiento}
\addcontentsline{toc}{subsubsection}{1.4 Fuerzas de Fricción / Rozamiento}

\begin{itemize}
  \item \textbf{Fricción estática máxima:}
  \[
    f_{s,\text{máx}} = \mu_s N \quad \implies \quad f_s \le \mu_s N
  \]
  \item \textbf{Fricción cinética:}
  \[
    f_k = \mu_k N
  \]
\end{itemize}

\subsubsection*{1.5 Movimiento Circular Uniforme (MCU) y Uniformemente Variado (MCUV)}
\addcontentsline{toc}{subsubsection}{1.5 Movimiento Circular Uniforme (MCU) y Uniformemente Variado (MCUV)}

\begin{itemize}
  \item \textbf{Relaciones lineales y angulares:}
  \[
    s = \theta \cdot r, \quad v = \omega \cdot r, \quad a_t = \alpha \cdot r
  \]
  \item \textbf{Aceleración centrípeta / normal:}
  \[
    a_c = \frac{v^2}{r} = \omega^2 r
  \]
  \item \textbf{Cinemática angular con $\alpha$ constante:}
  \[
    \omega(t) = \omega_0 + \alpha t, \quad \theta(t) = \theta_0 + \omega_0 t + \frac{1}{2} \alpha t^2, \quad \omega_f^2 = \omega_0^2 + 2\alpha (\theta_f - \theta_0)
  \]
\end{itemize}

\subsection*{Nivel Intermedio}
\addcontentsline{toc}{subsection}{Nivel Intermedio}

\subsubsection*{1.6 Cinemática y Dinámica Vectorial Tridimensional}
\addcontentsline{toc}{subsubsection}{1.6 Cinemática y Dinámica Vectorial Tridimensional}

\begin{itemize}
  \item \textbf{Vectores cinemáticos fundamentales:}
  \[
    \vec{v}(t) = \frac{d\vec{r}}{dt}, \quad \vec{a}(t) = \frac{d\vec{v}}{dt} = \frac{d^2\vec{r}}{dt^2}
  \]
  \item \textbf{Descomposición intrínseca de la aceleración:}
  \[
    \vec{a} = a_t \hat{u}_t + a_n \hat{u}_n = \frac{dv}{dt} \hat{u}_t + \frac{v^2}{\rho} \hat{u}_n
  \]
  \item \textbf{Coordenadas polares planas:}
  \[
    \vec{v} = \dot{r} \hat{u}_r + r \dot{\theta} \hat{u}_\theta
  \]
  \[
    \vec{a} = (\ddot{r} - r\dot{\theta}^2)\hat{u}_r + (r\ddot{\theta} + 2\dot{r}\dot{\theta})\hat{u}_\theta
  \]
\end{itemize}

\subsubsection*{1.7 Sistemas de Referencia No Inerciales (Fuerzas Ficticias)}
\addcontentsline{toc}{subsubsection}{1.7 Sistemas de Referencia No Inerciales (Fuerzas Ficticias)}

\begin{itemize}
  \item \textbf{Aceleración en sistema móvil ($S'$ con origen acelerado $\vec{A}_0 = \frac{d^2\vec{R}_0}{dt^2}$ respecto al marco inercial $S$, y rotando a $\vec{\omega}$):}
  \[
    \vec{a} = \vec{A}_0 + \vec{a}' + \dot{\vec{\omega}} \times \vec{r}' + 2(\vec{\omega} \times \vec{v}') + \vec{\omega} \times (\vec{\omega} \times \vec{r}')
  \]
  \item \textbf{Ecuación de movimiento en el sistema no inercial ($S'$):}
  \[
    m\vec{a}' = \vec{F}_{\text{real}} - m\vec{A}_0 - m\dot{\vec{\omega}} \times \vec{r}' - 2m(\vec{\omega} \times \vec{v}') - m\vec{\omega} \times (\vec{\omega} \times \vec{r}')
  \]
  \item \textbf{Fuerza de Coriolis:}
  \[
    \vec{F}_{\text{Coriolis}} = -2m (\vec{\omega} \times \vec{v}')
  \]
  \item \textbf{Fuerza Centrífuga:}
  \[
    \vec{F}_{\text{centrífuga}} = -m\vec{\omega} \times (\vec{\omega} \times \vec{r}')
  \]
\end{itemize}

\subsubsection*{1.8 Dinámica Rotacional del Cuerpo Rígido (Eje Fijo)}
\addcontentsline{toc}{subsubsection}{1.8 Dinámica Rotacional del Cuerpo Rígido (Eje Fijo)}

\begin{itemize}
  \item \textbf{Momento de torsión / Torque:}
  \[
    \vec{\tau} = \vec{r} \times \vec{F}
  \]
  \item \textbf{Momento de inercia:}
  \[
    I = \sum_i m_i r_i^2 = \int r^2 dm
  \]
  \item \textbf{Ecuación fundamental de la rotación:}
  \[
    \sum \tau = I \alpha
  \]
  \item \textbf{Teorema de los Ejes Paralelos (Steiner):}
  \[
    I = I_{\text{CM}} + M d^2
  \]
  \item \textbf{Teorema de los Ejes Perpendiculares (Cuerpos laminares planos en $xy$):}
  \[
    I_z = I_x + I_y
  \]
\end{itemize}

\subsection*{Nivel Universitario / Avanzado}
\addcontentsline{toc}{subsection}{Nivel Universitario / Avanzado}

\subsubsection*{1.9 Dinámica Tridimensional del Sólido Rígido}
\addcontentsline{toc}{subsubsection}{1.9 Dinámica Tridimensional del Sólido Rígido}

\begin{itemize}
  \item \textbf{Tensor de Inercia $\mathbf{I}$ (con productos de inercia $I_{ij} = \int x_i x_j dm$):}
  \[
    \mathbf{I} = \begin{bmatrix} I_{xx} & -I_{xy} & -I_{xz} \\ -I_{yx} & I_{yy} & -I_{yz} \\ -I_{zx} & -I_{zy} & I_{zz} \end{bmatrix}, \quad I_{xx} = \int (y^2 + z^2) \, dm, \quad I_{xy} = \int xy \, dm
  \]
  \[
    \left( \text{En notación tensorial: } I_{ij} = \int (r^2 \delta_{ij} - x_i x_j) \, dm \right)
  \]
  \item \textbf{Momento angular tridimensional:}
  \[
    \vec{L} = \mathbf{I} \vec{\omega}
  \]
  \item \textbf{Ecuaciones de Euler para el cuerpo rígido (en el sistema de ejes principales):}
  \[
    I_1 \dot{\omega}_1 - (I_2 - I_3)\omega_2 \omega_3 = \tau_1
  \]
  \[
    I_2 \dot{\omega}_2 - (I_3 - I_1)\omega_3 \omega_1 = \tau_2
  \]
  \[
    I_3 \dot{\omega}_3 - (I_1 - I_2)\omega_1 \omega_2 = \tau_3
  \]
\end{itemize}

\section*{2. Trabajo, Energía, Momento Lineal y Gravitación}
\addcontentsline{toc}{section}{2. Trabajo, Energía, Momento Lineal y Gravitación}

\subsection*{Nivel Básico}
\addcontentsline{toc}{subsection}{Nivel Básico}

\subsubsection*{2.1 Trabajo y Energía Mecánica}
\addcontentsline{toc}{subsubsection}{2.1 Trabajo y Energía Mecánica}

\begin{itemize}
  \item \textbf{Trabajo efectuado por fuerza constante:}
  \[
    W = \vec{F} \cdot \Delta\vec{r} = F \Delta r \cos(\theta)
  \]
  \item \textbf{Energía Cinética traslacional:}
  \[
    E_k = \frac{1}{2} m v^2
  \]
  \item \textbf{Energía Potencial Gravitatoria (cerca de la superficie terrestre):}
  \[
    E_{p,\text{grav}} = m g h
  \]
  \item \textbf{Energía Potencial Elástica (Ley de Hooke $F = -kx$):}
  \[
    E_{p,\text{elás}} = \frac{1}{2} k x^2
  \]
  \item \textbf{Teorema del Trabajo y la Energía Cinética:}
  \[
    W_{\text{neto}} = \Delta E_k = \frac{1}{2} m v_f^2 - \frac{1}{2} m v_0^2
  \]
  \item \textbf{Conservación de la Energía Mecánica:}
  \[
    E_m = E_k + E_p \implies \Delta E_m = W_{\text{no conservativas}}
  \]
\end{itemize}

\subsubsection*{2.2 Cantidad de Movimiento e Impulso}
\addcontentsline{toc}{subsubsection}{2.2 Cantidad de Movimiento e Impulso}

\begin{itemize}
  \item \textbf{Momento lineal:}
  \[
    \vec{p} = m \vec{v}
  \]
  \item \textbf{Impulso:}
  \[
    \vec{J} = \vec{F}_{\text{prom}} \Delta t = \Delta \vec{p}
  \]
  \item \textbf{Choques unidimensionales y Coeficiente de Restitución ($e$):}
  \[
    e = \frac{v_{2f} - v_{1f}}{v_{1i} - v_{2i}}
  \]
\end{itemize}

\subsection*{Nivel Intermedio}
\addcontentsline{toc}{subsection}{Nivel Intermedio}

\subsubsection*{2.3 Mecánica de Sistemas de Partículas}
\addcontentsline{toc}{subsubsection}{2.3 Mecánica de Sistemas de Partículas}

\begin{itemize}
  \item \textbf{Trabajo de fuerza variable general:}
  \[
    W = \int_{C} \vec{F} \cdot d\vec{r}
  \]
  \item \textbf{Campos de fuerza conservativos y Potencial:}
  \[
    \vec{F} = -\nabla U = -\left( \frac{\partial U}{\partial x}\hat{i} + \frac{\partial U}{\partial y}\hat{j} + \frac{\partial U}{\partial z}\hat{k} \right)
  \]
  \item \textbf{Centro de Masas (CM):}
  \[
    \vec{R}_{\text{CM}} = \frac{\sum m_i \vec{r}_i}{\sum m_i} = \frac{1}{M}\int \vec{r} \, dm
  \]
  \item \textbf{Teoremas de König:}
  \[
    E_k = \frac{1}{2} M v_{\text{CM}}^2 + E_k' \quad (\text{traslación CM } + \text{ respecto a CM})
  \]
  \[
    \vec{L} = \vec{R}_{\text{CM}} \times \vec{P}_{\text{total}} + \vec{L}'
  \]
\end{itemize}

\subsubsection*{2.4 Gravitación Universal y Fuerzas Centrales}
\addcontentsline{toc}{subsubsection}{2.4 Gravitación Universal y Fuerzas Centrales}

\begin{itemize}
  \item \textbf{Ley de Gravitación de Newton:}
  \[
    \vec{F}_g = -G \frac{m_1 m_2}{r^2} \hat{u}_r
  \]
  \item \textbf{Energía Potencial Gravitatoria:}
  \[
    U(r) = -G \frac{m_1 m_2}{r}
  \]
  \item \textbf{Leyes de Kepler:}
  \item \textbf{2da Ley (Velocidad areolar constante):}
  \[
    \frac{dA}{dt} = \frac{L}{2m} = \text{constante}
  \]
  \item \textbf{3ra Ley:}
  \[
    T^2 = \left( \frac{4\pi^2}{G(M + m)} \right) a^3
  \]
  \item \textbf{Velocidades orbitales y de escape:}
  \[
    v_{\text{orb}} = \sqrt{\frac{GM}{r}}, \quad v_{\text{esc}} = \sqrt{\frac{2GM}{R}}
  \]
\end{itemize}

\subsection*{Nivel Universitario / Avanzado}
\addcontentsline{toc}{subsection}{Nivel Universitario / Avanzado}

\subsubsection*{2.5 Problema de Dos Cuerpos y Fuerzas Centrales}
\addcontentsline{toc}{subsubsection}{2.5 Problema de Dos Cuerpos y Fuerzas Centrales}

\begin{itemize}
  \item \textbf{Masa reducida ($\mu$):}
  \[
    \mu = \frac{m_1 m_2}{m_1 + m_2}
  \]
  \item \textbf{Ecuación de movimiento relativa:}
  \[
    \mu \ddot{\vec{r}} = f(r) \hat{u}_r
  \]
  \item \textbf{Potencial Efectivo:}
  \[
    U_{\text{eff}}(r) = U(r) + \frac{L^2}{2\mu r^2}
  \]
  \item \textbf{Ecuación diferencial de la órbita (Ecuación de Binet para $u = 1/r$):}
  \[
    \frac{d^2 u}{d\theta^2} + u = -\frac{\mu}{L^2 u^2} f\left(\frac{1}{u}\right)
  \]
  \item \textbf{Ecuación de las cónicas orbitales ($p = \text{semilatus rectum}$, $e = \text{excentricidad}$ para $U(r) = -G m_1 m_2 / r$):}
  \[
    r(\theta) = \frac{p}{1 + e \cos(\theta)}, \quad p = \frac{L^2}{\mu G m_1 m_2}, \quad e = \sqrt{1 + \frac{2 E L^2}{\mu (G m_1 m_2)^2}}
  \]
\end{itemize}

\section*{3. Oscilaciones y Mecánica de Fluidos}
\addcontentsline{toc}{section}{3. Oscilaciones y Mecánica de Fluidos}

\subsection*{Nivel Básico}
\addcontentsline{toc}{subsection}{Nivel Básico}

\subsubsection*{3.1 Oscilador Armónico Simple (M.A.S.)}
\addcontentsline{toc}{subsubsection}{3.1 Oscilador Armónico Simple (M.A.S.)}

\begin{itemize}
  \item \textbf{Ecuación diferencial y solución:}
  \[
    \ddot{x} + \omega_0^2 x = 0 \implies x(t) = A \cos(\omega_0 t + \phi)
  \]
  \item \textbf{Frecuencia angular y periodo (sistema masa-resorte):}
  \[
    \omega_0 = \sqrt{\frac{k}{m}}, \quad T = 2\pi\sqrt{\frac{m}{k}}
  \]
  \item \textbf{Péndulo Simple:}
  \[
    T = 2\pi\sqrt{\frac{L}{g}}
  \]
  \item \textbf{Péndulo Físico:}
  \[
    T = 2\pi\sqrt{\frac{I}{m g d}}
  \]
\end{itemize}

\subsubsection*{3.2 Estática de Fluidos}
\addcontentsline{toc}{subsubsection}{3.2 Estática de Fluidos}

\begin{itemize}
  \item \textbf{Presión hidrostática:}
  \[
    P(h) = P_0 + \rho g h
  \]
  \item \textbf{Principio de Pascal:}
  \[
    \frac{F_1}{A_1} = \frac{F_2}{A_2}
  \]
  \item \textbf{Principio de Arquímedes (Fuerza de Empuje):}
  \[
    E = \rho_{\text{fluido}} g V_{\text{sumergido}}
  \]
\end{itemize}

\subsection*{Nivel Intermedio}
\addcontentsline{toc}{subsection}{Nivel Intermedio}

\subsubsection*{3.3 Oscilaciones Amortiguadas y Forzadas}
\addcontentsline{toc}{subsubsection}{3.3 Oscilaciones Amortiguadas y Forzadas}

\begin{itemize}
  \item \textbf{Oscilador amortiguado ($b = \text{coeficiente de fricción}$):}
  \[
    \ddot{x} + 2\gamma \dot{x} + \omega_0^2 x = 0, \quad \gamma = \frac{b}{2m}
  \]
  \[
    x(t) = A e^{-\gamma t} \cos(\omega_d t + \phi), \quad \omega_d = \sqrt{\omega_0^2 - \gamma^2} \quad (\text{subamortiguado } \gamma < \omega_0)
  \]
  \item \textbf{Factor de Calidad $Q$:}
  \[
    Q = \frac{\omega_0}{2\gamma}
  \]
  \item \textbf{Oscilaciones forzadas con fuerza impulsora $F_0 \cos(\omega t)$:}
  \[
    \ddot{x} + 2\gamma \dot{x} + \omega_0^2 x = \frac{F_0}{m} \cos(\omega t)
  \]
  \[
    A(\omega) = \frac{F_0 / m}{\sqrt{(\omega_0^2 - \omega^2)^2 + 4\gamma^2 \omega^2}}
  \]
\end{itemize}

\subsubsection*{3.4 Dinámica de Fluidos Ideales}
\addcontentsline{toc}{subsubsection}{3.4 Dinámica de Fluidos Ideales}

\begin{itemize}
  \item \textbf{Ecuación de Continuidad (Flujo estacionario e incompresible 1D a través de secciones transversales):}
  \[
    A_1 v_1 = A_2 v_2 = Q \quad (\text{Caudal volumétrico constante})
  \]
  \[
    \left( \text{Para flujo compresible estacionario: } \rho_1 A_1 v_1 = \rho_2 A_2 v_2 = \dot{m} = \text{constante} \right)
  \]
  \item \textbf{Ecuación de Bernoulli:}
  \[
    P_1 + \frac{1}{2}\rho v_1^2 + \rho g z_1 = P_2 + \frac{1}{2}\rho v_2^2 + \rho g z_2 = \text{constante}
  \]
  \item \textbf{Teorema de Torricelli:}
  \[
    v = \sqrt{2gh}
  \]
\end{itemize}

\subsection*{Nivel Universitario / Avanzado}
\addcontentsline{toc}{subsection}{Nivel Universitario / Avanzado}

\subsubsection*{3.5 Mecánica de Fluidos y Medios Continuos}
\addcontentsline{toc}{subsubsection}{3.5 Mecánica de Fluidos y Medios Continuos}

\begin{itemize}
  \item \textbf{Ecuación de Continuidad diferencial:}
  \[
    \frac{\partial \rho}{\partial t} + \nabla \cdot (\rho \vec{v}) = 0
  \]
  \item \textbf{Ecuación de Navier-Stokes (Fluidos Newtonianos incompresibles con viscosidad $\mu$):}
  \[
    \rho \left( \frac{\partial \vec{v}}{\partial t} + (\vec{v} \cdot \nabla)\vec{v} \right) = -\nabla P + \mu \nabla^2 \vec{v} + \rho \vec{g}
  \]
  \item \textbf{Tensor de Esfuerzos de Cauchy $\boldsymbol{\sigma}$:}
  \[
    \nabla \cdot \boldsymbol{\sigma} + \vec{f}_{\text{ext}} = \rho \frac{d\vec{v}}{dt}
  \]
  \item \textbf{Ecuación de Navier-Cauchy para Sólidos Elásticos Lineales (Parámetros de Lamé $\lambda, \mu$):}
  \[
    \rho \frac{\partial^2 \vec{u}}{\partial t^2} = (\lambda + \mu)\nabla(\nabla \cdot \vec{u}) + \mu \nabla^2 \vec{u} + \vec{f}_{\text{vol}}
  \]
\end{itemize}

\section*{4. Mecánica Analítica y Relativista}
\addcontentsline{toc}{section}{4. Mecánica Analítica y Relativista}

\subsection*{Nivel Intermedio / Universitario}
\addcontentsline{toc}{subsection}{Nivel Intermedio / Universitario}

\subsubsection*{4.1 Mecánica Lagrangiana}
\addcontentsline{toc}{subsubsection}{4.1 Mecánica Lagrangiana}

\begin{itemize}
  \item \textbf{Coordenadas generalizadas:}
  \[
    q = (q_1, q_2, \dots, q_n), \quad \dot{q} = (\dot{q}_1, \dot{q}_2, \dots, \dot{q}_n)
  \]
  \item \textbf{Función Lagrangiana:}
  \[
    L(q, \dot{q}, t) = T(q, \dot{q}, t) - V(q, t)
  \]
  \item \textbf{Principio de Mínima Acción de Hamilton:}
  \[
    \delta S = \delta \int_{t_1}^{t_2} L(q, \dot{q}, t) dt = 0
  \]
  \item \textbf{Ecuaciones de Euler-Lagrange:}
  \[
    \frac{d}{dt}\left( \frac{\partial L}{\partial \dot{q}_j} \right) - \frac{\partial L}{\partial q_j} = 0, \quad j = 1, \dots, n
  \]
  \item \textbf{Momento conjugado o canónico:}
  \[
    p_j = \frac{\partial L}{\partial \dot{q}_j}
  \]
  \item \textbf{Teorema de Noether:}
  \[
    \text{Si } \frac{\partial L}{\partial q_j} = 0 \implies p_j = \text{constante de movimiento}
  \]
\end{itemize}

\subsubsection*{4.2 Mecánica Hamiltoniana}
\addcontentsline{toc}{subsubsection}{4.2 Mecánica Hamiltoniana}

\begin{itemize}
  \item \textbf{Función Hamiltoniana (Transformada de Legendre):}
  \[
    H(q, p, t) = \sum_{j=1}^n p_j \dot{q}_j - L(q, \dot{q}, t)
  \]
  \item \textbf{Ecuaciones Canónicas de Hamilton:}
  \[
    \dot{q}_j = \frac{\partial H}{\partial p_j}, \quad \dot{p}_j = -\frac{\partial H}{\partial q_j}
  \]
  \item \textbf{Variación temporal del Hamiltoniano:}
  \[
    \frac{dH}{dt} = \frac{\partial H}{\partial t} = -\frac{\partial L}{\partial t}
  \]
\end{itemize}

\subsection*{Nivel Universitario Avanzado}
\addcontentsline{toc}{subsection}{Nivel Universitario Avanzado}

\subsubsection*{4.3 Formulaciones Canónicas Avanzadas}
\addcontentsline{toc}{subsubsection}{4.3 Formulaciones Canónicas Avanzadas}

\begin{itemize}
  \item \textbf{Corchetes de Poisson:}
  \[
    \{f, g\}_{q,p} = \sum_{j=1}^n \left( \frac{\partial f}{\partial q_j} \frac{\partial g}{\partial p_j} - \frac{\partial f}{\partial p_j} \frac{\partial g}{\partial q_j} \right)
  \]
  \item \textbf{Evolución temporal de una función observable:}
  \[
    \frac{df}{dt} = \{f, H\} + \frac{\partial f}{\partial t}
  \]
  \item \textbf{Teorema de Liouville (Conservación del volumen en el espacio fásico):}
  \[
    \frac{d\rho}{dt} = \frac{\partial \rho}{\partial t} + \{\rho, H\} = 0
  \]
  \item \textbf{Ecuación de Hamilton-Jacobi:}
  \[
    H\left( q_1, \dots, q_n, \frac{\partial S}{\partial q_1}, \dots, \frac{\partial S}{\partial q_n}, t \right) + \frac{\partial S}{\partial t} = 0
  \]
\end{itemize}

\subsubsection*{4.4 Mecánica Relativista (Relatividad Especial)}
\addcontentsline{toc}{subsubsection}{4.4 Mecánica Relativista (Relatividad Especial)}

\begin{itemize}
  \item \textbf{Factor de Lorentz:}
  \[
    \gamma = \frac{1}{\sqrt{1 - \frac{v^2}{c^2}}}
  \]
  \item \textbf{Momento lineal relativista:}
  \[
    \vec{p} = \gamma m_0 \vec{v}
  \]
  \item \textbf{Fuerza relativista:}
  \[
    \vec{F} = \frac{d\vec{p}}{dt} = \frac{d}{dt}(\gamma m_0 \vec{v})
  \]
  \item \textbf{Energía total, cinética y en reposo:}
  \[
    E = \gamma m_0 c^2 = E_k + m_0 c^2
  \]
  \[
    E_k = (\gamma - 1) m_0 c^2
  \]
  \item \textbf{Relación de dispersión Energía-Momento:}
  \[
    E^2 = (p c)^2 + (m_0 c^2)^2
  \]
  \item \textbf{Cuadrivectores en Cinemática y Dinámica:}
  \[
    X^\mu = (ct, \vec{r}), \quad U^\mu = \frac{dX^\mu}{d\tau} = \gamma(c, \vec{v})
  \]
  \[
    P^\mu = m_0 U^\mu = \left( \frac{E}{c}, \vec{p} \right), \quad P^\mu P_\mu = \frac{E^2}{c^2} - |\vec{p}|^2 = m_0^2 c^2
  \]
\end{itemize}

\section*{Compendio II: Oscilaciones y Ondas Mecánicas}
\addcontentsline{toc}{section}{Compendio II: Oscilaciones y Ondas Mecánicas}

\textit{(Nivel Básico a Universitario)}

\section*{5. Movimiento Oscilatorio (Oscilaciones)}
\addcontentsline{toc}{section}{5. Movimiento Oscilatorio (Oscilaciones)}

\subsection*{Nivel Básico}
\addcontentsline{toc}{subsection}{Nivel Básico}

\subsubsection*{5.1 Cinemática del Movimiento Armónico Simple (M.A.S.)}
\addcontentsline{toc}{subsubsection}{5.1 Cinemática del Movimiento Armónico Simple (M.A.S.)}

\begin{itemize}
  \item \textbf{Posición en función del tiempo:}
  \[
    x(t) = A \cos(\omega t + \phi)
  \]
  \item \textbf{Velocidad:}
  \[
    v(t) = \frac{dx}{dt} = -\omega A \sen(\omega t + \phi) = \pm \omega \sqrt{A^2 - x^2}
  \]
  \item \textbf{Aceleración:}
  \[
    a(t) = \frac{d^2x}{dt^2} = -\omega^2 A \cos(\omega t + \phi) = -\omega^2 x(t)
  \]
  \item \textbf{Valores máximos:}
  \[
    x_{\text{máx}} = A, \quad v_{\text{máx}} = \omega A, \quad a_{\text{máx}} = \omega^2 A
  \]
\end{itemize}

\subsubsection*{5.2 Parámetros Temporales y Frecuenciales}
\addcontentsline{toc}{subsubsection}{5.2 Parámetros Temporales y Frecuenciales}

\begin{itemize}
  \item \textbf{Periodo ($T$):}
  \[
    T = \frac{1}{f} = \frac{2\pi}{\omega}
  \]
  \item \textbf{Frecuencia ($f$):}
  \[
    f = \frac{1}{T} = \frac{\omega}{2\pi}
  \]
  \item \textbf{Frecuencia angular ($\omega$):}
  \[
    \omega = 2\pi f = \frac{2\pi}{T}
  \]
\end{itemize}

\subsubsection*{5.3 Energía del M.A.S. (Sistema Conservativo)}
\addcontentsline{toc}{subsubsection}{5.3 Energía del M.A.S. (Sistema Conservativo)}

\begin{itemize}
  \item \textbf{Energía Cinética:}
  \[
    E_k = \frac{1}{2} m v^2 = \frac{1}{2} m \omega^2 A^2 \sen^2(\omega t + \phi)
  \]
  \item \textbf{Energía Potencial Elástica:}
  \[
    E_p = \frac{1}{2} k x^2 = \frac{1}{2} k A^2 \cos^2(\omega t + \phi)
  \]
  \item \textbf{Energía Mecánica Total:}
  \[
    E_m = E_k + E_p = \frac{1}{2} k A^2 = \frac{1}{2} m \omega^2 A^2 = \text{constante}
  \]
\end{itemize}

\subsubsection*{5.4 Sistemas Oscilatorios Simples}
\addcontentsline{toc}{subsubsection}{5.4 Sistemas Oscilatorios Simples}

\begin{itemize}
  \item \textbf{Masa-Resorte:}
  \[
    \omega_0 = \sqrt{\frac{k}{m}}, \quad T = 2\pi\sqrt{\frac{m}{k}}
  \]
  \item \textbf{Péndulo Simple (Aproximación para ángulos pequeños $\sen\theta \approx \theta$):}
  \[
    \omega_0 = \sqrt{\frac{g}{L}}, \quad T = 2\pi\sqrt{\frac{L}{g}}
  \]
\end{itemize}

\subsection*{Nivel Intermedio}
\addcontentsline{toc}{subsection}{Nivel Intermedio}

\subsubsection*{5.5 Péndulos Compuestos y de Torsión}
\addcontentsline{toc}{subsubsection}{5.5 Péndulos Compuestos y de Torsión}

\begin{itemize}
  \item \textbf{Péndulo Físico (Centro de masa a distancia $d$ del eje de giro):}
  \[
    \omega_0 = \sqrt{\frac{m g d}{I}}, \quad T = 2\pi\sqrt{\frac{I}{m g d}}
  \]
  \item \textbf{Péndulo de Torsión (Constante de torsión $\kappa$):}
  \[
    \tau = -\kappa \theta \implies \omega_0 = \sqrt{\frac{\kappa}{I}}, \quad T = 2\pi\sqrt{\frac{I}{\kappa}}
  \]
\end{itemize}

\subsubsection*{5.6 Oscilaciones Amortiguadas}
\addcontentsline{toc}{subsubsection}{5.6 Oscilaciones Amortiguadas}

\begin{itemize}
  \item \textbf{Ecuación Diferencial Fundamental ($b = \text{coeficiente de amortiguamiento viscoso}$):}
  \[
    m \ddot{x} + b \dot{x} + k x = 0 \implies \ddot{x} + 2\gamma \dot{x} + \omega_0^2 x = 0, \quad \gamma = \frac{b}{2m}, \quad \omega_0 = \sqrt{\frac{k}{m}}
  \]
  \item \textbf{Régimen Subamortiguado ($\gamma < \omega_0$):}
  \[
    x(t) = A_0 e^{-\gamma t} \cos(\omega_d t + \phi), \quad \omega_d = \sqrt{\omega_0^2 - \gamma^2}
  \]
  \item \textbf{Régimen Críticamente Amortiguado ($\gamma = \omega_0$):}
  \[
    x(t) = (C_1 + C_2 t) e^{-\gamma t}
  \]
  \item \textbf{Régimen Sobreamortiguado ($\gamma > \omega_0$):}
  \[
    x(t) = C_1 e^{-(\gamma - \sqrt{\gamma^2 - \omega_0^2})t} + C_2 e^{-(\gamma + \sqrt{\gamma^2 - \omega_0^2})t}
  \]
  \item \textbf{Decremento Logarítmico ($\delta$):}
  \[
    \delta = \ln\left( \frac{x(t)}{x(t + T_d)} \right) = \gamma T_d = \frac{2\pi \gamma}{\omega_d}
  \]
  \item \textbf{Factor de Calidad ($Q$):}
  \[
    Q = \frac{\omega_0}{2\gamma} = \frac{\pi}{\delta} = 2\pi \frac{E}{\Delta E_{\text{ciclo}}}
  \]
\end{itemize}

\subsubsection*{5.7 Oscilaciones Forzadas y Resonancia}
\addcontentsline{toc}{subsubsection}{5.7 Oscilaciones Forzadas y Resonancia}

\begin{itemize}
  \item \textbf{Ecuación Diferencial con Excitación Armónica:}
  \[
    \ddot{x} + 2\gamma \dot{x} + \omega_0^2 x = \frac{F_0}{m} \cos(\omega t)
  \]
  \item \textbf{Solución de Estado Estacionario ($x_p(t) = A(\omega) \cos(\omega t - \delta)$):}
  \[
    A(\omega) = \frac{F_0 / m}{\sqrt{(\omega_0^2 - \omega^2)^2 + 4\gamma^2 \omega^2}}
  \]
  \[
    \tan(\delta) = \frac{2\gamma \omega}{\omega_0^2 - \omega^2}
  \]
  \item \textbf{Frecuencia de Resonancia de Amplitud:}
  \[
    \omega_{\text{res}} = \sqrt{\omega_0^2 - 2\gamma^2}
  \]
  \item \textbf{Potencia Media Absorbida por el Oscilador:}
  \[
    \langle P(\omega) \rangle = \frac{F_0^2 \gamma \omega^2}{m [(\omega_0^2 - \omega^2)^2 + 4\gamma^2 \omega^2]}
  \]
\end{itemize}

\subsection*{Nivel Universitario / Avanzado}
\addcontentsline{toc}{subsection}{Nivel Universitario / Avanzado}

\subsubsection*{5.8 Osciladores Acoplados y Modos Normales}
\addcontentsline{toc}{subsubsection}{5.8 Osciladores Acoplados y Modos Normales}

\begin{itemize}
  \item \textbf{Ecuación Matricial de Movimiento (Sistemas lineales de $N$ grados de libertad):}
  \[
    \mathbf{M} \ddot{\vec{x}} + \mathbf{K} \vec{x} = \vec{0}
  \]
  \item \textbf{Ecuación Secular / Polinomio Característico:}
  \[
    \det(\mathbf{K} - \omega^2 \mathbf{M}) = 0
  \]
  \item \textbf{Transformación a Coordenadas Normales ($\vec{\eta} = \mathbf{P}^{-1} \vec{x}$):}
  \[
    \ddot{\eta}_k + \omega_k^2 \eta_k = 0, \quad k = 1, 2, \dots, N
  \]
\end{itemize}

\subsubsection*{5.9 Oscilaciones No Lineales (Péndulo Simple de Gran Amplitud)}
\addcontentsline{toc}{subsubsection}{5.9 Oscilaciones No Lineales (Péndulo Simple de Gran Amplitud)}

\begin{itemize}
  \item \textbf{Ecuación Diferencial Exacta:}
  \[
    \ddot{\theta} + \omega_0^2 \sen(\theta) = 0
  \]
  \item \textbf{Periodo Exacto mediante Integrales Elípticas Completas de Primera Especie $K(m)$:}
  \[
    T = 4\sqrt{\frac{L}{g}} K\left(\sen^2\left(\frac{\theta_0}{2}\right)\right) = 4\sqrt{\frac{L}{g}} \int_0^{\pi/2} \frac{d\phi}{\sqrt{1 - \sen^2\left(\frac{\theta_0}{2}\right) \sen^2\phi}}
  \]
  \item \textbf{Aproximación de Borda (Series de Taylor):}
  \[
    T \approx 2\pi\sqrt{\frac{L}{g}} \left( 1 + \frac{1}{16}\theta_0^2 + \frac{11}{3072}\theta_0^4 + \dots \right)
  \]
\end{itemize}

\subsubsection*{5.10 Osciladores Autoexcitados (Ecuación de Van der Pol)}
\addcontentsline{toc}{subsubsection}{5.10 Osciladores Autoexcitados (Ecuación de Van der Pol)}

\begin{itemize}
  \item \textbf{Ecuación con amortiguamiento no lineal:}
  \[
    \ddot{x} - \mu(1 - x^2)\dot{x} + \omega_0^2 x = 0
  \]
\end{itemize}

\section*{6. Ondas Mecánicas en Medios Continuos}
\addcontentsline{toc}{section}{6. Ondas Mecánicas en Medios Continuos}

\subsection*{Nivel Básico}
\addcontentsline{toc}{subsection}{Nivel Básico}

\subsubsection*{6.1 Cinemática de Ondas Armónicas Unidimensionales}
\addcontentsline{toc}{subsubsection}{6.1 Cinemática de Ondas Armónicas Unidimensionales}

\begin{itemize}
  \item \textbf{Función de Onda Progresiva (hacia la derecha $-$ / hacia la izquierda $+$):}
  \[
    y(x, t) = A \sen(k x \mp \omega t + \phi)
  \]
  \item \textbf{Relación entre Parámetros Fundamentales:}
  \[
    k = \frac{2\pi}{\lambda}, \quad \omega = \frac{2\pi}{T} = 2\pi f, \quad v = \lambda f = \frac{\lambda}{T} = \frac{\omega}{k}
  \]
  \item \textbf{Velocidad y Aceleración Transversal de las Partículas del Medio:}
  \[
    v_y(x, t) = \frac{\partial y}{\partial t} = \mp \omega A \cos(k x \mp \omega t + \phi)
  \]
  \[
    a_y(x, t) = \frac{\partial^2 y}{\partial t^2} = -\omega^2 A \sen(k x \mp \omega t + \phi) = -\omega^2 y(x, t)
  \]
\end{itemize}

\subsection*{Nivel Intermedio}
\addcontentsline{toc}{subsection}{Nivel Intermedio}

\subsubsection*{6.2 Ecuación Diferencial de Onda 1D (d'Alembert)}
\addcontentsline{toc}{subsubsection}{6.2 Ecuación Diferencial de Onda 1D (d'Alembert)}

\begin{itemize}
  \item \textbf{Forma canónica:}
  \[
    \frac{\partial^2 y}{\partial x^2} = \frac{1}{v^2} \frac{\partial^2 y}{\partial t^2}
  \]
  \item \textbf{Solución General de d'Alembert:}
  \[
    y(x, t) = f(x - vt) + g(x + vt)
  \]
\end{itemize}

\subsubsection*{6.3 Velocidad de Propagación en Medios Elásticos}
\addcontentsline{toc}{subsubsection}{6.3 Velocidad de Propagación en Medios Elásticos}

\begin{itemize}
  \item \textbf{Cuerda Tensa (Tensión $T$, densidad lineal de masa $\mu = m/L$):}
  \[
    v = \sqrt{\frac{T}{\mu}}
  \]
  \item \textbf{Barra Sólida (Ondas longitudinales, módulo de Young $Y$, densidad $\rho$):}
  \[
    v = \sqrt{\frac{Y}{\rho}}
  \]
  \item \textbf{Fluido / Gas (Módulo de compresibilidad volumétrica $B$):}
  \[
    v = \sqrt{\frac{B}{\rho}}
  \]
  \item \textbf{Gas Ideal (Ecuación de Newton-Laplace para proceso adiabático):}
  \[
    v = \sqrt{\frac{\gamma R T}{M}} = \sqrt{\frac{\gamma P}{\rho}}
  \]
\end{itemize}

\subsubsection*{6.4 Transporte de Energía e Intensidad}
\addcontentsline{toc}{subsubsection}{6.4 Transporte de Energía e Intensidad}

\begin{itemize}
  \item \textbf{Densidad Lineal de Energía Mecánica Total:}
  \[
    \frac{dE}{dx} = \mu \omega^2 A^2 \cos^2(kx - \omega t)
  \]
  \item \textbf{Potencia Instantánea Transmitida en Cuerda:}
  \[
    P(x, t) = -T \left( \frac{\partial y}{\partial x} \right) \left( \frac{\partial y}{\partial t} \right) = \sqrt{\mu T} \, \omega^2 A^2 \cos^2(kx - \omega t)
  \]
  \item \textbf{Potencia Promedio Transmitida:}
  \[
    \langle P \rangle = \frac{1}{2} \mu v \omega^2 A^2 = \frac{1}{2} \sqrt{\mu T} \, \omega^2 A^2
  \]
  \item \textbf{Intensidad de Onda ($I$):}
  \[
    I = \frac{\langle P \rangle}{\text{Área}} = \frac{1}{2} \rho v \omega^2 A^2
  \]
  \item \textbf{Ley del Inverso del Cuadrado (Fuente puntual isótropa en medio no disipativo):}
  \[
    I(r) = \frac{P_{\text{fuente}}}{4\pi r^2} \implies \frac{I_1}{I_2} = \frac{r_2^2}{r_1^2}
  \]
\end{itemize}

\subsection*{Nivel Universitario / Avanzado}
\addcontentsline{toc}{subsection}{Nivel Universitario / Avanzado}

\subsubsection*{6.5 Dispersión de Ondas}
\addcontentsline{toc}{subsubsection}{6.5 Dispersión de Ondas}

\begin{itemize}
  \item \textbf{Relación de Dispersión:}
  \[
    \omega = \omega(k)
  \]
  \item \textbf{Velocidad de Fase:}
  \[
    v_p = \frac{\omega}{k}
  \]
  \item \textbf{Velocidad de Grupo:}
  \[
    v_g = \frac{d\omega}{dk} = v_p + k \frac{dv_p}{dk} = v_p - \lambda \frac{dv_p}{d\lambda}
  \]
\end{itemize}

\subsubsection*{6.6 Impedancia Mecánica y Fenómenos en Fronteras}
\addcontentsline{toc}{subsubsection}{6.6 Impedancia Mecánica y Fenómenos en Fronteras}

\begin{itemize}
  \item \textbf{Impedancia Característica del Medio:}
  \[
    Z = \rho v = \sqrt{\rho B} \quad (\text{Fluidos}), \quad Z = \mu v = \sqrt{\mu T} \quad (\text{Cuerdas})
  \]
  \item \textbf{Coeficiente de Reflexión en Amplitud (para ondas de desplazamiento / velocidad incidiendo de medio 1 a medio 2):}
  \[
    r = \frac{Z_1 - Z_2}{Z_1 + Z_2}
  \]
  \item \textbf{Coeficiente de Transmisión en Amplitud (desplazamiento / velocidad):}
  \[
    t = \frac{2Z_1}{Z_1 + Z_2}
  \]
  \item \textbf{Coeficientes de Reflexión y Transmisión de Potencia / Intensidad ($R + T_w = 1$):}
  \[
    R = r^2 = \left( \frac{Z_1 - Z_2}{Z_1 + Z_2} \right)^2, \quad T_w = \frac{Z_2}{Z_1} t^2 = \frac{4 Z_1 Z_2}{(Z_1 + Z_2)^2}
  \]
\end{itemize}

\subsubsection*{6.7 Ondas Elásticas en Medios Continuos 3D (Ecuación de Navier-Cauchy)}
\addcontentsline{toc}{subsubsection}{6.7 Ondas Elásticas en Medios Continuos 3D (Ecuación de Navier-Cauchy)}

\begin{itemize}
  \item \textbf{Ondas Longitudinales Primarias (Ondas P / Compresionales):}
  \[
    v_P = \sqrt{\frac{\lambda + 2\mu}{\rho}} = \sqrt{\frac{K + \frac{4}{3}G}{\rho}}
  \]
  \item \textbf{Ondas Transversales Secundarias (Ondas S / Cizalladura):}
  \[
    v_S = \sqrt{\frac{\mu}{\rho}} = \sqrt{\frac{G}{\rho}}
  \]
\end{itemize}

\section*{7. Interferencia, Ondas Estacionarias y Acústica}
\addcontentsline{toc}{section}{7. Interferencia, Ondas Estacionarias y Acústica}

\subsection*{Nivel Básico}
\addcontentsline{toc}{subsection}{Nivel Básico}

\subsubsection*{7.1 Superposición e Interferencia}
\addcontentsline{toc}{subsubsection}{7.1 Superposición e Interferencia}

\begin{itemize}
  \item \textbf{Interferencia de dos ondas armónicas con diferencia de fase $\Delta\phi$:}
  \[
    y_R(x, t) = \left[ 2A \cos\left(\frac{\Delta\phi}{2}\right) \right] \sen\left(kx - \omega t + \frac{\Delta\phi}{2}\right)
  \]
  \item \textbf{Condición de Interferencia Constructiva:}
  \[
    \Delta\phi = 2n\pi \implies \Delta r = n\lambda, \quad n \in \mathbb{Z}
  \]
  \item \textbf{Condición de Interferencia Destructiva:}
  \[
    \Delta\phi = (2n + 1)\pi \implies \Delta r = \left(n + \frac{1}{2}\right)\lambda, \quad n \in \mathbb{Z}
  \]
\end{itemize}

\subsubsection*{7.2 Batidos o Pulsaciones}
\addcontentsline{toc}{subsubsection}{7.2 Batidos o Pulsaciones}

\begin{itemize}
  \item \textbf{Superposición de ondas con frecuencias ligeramente diferentes $\omega_1 \approx \omega_2$:}
  \[
    y_R(t) = \left[ 2A \cos\left( \frac{\omega_1 - \omega_2}{2} t \right) \right] \cos\left( \frac{\omega_1 + \omega_2}{2} t \right)
  \]
  \item \textbf{Frecuencia de Batido:}
  \[
    f_{\text{batido}} = |f_1 - f_2|
  \]
\end{itemize}

\subsubsection*{7.3 Escala Decibélica y Nivel de Sonido}
\addcontentsline{toc}{subsubsection}{7.3 Escala Decibélica y Nivel de Sonido}

\begin{itemize}
  \item \textbf{Nivel de Intensidad Sonora ($\beta$):}
  \[
    \beta = 10 \log_{10}\left( \frac{I}{I_0} \right), \quad I_0 = 10^{-12} \, \text{W/m}^2
  \]
  \item \textbf{Nivel de Presión Sonora ($SPL$):}
  \[
    L_p = 20 \log_{10}\left( \frac{p_{\text{rms}}}{p_0} \right), \quad p_0 = 20 \, \mu\text{Pa}
  \]
\end{itemize}

\subsection*{Nivel Intermedio}
\addcontentsline{toc}{subsection}{Nivel Intermedio}

\subsubsection*{7.4 Ondas Estacionarias Unidimensionales}
\addcontentsline{toc}{subsubsection}{7.4 Ondas Estacionarias Unidimensionales}

\begin{itemize}
  \item \textbf{Ecuación Analítica:}
  \[
    y(x, t) = [2A \sen(kx)] \cos(\omega t)
  \]
  \item \textbf{Posición de Nodos (Amplitud nula):}
  \[
    \sen(kx) = 0 \implies x_n = n \frac{\lambda}{2}, \quad n = 0, 1, 2, \dots
  \]
  \item \textbf{Posición de Antinodos o Vientres (Amplitud máxima $2A$):}
  \[
    |\sen(kx)| = 1 \implies x_a = \left( n + \frac{1}{2} \right) \frac{\lambda}{2}, \quad n = 0, 1, 2, \dots
  \]
\end{itemize}

\subsubsection*{7.5 Modos Normales y Frecuencias Resonantes en Sistemas Acotados}
\addcontentsline{toc}{subsubsection}{7.5 Modos Normales y Frecuencias Resonantes en Sistemas Acotados}

\begin{itemize}
  \item \textbf{Cuerda Fija en Ambos Extremos o Tubo Abierto-Abierto:}
  \[
    \lambda_n = \frac{2L}{n}, \quad f_n = n \frac{v}{2L} = n f_1, \quad n = 1, 2, 3, \dots
  \]
  \item \textbf{Cuerda con un Extremo Fijo y uno Libre o Tubo Abierto-Cerrado:}
  \[
    \lambda_n = \frac{4L}{2n - 1}, \quad f_n = (2n - 1) \frac{v}{4L} = (2n - 1) f_1, \quad n = 1, 2, 3, \dots
  \]
\end{itemize}

\subsubsection*{7.6 Efecto Doppler Clásico (Velocidad del sonido $v$)}
\addcontentsline{toc}{subsubsection}{7.6 Efecto Doppler Clásico (Velocidad del sonido $v$)}

\begin{itemize}
  \item \textbf{Fórmula General:}
  \[
    f_o = f_s \left( \frac{v \pm v_o}{v \mp v_s} \right)
  \]
  \[
    \text{Signo superior: acercamiento} \quad | \quad \text{Signo inferior: alejamiento}
  \]
\end{itemize}

\subsubsection*{7.7 Ondas de Choque y Cono de Mach}
\addcontentsline{toc}{subsubsection}{7.7 Ondas de Choque y Cono de Mach}

\begin{itemize}
  \item \textbf{Ángulo de Mach ($\theta_M$):}
  \[
    \sen(\theta_M) = \frac{v}{v_s} = \frac{1}{M}, \quad M = \frac{v_s}{v} > 1
  \]
\end{itemize}

\subsection*{Nivel Universitario / Avanzado}
\addcontentsline{toc}{subsection}{Nivel Universitario / Avanzado}

\subsubsection*{7.8 Acústica Física y Ondas de Presión}
\addcontentsline{toc}{subsubsection}{7.8 Acústica Física y Ondas de Presión}

\begin{itemize}
  \item \textbf{Onda de Desplazamiento Molecular:}
  \[
    s(x, t) = s_{\text{máx}} \cos(kx - \omega t)
  \]
  \item \textbf{Onda de Presión Acústica Excedente:}
  \[
    \Delta p(x, t) = -B \frac{\partial s}{\partial x} = \Delta p_{\text{máx}} \sen(kx - \omega t)
  \]
  \item \textbf{Amplitud de Presión:}
  \[
    \Delta p_{\text{máx}} = B k s_{\text{máx}} = \rho v \omega s_{\text{máx}} = Z \omega s_{\text{máx}}
  \]
  \item \textbf{Relación entre Intensidad y Presión Acústica:}
  \[
    I = \frac{\Delta p_{\text{máx}}^2}{2\rho v} = \frac{p_{\text{rms}}^2}{\rho v}
  \]
\end{itemize}

\subsubsection*{7.9 Ecuación de Onda Acústica Tridimensional}
\addcontentsline{toc}{subsubsection}{7.9 Ecuación de Onda Acústica Tridimensional}

\begin{itemize}
  \item \textbf{Ecuación Diferencial para la Presión:}
  \[
    \nabla^2 p - \frac{1}{c^2} \frac{\partial^2 p}{\partial t^2} = 0
  \]
  \item \textbf{Solución para Ondas Esféricas Monocromáticas:}
  \[
    p(r, t) = \frac{A}{r} e^{i(kr - \omega t)}
  \]
  \item \textbf{Potencial de Velocidad Acústico ($\vec{u} = -\nabla \Phi$):}
  \[
    \nabla^2 \Phi - \frac{1}{c^2} \frac{\partial^2 \Phi}{\partial t^2} = 0, \quad p = \rho_0 \frac{\partial \Phi}{\partial t}
  \]
\end{itemize}

\subsubsection*{7.10 Efecto Doppler Vectorial en Dirección Arbitraria}
\addcontentsline{toc}{subsubsection}{7.10 Efecto Doppler Vectorial en Dirección Arbitraria}

\begin{itemize}
  \item \textbf{Frecuencia Observada:}
  \[
    f_o = f_s \left( \frac{c - \vec{v}_o \cdot \hat{n}}{c - \vec{v}_s \cdot \hat{n}} \right)
  \]
  \[
    (\hat{n} = \text{vector unitario dirigido de la fuente al observador})
  \]
\end{itemize}

\subsubsection*{7.11 Atenuación y Absorción Acústica Viscotérmica}
\addcontentsline{toc}{subsubsection}{7.11 Atenuación y Absorción Acústica Viscotérmica}

\begin{itemize}
  \item \textbf{Decaimiento Espacial de Amplitud:}
  \[
    A(x) = A_0 e^{-\alpha x}
  \]
  \item \textbf{Coeficiente de Atenuación Clásico de Stokes-Kirchhoff:}
  \[
    \alpha = \frac{\omega^2}{2\rho c^3} \left[ \frac{4}{3}\eta + \eta_v + \frac{(\gamma - 1)\kappa}{C_p} \right]
  \]
\end{itemize}

\section*{Compendio III: Termodinámica Clásica, Química y Estadística}
\addcontentsline{toc}{section}{Compendio III: Termodinámica Clásica, Química y Estadística}

\textit{(Nivel Básico a Universitario)}

\section*{8. Conceptos Fundamentales, Gases y Ecuaciones de Estado}
\addcontentsline{toc}{section}{8. Conceptos Fundamentales, Gases y Ecuaciones de Estado}

\subsection*{Nivel Básico}
\addcontentsline{toc}{subsection}{Nivel Básico}

\subsubsection*{8.1 Escalas Termométricas y Temperatura}
\addcontentsline{toc}{subsubsection}{8.1 Escalas Termométricas y Temperatura}

\begin{itemize}
  \item \textbf{Escala Celsius a Kelvin:}
  \[
    T(\text{K}) = T(^\circ\text{C}) + 273.15
  \]
  \item \textbf{Escala Celsius a Fahrenheit:}
  \[
    T(^\circ\text{F}) = \frac{9}{5} T(^\circ\text{C}) + 32
  \]
  \item \textbf{Escala Fahrenheit a Rankine (Temperatura absoluta $T_{\text{R}}$):}
  \[
    T_{\text{R}}(\text{Ra}) = T(^\circ\text{F}) + 459.67 = \frac{9}{5} T(\text{K})
  \]
\end{itemize}

\subsubsection*{8.2 Dilatación Térmica}
\addcontentsline{toc}{subsubsection}{8.2 Dilatación Térmica}

\begin{itemize}
  \item \textbf{Dilatación Lineal:}
  \[
    \Delta L = \alpha L_0 \Delta T \implies L(T) = L_0 (1 + \alpha \Delta T)
  \]
  \item \textbf{Dilatación Superficial ($\gamma \approx 2\alpha$):}
  \[
    \Delta A = \gamma A_0 \Delta T \implies A(T) = A_0 (1 + \gamma \Delta T)
  \]
  \item \textbf{Dilatación Volumétrica ($\beta \approx 3\alpha$ para sólidos isótropos):}
  \[
    \Delta V = \beta V_0 \Delta T \implies V(T) = V_0 (1 + \beta \Delta T)
  \]
\end{itemize}

\subsubsection*{8.3 Calorimetría y Capacidad Calorífica}
\addcontentsline{toc}{subsubsection}{8.3 Calorimetría y Capacidad Calorífica}

\begin{itemize}
  \item \textbf{Calor sensible:}
  \[
    Q = m c \Delta T = n C_m \Delta T = C \Delta T
  \]
  \item \textbf{Calor latente en cambio de fase:}
  \[
    Q = m L
  \]
  \item \textbf{Principio de Conservación de la Energía (Equilibrio térmico):}
  \[
    \sum Q_{\text{ganado}} + \sum Q_{\text{perdido}} = 0 \implies \sum m_i c_i (T_{\text{eq}} - T_i) = 0
  \]
\end{itemize}

\subsubsection*{8.4 Gas Ideal y Leyes Empíricas}
\addcontentsline{toc}{subsubsection}{8.4 Gas Ideal y Leyes Empíricas}

\begin{itemize}
  \item \textbf{Ecuación de Estado del Gas Ideal:}
  \[
    P V = n R T = N k_B T, \quad R = N_A k_B \approx 8.314 \, \frac{\text{J}}{\text{mol}\cdot\text{K}}
  \]
  \item \textbf{Ley de Boyle-Mariotte ($T = \text{constante}$):}
  \[
    P_1 V_1 = P_2 V_2
  \]
  \item \textbf{Ley de Charles ($P = \text{constante}$):}
  \[
    \frac{V_1}{T_1} = \frac{V_2}{T_2}
  \]
  \item \textbf{Ley de Gay-Lussac ($V = \text{constante}$):}
  \[
    \frac{P_1}{T_1} = \frac{P_2}{T_2}
  \]
  \item \textbf{Ley Combinada:}
  \[
    \frac{P_1 V_1}{T_1} = \frac{P_2 V_2}{T_2}
  \]
\end{itemize}

\subsection*{Nivel Intermedio}
\addcontentsline{toc}{subsection}{Nivel Intermedio}

\subsubsection*{8.5 Mezclas de Gases y Propiedades Molares}
\addcontentsline{toc}{subsubsection}{8.5 Mezclas de Gases y Propiedades Molares}

\begin{itemize}
  \item \textbf{Ley de Dalton de las Presiones Parciales:}
  \[
    P_{\text{total}} = \sum_{i=1}^k P_i, \quad P_i = x_i P_{\text{total}}, \quad x_i = \frac{n_i}{n_{\text{total}}}
  \]
  \item \textbf{Ley de Amagat de los Volúmenes Parciales:}
  \[
    V_{\text{total}} = \sum_{i=1}^k V_i, \quad V_i = x_i V_{\text{total}}
  \]
  \item \textbf{Masa molar aparente de una mezcla:}
  \[
    \bar{M} = \sum_{i=1}^k x_i M_i
  \]
\end{itemize}

\subsubsection*{8.6 Coeficientes Termoelásticos}
\addcontentsline{toc}{subsubsection}{8.6 Coeficientes Termoelásticos}

\begin{itemize}
  \item \textbf{Coeficiente de Dilatación Térmica Isobárica ($\alpha$ o $\beta$):}
  \[
    \alpha = \frac{1}{V} \left( \frac{\partial V}{\partial T} \right)_P
  \]
  \item \textbf{Compresibilidad Isotérmica ($\kappa_T$ o $\beta_T$):}
  \[
    \kappa_T = -\frac{1}{V} \left( \frac{\partial V}{\partial P} \right)_T
  \]
  \item \textbf{Compresibilidad Adiabática / Isentrópica ($\kappa_S$):}
  \[
    \kappa_S = -\frac{1}{V} \left( \frac{\partial V}{\partial P} \right)_S
  \]
  \item \textbf{Coeficiente Piezométrico / Tensión Térmica ($\beta_P$):}
  \[
    \beta_P = \frac{1}{P} \left( \frac{\partial P}{\partial T} \right)_V = \frac{\alpha}{P \kappa_T}
  \]
\end{itemize}

\subsubsection*{8.7 Ecuaciones de Estado para Gases Reales}
\addcontentsline{toc}{subsubsection}{8.7 Ecuaciones de Estado para Gases Reales}

\begin{itemize}
  \item \textbf{Ecuación de Van der Waals:}
  \[
    \left( P + \frac{a n^2}{V^2} \right)(V - n b) = n R T \iff \left( P + \frac{a}{V_m^2} \right)(V_m - b) = R T
  \]
  \item \textbf{Constantes Críticas de Van der Waals:}
  \[
    T_c = \frac{8a}{27Rb}, \quad P_c = \frac{a}{27b^2}, \quad V_{m,c} = 3b
  \]
  \item \textbf{Factor de Compresibilidad Crítico de Van der Waals:}
  \[
    Z_c = \frac{P_c V_{m,c}}{R T_c} = \frac{3}{8} = 0.375
  \]
  \item \textbf{Ecuación Reducida de Van der Waals ($P_r = P/P_c, \, T_r = T/T_c, \, V_r = V_m/V_{m,c}$):}
  \[
    \left( P_r + \frac{3}{V_r^2} \right)(3V_r - 1) = 8T_r
  \]
  \item \textbf{Factor de Compresibilidad y Expansión Virial:}
  \[
    Z = \frac{P V_m}{R T} = 1 + \frac{B(T)}{V_m} + \frac{C(T)}{V_m^2} + \dots = 1 + B'(T)P + C'(T)P^2 + \dots
  \]
\end{itemize}

\subsection*{Nivel Universitario / Avanzado}
\addcontentsline{toc}{subsection}{Nivel Universitario / Avanzado}

\subsubsection*{8.8 Otras Ecuaciones de Estado Reales}
\addcontentsline{toc}{subsubsection}{8.8 Otras Ecuaciones de Estado Reales}

\begin{itemize}
  \item \textbf{Redlich-Kwong:}
  \[
    P = \frac{R T}{V_m - b} - \frac{a}{\sqrt{T} V_m (V_m + b)}
  \]
  \item \textbf{Peng-Robinson:}
  \[
    P = \frac{R T}{V_m - b} - \frac{a(T)}{V_m^2 + 2bV_m - b^2}
  \]
  \item \textbf{Dieterici:}
  \[
    P = \frac{R T}{V_m - b} \exp\left( -\frac{a}{R T V_m} \right)
  \]
\end{itemize}

\subsubsection*{8.9 Teoría Cinética Molecular de los Gases}
\addcontentsline{toc}{subsubsection}{8.9 Teoría Cinética Molecular de los Gases}

\begin{itemize}
  \item \textbf{Presión Cinética:}
  \[
    P = \frac{1}{3} \rho \langle v^2 \rangle = \frac{1}{3} \frac{N m}{V} v_{\text{rms}}^2
  \]
  \item \textbf{Energía Cinética Media Translacional por Molécula:}
  \[
    \langle \epsilon_k \rangle = \frac{1}{2} m \langle v^2 \rangle = \frac{3}{2} k_B T
  \]
  \item \textbf{Velocidades Moleculares Características:}
  \[
    v_{\text{rms}} = \sqrt{\langle v^2 \rangle} = \sqrt{\frac{3 k_B T}{m}} = \sqrt{\frac{3 R T}{M}}
  \]
  \[
    v_{\text{prom}} = \langle v \rangle = \sqrt{\frac{8 k_B T}{\pi m}} = \sqrt{\frac{8 R T}{\pi M}}
  \]
  \[
    v_{\text{mp}} = \sqrt{\frac{2 k_B T}{m}} = \sqrt{\frac{2 R T}{M}}
  \]
  \item \textbf{Distribución de Rapideces de Maxwell-Boltzmann:}
  \[
    f(v) = 4\pi \left( \frac{m}{2\pi k_B T} \right)^{3/2} v^2 \exp\left( -\frac{m v^2}{2 k_B T} \right)
  \]
  \item \textbf{Teorema de Equipartición de la Energía ($f = \text{grados de libertad activos}$):}
  \[
    U = \frac{f}{2} N k_B T = \frac{f}{2} n R T, \quad C_{v,m} = \frac{f}{2} R, \quad C_{p,m} = \left(\frac{f}{2} + 1\right) R
  \]
\end{itemize}

\section*{9. Primera y Segunda Ley de la Termodinámica}
\addcontentsline{toc}{section}{9. Primera y Segunda Ley de la Termodinámica}

\subsection*{Nivel Básico}
\addcontentsline{toc}{subsection}{Nivel Básico}

\subsubsection*{9.1 Primera Ley de la Termodinámica (Sistemas Cerrados)}
\addcontentsline{toc}{subsubsection}{9.1 Primera Ley de la Termodinámica (Sistemas Cerrados)}

\begin{itemize}
  \item \textbf{Forma Integral (Convención física: trabajo $W$ realizado por el sistema):}
  \[
    \Delta U = Q - W
  \]
  \item \textbf{Forma Diferencial General (Sistema cerrado):}
  \[
    dU = \delta Q - \delta W
  \]
  \item \textbf{Para procesos cuasiestáticos de expansión/compresión hidrostática ($\delta W = P \, dV$):}
  \[
    dU = \delta Q - P \, dV
  \]
  \item \textbf{Trabajo Cuasiestático de Expansión/Compresión:}
  \[
    W = \int_{V_i}^{V_f} P \, dV
  \]
\end{itemize}

\subsubsection*{9.2 Procesos Termodinámicos Cuasiestáticos en Gases Ideales}
\addcontentsline{toc}{subsubsection}{9.2 Procesos Termodinámicos Cuasiestáticos en Gases Ideales}

\begin{itemize}
  \item \textbf{Proceso Isocórico ($V = \text{constante}, \, dV = 0$):}
  \[
    W = 0, \quad Q = \Delta U = n C_v \Delta T
  \]
  \item \textbf{Proceso Isobárico ($P = \text{constante}$):}
  \[
    W = P(V_f - V_i) = n R \Delta T, \quad Q = n C_p \Delta T, \quad \Delta U = n C_v \Delta T
  \]
  \item \textbf{Proceso Isotérmico ($T = \text{constante}, \, \Delta T = 0$):}
  \[
    \Delta U = 0 \implies Q = W = n R T \ln\left(\frac{V_f}{V_i}\right) = n R T \ln\left(\frac{P_i}{P_f}\right)
  \]
  \item \textbf{Proceso Adiabático Reversible ($Q = 0, \, \delta Q = 0$):}
  \[
    Q = 0 \implies \Delta U = -W = n C_v (T_f - T_i) = \frac{P_f V_f - P_i V_i}{1 - \gamma}
  \]
  \[
    P V^\gamma = \text{constante}, \quad T V^{\gamma - 1} = \text{constante}, \quad T^\gamma P^{1 - \gamma} = \text{constante}
  \]
  \item \textbf{Relación de Mayer y Coeficiente de Dilatación Adiabática:}
  \[
    C_p - C_v = R, \quad \gamma = \frac{C_p}{C_v}
  \]
\end{itemize}

\subsection*{Nivel Intermedio}
\addcontentsline{toc}{subsection}{Nivel Intermedio}

\subsubsection*{9.3 Procesos Politrópicos ($P V^n = \text{constante}$)}
\addcontentsline{toc}{subsubsection}{9.3 Procesos Politrópicos ($P V^n = \text{constante}$)}

\begin{itemize}
  \item \textbf{Trabajo:}
  \[
    W = \frac{P_f V_f - P_i V_i}{1 - n} = \frac{n R (T_f - T_i)}{1 - n} \quad (n \neq 1)
  \]
  \item \textbf{Capacidad Calorífica Politrópica:}
  \[
    C_n = C_v + \frac{R}{1 - n} = C_v \left( \frac{n - \gamma}{n - 1} \right)
  \]
\end{itemize}

\subsubsection*{9.4 Segunda Ley de la Termodinámica y Entropía}
\addcontentsline{toc}{subsubsection}{9.4 Segunda Ley de la Termodinámica y Entropía}

\begin{itemize}
  \item \textbf{Definición de Entropía (Clausius):}
  \[
    dS = \frac{\delta Q_{\text{rev}}}{T} \implies \Delta S = \int_i^f \frac{\delta Q_{\text{rev}}}{T}
  \]
  \item \textbf{Desigualdad de Clausius:}
  \[
    \oint \frac{\delta Q}{T} \le 0 \quad (\text{igualdad para ciclos reversibles})
  \]
  \item \textbf{Principio del Incremento de Entropía:}
  \[
    \Delta S_{\text{universo}} = \Delta S_{\text{sistema}} + \Delta S_{\text{entorno}} \ge 0
  \]
  \item \textbf{Variación de Entropía en un Gas Ideal:}
  \[
    \Delta S = n C_v \ln\left( \frac{T_f}{T_i} \right) + n R \ln\left( \frac{V_f}{V_i} \right) = n C_p \ln\left( \frac{T_f}{T_i} \right) - n R \ln\left( \frac{P_f}{P_i} \right)
  \]
  \item \textbf{Variación de Entropía en Sustancias Incompresibles (Líquidos/Sólidos):}
  \[
    \Delta S = m c \ln\left( \frac{T_f}{T_i} \right)
  \]
  \item \textbf{Entropía de Mezcla de Gases Ideales:}
  \[
    \Delta S_{\text{mezcla}} = -n_{\text{total}} R \sum_{i=1}^k x_i \ln(x_i)
  \]
\end{itemize}

\subsubsection*{9.5 Máquinas Térmicas, Refrigeradores y Ciclos de Potencia}
\addcontentsline{toc}{subsubsection}{9.5 Máquinas Térmicas, Refrigeradores y Ciclos de Potencia}

\begin{itemize}
  \item \textbf{Eficiencia Térmica General de una Máquina Térmica:}
  \[
    \eta = \frac{W_{\text{neto}}}{Q_H} = \frac{Q_H - |Q_C|}{Q_H} = 1 - \frac{|Q_C|}{Q_H}
  \]
  \item \textbf{Eficiencia de Carnot (Límite máximo entre dos focos):}
  \[
    \eta_{\text{Carnot}} = 1 - \frac{T_C}{T_H}
  \]
  \item \textbf{Coeficiente de Desempeño de Refrigeradores ($\text{COP}_R$):}
  \[
    \text{COP}_R = \beta = \frac{Q_C}{W_{\text{neto}}} = \frac{Q_C}{Q_H - Q_C} \implies \text{COP}_{R,\text{Carnot}} = \frac{T_C}{T_H - T_C}
  \]
  \item \textbf{Coeficiente de Desempeño de Bombas de Calor ($\text{COP}_{\text{HP}}$):}
  \[
    \text{COP}_{\text{HP}} = \gamma' = \frac{Q_H}{W_{\text{neto}}} = \text{COP}_R + 1 \implies \text{COP}_{\text{HP},\text{Carnot}} = \frac{T_H}{T_H - T_C}
  \]
  \item \textbf{Ciclo Otto de Aire Estándar (Relación de compresión $r = V_{\text{máx}}/V_{\text{mín}}$):}
  \[
    \eta_{\text{Otto}} = 1 - \frac{1}{r^{\gamma - 1}}
  \]
  \item \textbf{Ciclo Diesel (Relación de corte de admisión $r_c = V_3/V_2$):}
  \[
    \eta_{\text{Diesel}} = 1 - \frac{1}{r^{\gamma - 1}} \left[ \frac{r_c^\gamma - 1}{\gamma(r_c - 1)} \right]
  \]
  \item \textbf{Ciclo Brayton / Joule (Relación de presiones $r_p = P_2/P_1$):}
  \[
    \eta_{\text{Brayton}} = 1 - \frac{1}{r_p^{(\gamma - 1)/\gamma}}
  \]
\end{itemize}

\subsection*{Nivel Universitario / Avanzado}
\addcontentsline{toc}{subsection}{Nivel Universitario / Avanzado}

\subsubsection*{9.6 Balances de Energía y Entropía en Volúmenes de Control (Sistemas Abiertos)}
\addcontentsline{toc}{subsubsection}{9.6 Balances de Energía y Entropía en Volúmenes de Control (Sistemas Abiertos)}

\begin{itemize}
  \item \textbf{Conservación de Masa:}
  \[
    \frac{dm_{\text{VC}}}{dt} = \sum_{\text{ent}} \dot{m}_i - \sum_{\text{sal}} \dot{m}_e
  \]
  \item \textbf{Primera Ley en Volumen de Control:}
  \[
    \frac{dE_{\text{VC}}}{dt} = \dot{Q} - \dot{W}_{\text{eje}} + \sum_{\text{ent}} \dot{m}_i \left( h_i + \frac{v_i^2}{2} + g z_i \right) - \sum_{\text{sal}} \dot{m}_e \left( h_e + \frac{v_e^2}{2} + g z_e \right)
  \]
  \item \textbf{Primera Ley en Estado Estacionario ($\frac{dE_{\text{VC}}}{dt} = 0$, 1 entrada y 1 salida):}
  \[
    q - w_{\text{eje}} = (h_e - h_i) + \frac{v_e^2 - v_i^2}{2} + g(z_e - z_i)
  \]
  \item \textbf{Balance de Entropía en Volumen de Control:}
  \[
    \frac{dS_{\text{VC}}}{dt} = \sum_k \frac{\dot{Q}_k}{T_k} + \sum_{\text{ent}} \dot{m}_i s_i - \sum_{\text{sal}} \dot{m}_e s_e + \dot{S}_{\text{gen}}, \quad \dot{S}_{\text{gen}} \ge 0
  \]
\end{itemize}

\subsubsection*{9.7 Análisis Exergético (Disponibilidad) y Destrucción de Exergía}
\addcontentsline{toc}{subsubsection}{9.7 Análisis Exergético (Disponibilidad) y Destrucción de Exergía}

\begin{itemize}
  \item \textbf{Exergía de Sistema Cerrado:}
  \[
    \Phi = (U - U_0) + P_0(V - V_0) - T_0(S - S_0)
  \]
  \item \textbf{Exergía de Flujo (Sistema Abierto):}
  \[
    \psi = (h - h_0) - T_0(s - s_0) + \frac{v^2}{2} + gz
  \]
  \item \textbf{Teorema de Gouy-Stodola (Exergía Destruida / Pérdida de Trabajo Útil):}
  \[
    \dot{X}_{\text{destruida}} = I = T_0 \dot{S}_{\text{gen}}
  \]
\end{itemize}

\section*{10. Potenciales Termodinámicos y Relaciones de Maxwell}
\addcontentsline{toc}{section}{10. Potenciales Termodinámicos y Relaciones de Maxwell}

\subsection*{Nivel Intermedio}
\addcontentsline{toc}{subsection}{Nivel Intermedio}

\subsubsection*{10.1 Ecuaciones Fundamentales y Potenciales Termodinámicos}
\addcontentsline{toc}{subsubsection}{10.1 Ecuaciones Fundamentales y Potenciales Termodinámicos}

\begin{itemize}
  \item \textbf{Energía Interna $U(S, V, N)$:}
  \[
    dU = T dS - P dV + \sum_{i=1}^k \mu_i dN_i
  \]
  \item \textbf{Entalpía $H(S, P, N) = U + PV$:}
  \[
    dH = T dS + V dP + \sum_{i=1}^k \mu_i dN_i
  \]
  \item \textbf{Energía Libre de Helmholtz $A(T, V, N) = F = U - TS$:}
  \[
    dA = -S dT - P dV + \sum_{i=1}^k \mu_i dN_i
  \]
  \item \textbf{Energía Libre de Gibbs $G(T, P, N) = H - TS$:}
  \[
    dG = -S dT + V dP + \sum_{i=1}^k \mu_i dN_i
  \]
  \item \textbf{Gran Potencial Termodinámico / Potencial Gran Canónico ($\Omega = \Phi_G(T, V, \mu) = U - TS - \mu N = -PV$):}
  \[
    d\Omega = -S dT - P dV - N d\mu
  \]
\end{itemize}

\subsubsection*{10.2 Relaciones de Maxwell}
\addcontentsline{toc}{subsubsection}{10.2 Relaciones de Maxwell}

\begin{itemize}
  \item \textbf{A partir de $dU$:}
  \[
    \left( \frac{\partial T}{\partial V} \right)_S = -\left( \frac{\partial P}{\partial S} \right)_V
  \]
  \item \textbf{A partir de $dH$:}
  \[
    \left( \frac{\partial T}{\partial P} \right)_S = \left( \frac{\partial V}{\partial S} \right)_P
  \]
  \item \textbf{A partir de $dA$:}
  \[
    \left( \frac{\partial S}{\partial V} \right)_T = \left( \frac{\partial P}{\partial T} \right)_V
  \]
  \item \textbf{A partir de $dG$:}
  \[
    \left( \frac{\partial S}{\partial P} \right)_T = -\left( \frac{\partial V}{\partial T} \right)_P
  \]
\end{itemize}

\subsubsection*{10.3 Ecuaciones $T dS$ y Propiedades de Calores Específicos}
\addcontentsline{toc}{subsubsection}{10.3 Ecuaciones $T dS$ y Propiedades de Calores Específicos}

\begin{itemize}
  \item \textbf{Primera Ecuación $T dS$:}
  \[
    T dS = C_v dT + T \left( \frac{\partial P}{\partial T} \right)_V dV = C_v dT + \frac{T \alpha}{\kappa_T} dV
  \]
  \item \textbf{Segunda Ecuación $T dS$:}
  \[
    T dS = C_p dT - T \left( \frac{\partial V}{\partial T} \right)_P dP = C_p dT - T V \alpha dP
  \]
  \item \textbf{Relación General $C_p - C_v$:}
  \[
    C_p - C_v = T \left( \frac{\partial P}{\partial T} \right)_V \left( \frac{\partial V}{\partial T} \right)_P = \frac{T V \alpha^2}{\kappa_T}
  \]
  \item \textbf{Relación de Compresibilidades:}
  \[
    \frac{C_p}{C_v} = \frac{\kappa_T}{\kappa_S} = \gamma
  \]
\end{itemize}

\subsection*{Nivel Universitario / Avanzado}
\addcontentsline{toc}{subsection}{Nivel Universitario / Avanzado}

\subsubsection*{10.4 Ecuaciones de Estado de la Energía Interna y Efecto Joule-Thomson}
\addcontentsline{toc}{subsubsection}{10.4 Ecuaciones de Estado de la Energía Interna y Efecto Joule-Thomson}

\begin{itemize}
  \item \textbf{Primera Ecuación de Estado de la Energía:}
  \[
    \left( \frac{\partial U}{\partial V} \right)_T = T \left( \frac{\partial P}{\partial T} \right)_V - P = \frac{T \alpha}{\kappa_T} - P
  \]
  \item \textbf{Segunda Ecuación de Estado de la Energía (Entalpía):}
  \[
    \left( \frac{\partial H}{\partial P} \right)_T = V - T \left( \frac{\partial V}{\partial T} \right)_P = V(1 - T\alpha)
  \]
  \item \textbf{Coeficiente de Joule-Thomson ($\mu_{\text{JT}}$):}
  \[
    \mu_{\text{JT}} = \left( \frac{\partial T}{\partial P} \right)_H = \frac{1}{C_p} \left[ T \left( \frac{\partial V}{\partial T} \right)_P - V \right] = \frac{V}{C_p}(T\alpha - 1)
  \]
  \item \textbf{Temperatura de Inversión de Joule-Thomson ($\mu_{\text{JT}} = 0$):}
  \[
    T_{\text{inv}} = \frac{V}{\left( \frac{\partial V}{\partial T} \right)_P} = \frac{1}{\alpha}
  \]
\end{itemize}

\subsubsection*{10.5 Ecuación de Gibbs-Duhem y Propiedades Molares Parciales}
\addcontentsline{toc}{subsubsection}{10.5 Ecuación de Gibbs-Duhem y Propiedades Molares Parciales}

\begin{itemize}
  \item \textbf{Ecuación de Gibbs-Duhem:}
  \[
    S dT - V dP + \sum_{i=1}^k N_i d\mu_i = 0 \implies \sum_{i=1}^k x_i d\mu_i = 0 \quad (\text{a } T, P \text{ ctes.})
  \]
  \item \textbf{Propiedad Molar Parcial general ($\bar{M}_i$):}
  \[
    \bar{M}_i = \left( \frac{\partial M}{\partial n_i} \right)_{T, P, n_{j \neq i}} \implies M = \sum_{i=1}^k n_i \bar{M}_i
  \]
  \item \textbf{Potencial Químico como Propiedad Molar Parcial:}
  \[
    \mu_i = \bar{G}_i = \left( \frac{\partial G}{\partial n_i} \right)_{T, P, n_{j \neq i}}
  \]
\end{itemize}

\subsubsection*{10.6 Transiciones de Fase y Equilibrio Químico}
\addcontentsline{toc}{subsubsection}{10.6 Transiciones de Fase y Equilibrio Químico}

\begin{itemize}
  \item \textbf{Ecuación de Clapeyron (Cualquier transición de fase de 1er orden):}
  \[
    \frac{dP}{dT} = \frac{\Delta s}{\Delta v} = \frac{\Delta h}{T \Delta v}
  \]
  \item \textbf{Ecuación de Clausius-Clapeyron (Vaporización/Sublimación con $v_g \gg v_l$ y gas ideal):}
  \[
    \frac{d \ln P}{dT} = \frac{\Delta h_{\text{vap}}}{R T^2} \implies \ln\left( \frac{P_2}{P_1} \right) = -\frac{\Delta h_{\text{vap}}}{R} \left( \frac{1}{T_2} - \frac{1}{T_1} \right)
  \]
  \item \textbf{Regla de las Fases de Gibbs:}
  \[
    F = C - \mathcal{P} + 2
  \]
  \[
    (F = \text{grados de libertad}, \, C = \text{número de componentes}, \, \mathcal{P} = \text{número de fases})
  \]
  \item \textbf{Ecuaciones de Ehrenfest (Transiciones de fase de 2do orden):}
  \[
    \frac{dP}{dT} = \frac{\Delta C_p}{T V \Delta \alpha} = \frac{\Delta \alpha}{\Delta \kappa_T}
  \]
  \item \textbf{Isoterma de Van 't Hoff para Equilibrio Químico:}
  \[
    \Delta G^\circ = -R T \ln(K_p) \implies \frac{d \ln K_p}{dT} = \frac{\Delta H^\circ}{R T^2}
  \]
\end{itemize}

\subsubsection*{10.7 Tercera Ley de la Termodinámica (Teorema del Calor de Nernst)}
\addcontentsline{toc}{subsubsection}{10.7 Tercera Ley de la Termodinámica (Teorema del Calor de Nernst)}

\begin{itemize}
  \item \textbf{Comportamiento asíntotico en el cero absoluto:}
  \[
    \lim_{T \to 0} S = 0 \quad (\text{para un cristal perfecto})
  \]
  \[
    \lim_{T \to 0} C_p = \lim_{T \to 0} C_v = 0, \quad \lim_{T \to 0} \alpha = 0, \quad \lim_{T \to 0} \left( \frac{\partial P}{\partial T} \right)_V = 0
  \]
\end{itemize}

\section*{11. Termodinámica Estadística}
\addcontentsline{toc}{section}{11. Termodinámica Estadística}

\subsection*{Nivel Universitario / Avanzado}
\addcontentsline{toc}{subsection}{Nivel Universitario / Avanzado}

\subsubsection*{11.1 Colectividades y Funciones de Partición}
\addcontentsline{toc}{subsubsection}{11.1 Colectividades y Funciones de Partición}

\begin{itemize}
  \item \textbf{Entropía Estadística de Boltzmann:}
  \[
    S = k_B \ln(\Omega)
  \]
  \item \textbf{Entropía de Gibbs (Distribución general):}
  \[
    S = -k_B \sum_i P_i \ln(P_i)
  \]
  \item \textbf{Factor de Boltzmann ($\beta$):}
  \[
    \beta = \frac{1}{k_B T}
  \]
\end{itemize}

\subsubsection*{11.2 Colectividad Canónica (Sistema cerrado a $N, V, T$ fijos)}
\addcontentsline{toc}{subsubsection}{11.2 Colectividad Canónica (Sistema cerrado a $N, V, T$ fijos)}

\begin{itemize}
  \item \textbf{Función de Partición Canónica ($Z$ o $Q$):}
  \[
    Z = \sum_i g_i \exp(-\beta E_i) = \sum_i \exp(-\beta E_i)
  \]
  \item \textbf{Probabilidad de ocupación del microestado $i$:}
  \[
    P_i = \frac{\exp(-\beta E_i)}{Z}
  \]
  \item \textbf{Conexión con los Potenciales Termodinámicos:}
  \item \textbf{Energía Libre de Helmholtz:}
  \[
    F = -k_B T \ln(Z)
  \]
  \item \textbf{Energía Interna:}
  \[
    U = \langle E \rangle = -\frac{\partial \ln Z}{\partial \beta} = k_B T^2 \left( \frac{\partial \ln Z}{\partial T} \right)_{V, N}
  \]
  \item \textbf{Entropía:}
  \[
    S = k_B \ln(Z) + \frac{U}{T} = -\left( \frac{\partial F}{\partial T} \right)_{V, N}
  \]
  \item \textbf{Presión:}
  \[
    P = k_B T \left( \frac{\partial \ln Z}{\partial V} \right)_{T, N} = -\left( \frac{\partial F}{\partial V} \right)_{T, N}
  \]
  \item \textbf{Potencial Químico:}
  \[
    \mu = -k_B T \left( \frac{\partial \ln Z}{\partial N} \right)_{T, V} = \left( \frac{\partial F}{\partial N} \right)_{T, V}
  \]
\end{itemize}

\subsubsection*{11.3 Colectividad Gran Canónica (Sistema abierto a $\mu, V, T$ fijos)}
\addcontentsline{toc}{subsubsection}{11.3 Colectividad Gran Canónica (Sistema abierto a $\mu, V, T$ fijos)}

\begin{itemize}
  \item \textbf{Gran Función de Partición ($\Xi$):}
  \[
    \Xi = \sum_{N=0}^\infty \sum_i \exp\left[ -\beta (E_{i,N} - \mu N) \right] = \sum_{N=0}^\infty \lambda^N Z_N, \quad \lambda = e^{\beta \mu} \quad (\text{fugacidad})
  \]
  \item \textbf{Conexión con el Gran Potencial:}
  \[
    \Phi_G = -P V = -k_B T \ln(\Xi)
  \]
  \item \textbf{Número medio de partículas:}
  \[
    \langle N \rangle = k_B T \left( \frac{\partial \ln \Xi}{\partial \mu} \right)_{T, V} = \frac{1}{\beta} \left( \frac{\partial \ln \Xi}{\partial \mu} \right)
  \]
\end{itemize}

\subsubsection*{11.4 Estadísticas Cuánticas (Ocupación media del estado monoparticular $\epsilon_i$)}
\addcontentsline{toc}{subsubsection}{11.4 Estadísticas Cuánticas (Ocupación media del estado monoparticular $\epsilon_i$)}

\begin{itemize}
  \item \textbf{Estadística de Maxwell-Boltzmann (Partículas clásicas distinguibles):}
  \[
    \langle n_i \rangle_{\text{MB}} = \frac{1}{\exp[\beta(\epsilon_i - \mu)]}
  \]
  \item \textbf{Estadística de Fermi-Dirac (Fermiones, espín semi-entero, Principio de Exclusión):}
  \[
    \langle n_i \rangle_{\text{FD}} = \frac{1}{\exp[\beta(\epsilon_i - \mu)] + 1}
  \]
  \[
    \text{Energía de Fermi en } T = 0\text{ K}: \quad E_F = \frac{\hbar^2}{2m} (3\pi^2 n)^{2/3}, \quad n = \frac{N}{V}
  \]
  \item \textbf{Estadística de Bose-Einstein (Bosones, espín entero):}
  \[
    \langle n_i \rangle_{\text{BE}} = \frac{1}{\exp[\beta(\epsilon_i - \mu)] - 1}, \quad \mu \le \epsilon_0
  \]
  \[
    \text{Temperatura Crítica de Condensación de Bose-Einstein}: \quad T_c = \frac{2\pi \hbar^2}{m k_B} \left( \frac{n}{\zeta(3/2)} \right)^{2/3}, \quad \zeta(3/2) \approx 2.612
  \]
\end{itemize}

\section*{Compendio IV: Electricidad y Magnetismo}
\addcontentsline{toc}{section}{Compendio IV: Electricidad y Magnetismo}

\textit{(Nivel Básico a Universitario)}

\section*{12. Electrostática y Medios Dieléctricos}
\addcontentsline{toc}{section}{12. Electrostática y Medios Dieléctricos}

\subsection*{Nivel Básico}
\addcontentsline{toc}{subsection}{Nivel Básico}

\subsubsection*{12.1 Ley de Coulomb y Fuerza Eléctrica}
\addcontentsline{toc}{subsubsection}{12.1 Ley de Coulomb y Fuerza Eléctrica}

\begin{itemize}
  \item \textbf{Magnitud de la fuerza entre dos cargas puntuales:}
  \[
    F = \frac{1}{4\pi\varepsilon_0} \frac{|q_1 q_2|}{r^2} = k_e \frac{|q_1 q_2|}{r^2}, \quad k_e \approx 8.9875 \times 10^9 \, \frac{\text{N}\cdot\text{m}^2}{\text{C}^2}
  \]
  \item \textbf{Forma vectorial de la fuerza sobre $q_1$ debida a $q_2$:}
  \[
    \vec{F}_{12} = \frac{1}{4\pi\varepsilon_0} \frac{q_1 q_2}{r_{21}^2} \hat{r}_{21}
  \]
\end{itemize}

\subsubsection*{12.2 Campo Eléctrico ($\vec{E}$)}
\addcontentsline{toc}{subsubsection}{12.2 Campo Eléctrico ($\vec{E}$)}

\begin{itemize}
  \item \textbf{Definición por carga de prueba:}
  \[
    \vec{E} = \frac{\vec{F}}{q_0}
  \]
  \item \textbf{Campo de una carga puntual:}
  \[
    \vec{E} = \frac{1}{4\pi\varepsilon_0} \frac{q}{r^2} \hat{r}
  \]
  \item \textbf{Principio de superposición para $N$ cargas:}
  \[
    \vec{E}_{\text{total}} = \sum_{i=1}^N \vec{E}_i = \frac{1}{4\pi\varepsilon_0} \sum_{i=1}^N \frac{q_i}{r_i^2} \hat{r}_i
  \]
\end{itemize}

\subsubsection*{12.3 Potencial y Energía Potencial Electrostática}
\addcontentsline{toc}{subsubsection}{12.3 Potencial y Energía Potencial Electrostática}

\begin{itemize}
  \item \textbf{Potencial eléctrico de una carga puntual:}
  \[
    V(r) = \frac{1}{4\pi\varepsilon_0} \frac{q}{r}
  \]
  \item \textbf{Trabajo y diferencia de potencial:}
  \[
    W_{A \to B} = -q_0 \Delta V = q_0 (V_A - V_B)
  \]
  \item \textbf{Energía potencial electrostática de dos cargas:}
  \[
    U = \frac{1}{4\pi\varepsilon_0} \frac{q_1 q_2}{r}
  \]
\end{itemize}

\subsubsection*{12.4 Capacitancia y Condensadores}
\addcontentsline{toc}{subsubsection}{12.4 Capacitancia y Condensadores}

\begin{itemize}
  \item \textbf{Definición general:}
  \[
    C = \frac{Q}{V}
  \]
  \item \textbf{Condensador de placas plano-paralelas:}
  \[
    C = \frac{\varepsilon_0 A}{d}
  \]
  \item \textbf{Condensador coaxial / cilíndrico (longitud $L$, radios $a < b$):}
  \[
    C = \frac{2\pi\varepsilon_0 L}{\ln(b/a)}
  \]
  \item \textbf{Condensador esférico (radios $a < b$):}
  \[
    C = 4\pi\varepsilon_0 \left( \frac{a b}{b - a} \right)
  \]
  \item \textbf{Asociación de Condensadores:}
  \[
    \text{En Paralelo: } C_{\text{eq}} = \sum_{i=1}^N C_i \quad | \quad \text{En Serie: } \frac{1}{C_{\text{eq}}} = \sum_{i=1}^N \frac{1}{C_i}
  \]
  \item \textbf{Energía almacenada en un condensador:}
  \[
    U = \frac{1}{2} Q V = \frac{1}{2} C V^2 = \frac{Q^2}{2C}
  \]
\end{itemize}

\subsection*{Nivel Intermedio}
\addcontentsline{toc}{subsection}{Nivel Intermedio}

\subsubsection*{12.5 Distribuciones Continuas de Carga}
\addcontentsline{toc}{subsubsection}{12.5 Distribuciones Continuas de Carga}

\begin{itemize}
  \item \textbf{Campo eléctrico para distribución volumétrica continua:}
  \[
    \vec{E}(\vec{r}) = \frac{1}{4\pi\varepsilon_0} \int_V \frac{\rho(\vec{r}')(\vec{r} - \vec{r}')}{|\vec{r} - \vec{r}'|^3} \, dV'
  \]
  \[
    \left( \text{Superficial: } \vec{E}(\vec{r}) = \frac{1}{4\pi\varepsilon_0} \int_S \frac{\sigma(\vec{r}')(\vec{r} - \vec{r}')}{|\vec{r} - \vec{r}'|^3} \, dA', \quad \text{Lineal: } \vec{E}(\vec{r}) = \frac{1}{4\pi\varepsilon_0} \int_C \frac{\lambda(\vec{r}')(\vec{r} - \vec{r}')}{|\vec{r} - \vec{r}'|^3} \, d\ell' \right)
  \]
  \item \textbf{Potencial eléctrico continuo:}
  \[
    V(\vec{r}) = \frac{1}{4\pi\varepsilon_0} \int_V \frac{\rho(\vec{r}')}{|\vec{r} - \vec{r}'|} \, dV'
  \]
\end{itemize}

\subsubsection*{12.6 Ley de Gauss en Forma Integral}
\addcontentsline{toc}{subsubsection}{12.6 Ley de Gauss en Forma Integral}

\begin{itemize}
  \item \textbf{Flujo Eléctrico ($\Phi_E$):}
  \[
    \Phi_E = \iint_S \vec{E} \cdot d\vec{A}
  \]
  \item \textbf{Ley de Gauss:}
  \[
    \Phi_E = \oiint_{\partial V} \vec{E} \cdot d\vec{A} = \frac{Q_{\text{enc}}}{\varepsilon_0}
  \]
\end{itemize}

\subsubsection*{12.7 Relación Diferencial entre Campo y Potencial}
\addcontentsline{toc}{subsubsection}{12.7 Relación Diferencial entre Campo y Potencial}

\begin{itemize}
  \item \textbf{Campo como gradiente del potencial:}
  \[
    \vec{E} = -\nabla V = -\left( \frac{\partial V}{\partial x}\hat{i} + \frac{\partial V}{\partial y}\hat{j} + \frac{\partial V}{\partial z}\hat{k} \right)
  \]
  \item \textbf{Diferencia de potencial por integral de línea:}
  \[
    V_B - V_A = -\int_A^B \vec{E} \cdot d\vec{\ell}
  \]
\end{itemize}

\subsubsection*{12.8 Dipolo Eléctrico}
\addcontentsline{toc}{subsubsection}{12.8 Dipolo Eléctrico}

\begin{itemize}
  \item \textbf{Momento dipolar eléctrico:}
  \[
    \vec{p} = q \vec{d}
  \]
  \item \textbf{Potencial de un dipolo puntual en el origen:}
  \[
    V(r, \theta) = \frac{1}{4\pi\varepsilon_0} \frac{\vec{p} \cdot \hat{r}}{r^2} = \frac{1}{4\pi\varepsilon_0} \frac{p \cos\theta}{r^2}
  \]
  \item \textbf{Campo eléctrico del dipolo:}
  \[
    \vec{E}(r, \theta) = \frac{1}{4\pi\varepsilon_0 r^3} \left[ 3(\vec{p} \cdot \hat{r})\hat{r} - \vec{p} \right] = \frac{p}{4\pi\varepsilon_0 r^3} (2\cos\theta \hat{u}_r + \sen\theta \hat{u}_\theta)
  \]
  \item \textbf{Torque y energía de un dipolo en campo externo:}
  \[
    \vec{\tau} = \vec{p} \times \vec{E}, \quad U = -\vec{p} \cdot \vec{E}, \quad \vec{F} = (\vec{p} \cdot \nabla)\vec{E}
  \]
\end{itemize}

\subsubsection*{12.9 Dieléctricos y Polarización}
\addcontentsline{toc}{subsubsection}{12.9 Dieléctricos y Polarización}

\begin{itemize}
  \item \textbf{Vector Desplazamiento Eléctrico ($\vec{D}$):}
  \[
    \vec{D} = \varepsilon_0 \vec{E} + \vec{P} = \varepsilon \vec{E} = \varepsilon_0 \varepsilon_r \vec{E}
  \]
  \item \textbf{Vector Polarización ($\vec{P}$) y Susceptibilidad Eléctrica ($\chi_e$):}
  \[
    \vec{P} = \varepsilon_0 \chi_e \vec{E}, \quad \varepsilon_r = 1 + \chi_e
  \]
  \item \textbf{Densidades de Carga Ligada / Polarización:}
  \[
    \rho_b = -\nabla \cdot \vec{P}, \quad \sigma_b = \vec{P} \cdot \hat{n}
  \]
  \item \textbf{Densidad volumétrica de energía electrostática:}
  \[
    u_e = \frac{1}{2} \vec{D} \cdot \vec{E} = \frac{1}{2} \varepsilon E^2
  \]
\end{itemize}

\subsection*{Nivel Universitario / Avanzado}
\addcontentsline{toc}{subsection}{Nivel Universitario / Avanzado}

\subsubsection*{12.10 Ecuaciones Diferenciales de la Electrostática}
\addcontentsline{toc}{subsubsection}{12.10 Ecuaciones Diferenciales de la Electrostática}

\begin{itemize}
  \item \textbf{Ley de Gauss Diferencial:}
  \[
    \nabla \cdot \vec{E} = \frac{\rho}{\varepsilon_0} \quad \implies \quad \nabla \cdot \vec{D} = \rho_f
  \]
  \item \textbf{Carácter conservativo:}
  \[
    \nabla \times \vec{E} = \vec{0}
  \]
  \item \textbf{Ecuación de Poisson:}
  \[
    \nabla^2 V = -\frac{\rho}{\varepsilon_0}
  \]
  \item \textbf{Ecuación de Laplace (en regiones libres de carga $\rho = 0$):}
  \[
    \nabla^2 V = 0
  \]
\end{itemize}

\subsubsection*{12.11 Condiciones de Frontera Electrostáticas}
\addcontentsline{toc}{subsubsection}{12.11 Condiciones de Frontera Electrostáticas}

\begin{itemize}
  \item \textbf{Componente tangencial del campo eléctrico:}
  \[
    \hat{n} \times (\vec{E}_1 - \vec{E}_2) = \vec{0} \implies E_{1t} = E_{2t}
  \]
  \item \textbf{Componente normal del desplazamiento eléctrico:}
  \[
    \hat{n} \cdot (\vec{D}_1 - \vec{D}_2) = \sigma_f \implies D_{1n} - D_{2n} = \sigma_f
  \]
\end{itemize}

\subsubsection*{12.12 Expansión Multipolar del Potencial Electrostático}
\addcontentsline{toc}{subsubsection}{12.12 Expansión Multipolar del Potencial Electrostático}

\begin{itemize}
  \item \textbf{Serie de potencias:}
  \[
    V(\vec{r}) = \frac{1}{4\pi\varepsilon_0} \left[ \frac{Q_{\text{total}}}{r} + \frac{\vec{p} \cdot \hat{r}}{r^2} + \frac{1}{2 r^3} \sum_{i,j} Q_{ij} \hat{r}_i \hat{r}_j + \dots \right]
  \]
  \item \textbf{Tensor del momento cuadrupolar ($Q_{ij}$):}
  \[
    Q_{ij} = \int (3 x'_i x'_j - r'^2 \delta_{ij}) \rho(\vec{r}') dV'
  \]
\end{itemize}

\subsubsection*{12.13 Energía Electrostática Total de un Sistema Continuo}
\addcontentsline{toc}{subsubsection}{12.13 Energía Electrostática Total de un Sistema Continuo}

\begin{itemize}
  \item \textbf{En términos de fuentes:}
  \[
    U = \frac{1}{2} \int_V \rho(\vec{r}) V(\vec{r}) dV
  \]
  \item \textbf{En términos de campo:}
  \[
    U = \frac{1}{2} \int_{\text{todo el espacio}} \vec{D} \cdot \vec{E} \, dV = \frac{\varepsilon_0}{2} \int E^2 \, dV
  \]
\end{itemize}

\section*{13. Electrodinámica, Corriente y Circuitos Eléctricos}
\addcontentsline{toc}{section}{13. Electrodinámica, Corriente y Circuitos Eléctricos}

\subsection*{Nivel Básico}
\addcontentsline{toc}{subsection}{Nivel Básico}

\subsubsection*{13.1 Corriente, Resistencia y Ley de Ohm}
\addcontentsline{toc}{subsubsection}{13.1 Corriente, Resistencia y Ley de Ohm}

\begin{itemize}
  \item \textbf{Intensidad de corriente eléctrica:}
  \[
    I = \frac{dQ}{dt}
  \]
  \item \textbf{Ley de Ohm Macroscópica:}
  \[
    V = I R
  \]
  \item \textbf{Resistencia en función de la geometría:}
  \[
    R = \rho_e \frac{L}{A} = \frac{L}{\sigma_c A}
  \]
  \item \textbf{Variación de la resistividad con la temperatura:}
  \[
    \rho_e(T) = \rho_0 [1 + \alpha (T - T_0)]
  \]
\end{itemize}

\subsubsection*{13.2 Potencia Eléctrica y Efecto Joule}
\addcontentsline{toc}{subsubsection}{13.2 Potencia Eléctrica y Efecto Joule}

\begin{itemize}
  \item \textbf{Potencia disipada / consumida:}
  \[
    P = V I = I^2 R = \frac{V^2}{R}
  \]
  \item \textbf{Energía disipada por efecto Joule:}
  \[
    W = \int P \, dt = I^2 R t \quad (\text{para } I \text{ constante})
  \]
\end{itemize}

\subsubsection*{13.3 Fuerza Electromotriz (FEM) y Asociación de Resistencias}
\addcontentsline{toc}{subsubsection}{13.3 Fuerza Electromotriz (FEM) y Asociación de Resistencias}

\begin{itemize}
  \item \textbf{FEM ($\mathcal{E}$) y voltaje terminal con resistencia interna $r$:}
  \[
    V_{\text{terminal}} = \mathcal{E} - I r
  \]
  \item \textbf{Asociación de Resistencias:}
  \[
    \text{En Serie: } R_{\text{eq}} = \sum_{i=1}^N R_i \quad | \quad \text{En Paralelo: } \frac{1}{R_{\text{eq}}} = \sum_{i=1}^N \frac{1}{R_i}
  \]
\end{itemize}

\subsection*{Nivel Intermedio}
\addcontentsline{toc}{subsection}{Nivel Intermedio}

\subsubsection*{13.4 Densidad de Corriente y Ley de Ohm Microscópica}
\addcontentsline{toc}{subsubsection}{13.4 Densidad de Corriente y Ley de Ohm Microscópica}

\begin{itemize}
  \item \textbf{Densidad de corriente ($\vec{J}$):}
  \[
    \vec{J} = n q \vec{v}_d, \quad I = \iint_S \vec{J} \cdot d\vec{A}
  \]
  \item \textbf{Ley de Ohm Puntual / Microscópica:}
  \[
    \vec{J} = \sigma_c \vec{E} = \frac{1}{\rho_e} \vec{E}
  \]
  \item \textbf{Modelo de Drude para la conductividad:}
  \[
    \sigma_c = \frac{n q^2 \tau_c}{m_e}
  \]
\end{itemize}

\subsubsection*{13.5 Leyes de Kirchhoff}
\addcontentsline{toc}{subsubsection}{13.5 Leyes de Kirchhoff}

\begin{itemize}
  \item \textbf{Ley de Corrientes de Kirchhoff (Nodos - Conservación de carga):}
  \[
    \sum_{k=1}^n I_k = 0
  \]
  \item \textbf{Ley de Voltajes de Kirchhoff (Mallas - Conservación de energía):}
  \[
    \sum_{k=1}^m V_k = \sum_{k=1}^m \mathcal{E}_k - \sum_{k=1}^m I_k R_k = 0
  \]
\end{itemize}

\subsubsection*{13.6 Circuitos Transitorios de Primer Orden en Corriente Continua (DC)}
\addcontentsline{toc}{subsubsection}{13.6 Circuitos Transitorios de Primer Orden en Corriente Continua (DC)}

\begin{itemize}
  \item \textbf{Circuito RC - Carga del Condensador ($\tau = RC$):}
  \[
    q(t) = C\mathcal{E} (1 - e^{-t/\tau}), \quad i(t) = \frac{\mathcal{E}}{R} e^{-t/\tau}, \quad V_C(t) = \mathcal{E}(1 - e^{-t/\tau})
  \]
  \item \textbf{Circuito RC - Descarga del Condensador:}
  \[
    q(t) = Q_0 e^{-t/\tau}, \quad i(t) = -\frac{Q_0}{\tau} e^{-t/\tau}, \quad V_C(t) = V_0 e^{-t/\tau}
  \]
  \item \textbf{Circuito RL - Crecimiento de Corriente ($\tau = L/R$):}
  \[
    i(t) = \frac{\mathcal{E}}{R} (1 - e^{-t/\tau}), \quad V_L(t) = \mathcal{E} e^{-t/\tau}
  \]
  \item \textbf{Circuito RL - Decaimiento de Corriente:}
  \[
    i(t) = I_0 e^{-t/\tau}
  \]
\end{itemize}

\subsection*{Nivel Universitario / Avanzado}
\addcontentsline{toc}{subsection}{Nivel Universitario / Avanzado}

\subsubsection*{13.7 Ecuación de Continuidad de la Carga}
\addcontentsline{toc}{subsubsection}{13.7 Ecuación de Continuidad de la Carga}

\begin{itemize}
  \item \textbf{Conservación local de la carga:}
  \[
    \nabla \cdot \vec{J} + \frac{\partial \rho}{\partial t} = 0
  \]
  \item \textbf{Régimen Estacionario ($\partial\rho/\partial t = 0$):}
  \[
    \nabla \cdot \vec{J} = 0 \implies \oiint_S \vec{J} \cdot d\vec{A} = 0
  \]
\end{itemize}

\subsubsection*{13.8 Análisis de Circuitos en Régimen Senoidal Permanente (AC - Fasores)}
\addcontentsline{toc}{subsubsection}{13.8 Análisis de Circuitos en Régimen Senoidal Permanente (AC - Fasores)}

\begin{itemize}
  \item \textbf{Representación fasorial de tensiones y corrientes:}
  \[
    v(t) = V_m \cos(\omega t + \phi) \iff \mathbf{V} = V_{\text{rms}} e^{j\phi} = \frac{V_m}{\sqrt{2}} \angle \phi
  \]
  \item \textbf{Impedancia Compleja ($Z = R + jX$):}
  \[
    Z_R = R, \quad Z_L = j\omega L = \omega L \angle 90^\circ, \quad Z_C = \frac{1}{j\omega C} = -\frac{j}{\omega C} = \frac{1}{\omega C} \angle -90^\circ
  \]
  \item \textbf{Ley de Ohm Fasorial:}
  \[
    \mathbf{V} = \mathbf{I} \mathbf{Z}
  \]
  \item \textbf{Potencia Compleja ($\mathbf{S}$):}
  \[
    \mathbf{S} = \mathbf{V}_{\text{rms}} \mathbf{I}_{\text{rms}}^* = P + j Q = |\mathbf{S}| \cos\theta + j |\mathbf{S}| \sen\theta
  \]
  \[
    P = V_{\text{rms}} I_{\text{rms}} \cos(\theta_v - \theta_i) \quad (\text{Potencia Activa, W})
  \]
  \[
    Q = V_{\text{rms}} I_{\text{rms}} \sen(\theta_v - \theta_i) \quad (\text{Potencia Reactiva, VAR})
  \]
  \[
    |\mathbf{S}| = V_{\text{rms}} I_{\text{rms}} \quad (\text{Potencia Aparente, VA})
  \]
  \item \textbf{Factor de Potencia:}
  \[
    \text{FP} = \cos(\theta_v - \theta_i) = \frac{P}{|\mathbf{S}|}
  \]
\end{itemize}

\section*{14. Magnetostática y Materia Magnética}
\addcontentsline{toc}{section}{14. Magnetostática y Materia Magnética}

\subsection*{Nivel Básico}
\addcontentsline{toc}{subsection}{Nivel Básico}

\subsubsection*{14.1 Fuerza Magnética y Fuerza de Lorentz}
\addcontentsline{toc}{subsubsection}{14.1 Fuerza Magnética y Fuerza de Lorentz}

\begin{itemize}
  \item \textbf{Fuerza magnética sobre una carga en movimiento:}
  \[
    \vec{F}_B = q (\vec{v} \times \vec{B}) \implies F_B = |q| v B \sen\theta
  \]
  \item \textbf{Fuerza de Lorentz completa:}
  \[
    \vec{F} = q (\vec{E} + \vec{v} \times \vec{B})
  \]
  \item \textbf{Movimiento ciclotrónico (carga perpendicular a campo uniforme):}
  \[
    r = \frac{m v}{|q| B}, \quad \omega_c = \frac{|q| B}{m}, \quad T = \frac{2\pi m}{|q| B}
  \]
\end{itemize}

\subsubsection*{14.2 Fuerza Magnética sobre Conductores}
\addcontentsline{toc}{subsubsection}{14.2 Fuerza Magnética sobre Conductores}

\begin{itemize}
  \item \textbf{Conductor rectilíneo con corriente $I$:}
  \[
    \vec{F} = I (\vec{L} \times \vec{B})
  \]
  \item \textbf{Conductor de geometría arbitraria:}
  \[
    \vec{F} = \int_C I (d\vec{\ell} \times \vec{B})
  \]
  \item \textbf{Fuerza magnética entre dos conductores paralelos largos:}
  \[
    \frac{F}{L} = \frac{\mu_0 I_1 I_2}{2\pi d}, \quad \mu_0 = 4\pi \times 10^{-7} \, \frac{\text{T}\cdot\text{m}}{\text{A}}
  \]
\end{itemize}

\subsubsection*{14.3 Dipolo Magnético}
\addcontentsline{toc}{subsubsection}{14.3 Dipolo Magnético}

\begin{itemize}
  \item \textbf{Momento dipolar magnético de una espira plana:}
  \[
    \vec{\mu} = N I A \hat{n}
  \]
  \item \textbf{Torque sobre una espira en campo magnético:}
  \[
    \vec{\tau} = \vec{\mu} \times \vec{B}
  \]
  \item \textbf{Energía potencial magnética del dipolo:}
  \[
    U = -\vec{\mu} \cdot \vec{B}
  \]
\end{itemize}

\subsection*{Nivel Intermedio}
\addcontentsline{toc}{subsection}{Nivel Intermedio}

\subsubsection*{14.4 Ley de Biot-Savart}
\addcontentsline{toc}{subsubsection}{14.4 Ley de Biot-Savart}

\begin{itemize}
  \item \textbf{Campo magnético de un elemento de corriente:}
  \[
    d\vec{B} = \frac{\mu_0 I}{4\pi} \frac{d\vec{\ell} \times \hat{r}}{r^2} = \frac{\mu_0 I}{4\pi} \frac{d\vec{\ell} \times (\vec{r} - \vec{r}')}{|\vec{r} - \vec{r}'|^3}
  \]
  \item \textbf{Campo de un conductor rectilíneo infinito:}
  \[
    B = \frac{\mu_0 I}{2\pi R}
  \]
  \item \textbf{Campo en el eje de una espira circular (radio $R$, distancia $z$):}
  \[
    B(z) = \frac{\mu_0 I R^2}{2(R^2 + z^2)^{3/2}}
  \]
\end{itemize}

\subsubsection*{14.5 Ley de Ampère y Flujo Magnético}
\addcontentsline{toc}{subsubsection}{14.5 Ley de Ampère y Flujo Magnético}

\begin{itemize}
  \item \textbf{Ley de Ampère en forma integral:}
  \[
    \oint_C \vec{B} \cdot d\vec{\ell} = \mu_0 I_{\text{enc}}
  \]
  \item \textbf{Campo en el interior de un solenoide ideal ($n = N/L$):}
  \[
    B = \mu_0 n I
  \]
  \item \textbf{Campo en el interior de un toroide (radio medio $r$):}
  \[
    B = \frac{\mu_0 N I}{2\pi r}
  \]
  \item \textbf{Flujo Magnético ($\Phi_B$):}
  \[
    \Phi_B = \iint_S \vec{B} \cdot d\vec{A}
  \]
  \item \textbf{Ley de Gauss para el Magnetismo (No existencia de monopolos magnéticos aislados):}
  \[
    \oiint_{\partial V} \vec{B} \cdot d\vec{A} = 0
  \]
\end{itemize}

\subsubsection*{14.6 Efecto Hall}
\addcontentsline{toc}{subsubsection}{14.6 Efecto Hall}

\begin{itemize}
  \item \textbf{Voltaje Hall en un conductor rectangular (ancho $w$, espesor $d$):}
  \[
    V_H = \frac{I B}{n q d} = R_H \frac{I B}{d}, \quad R_H = \frac{1}{n q} \quad (\text{Coeficiente de Hall})
  \]
\end{itemize}

\subsubsection*{14.7 Medios Magnéticos y Magnetización}
\addcontentsline{toc}{subsubsection}{14.7 Medios Magnéticos y Magnetización}

\begin{itemize}
  \item \textbf{Vector Intensidad Magnética ($\vec{H}$):}
  \[
    \vec{H} = \frac{1}{\mu_0}\vec{B} - \vec{M} \implies \vec{B} = \mu_0 (\vec{H} + \vec{M}) = \mu \vec{H} = \mu_0 \mu_r \vec{H}
  \]
  \item \textbf{Vector Magnetización ($\vec{M}$) y Susceptibilidad Magnética ($\chi_m$):}
  \[
    \vec{M} = \chi_m \vec{H}, \quad \mu_r = 1 + \chi_m
  \]
  \item \textbf{Corrientes de Magnetización / Amperianas:}
  \[
    \vec{J}_b = \nabla \times \vec{M}, \quad \vec{K}_b = \vec{M} \times \hat{n}
  \]
\end{itemize}

\subsection*{Nivel Universitario / Avanzado}
\addcontentsline{toc}{subsection}{Nivel Universitario / Avanzado}

\subsubsection*{14.8 Ecuaciones Diferenciales de la Magnetostática}
\addcontentsline{toc}{subsubsection}{14.8 Ecuaciones Diferenciales de la Magnetostática}

\begin{itemize}
  \item \textbf{Divergencia nula de la inducción magnética:}
  \[
    \nabla \cdot \vec{B} = 0
  \]
  \item \textbf{Ley de Ampère Diferencial:}
  \[
    \nabla \times \vec{B} = \mu_0 \vec{J} \quad \implies \quad \nabla \times \vec{H} = \vec{J}_f
  \]
\end{itemize}

\subsubsection*{14.9 Potencial Vector Magnético ($\vec{A}$)}
\addcontentsline{toc}{subsubsection}{14.9 Potencial Vector Magnético ($\vec{A}$)}

\begin{itemize}
  \item \textbf{Definición ($\vec{B}$ es solenoidal):}
  \[
    \vec{B} = \nabla \times \vec{A}
  \]
  \item \textbf{Calibre / Gauge de Coulomb ($\nabla \cdot \vec{A} = 0$):}
  \[
    \nabla^2 \vec{A} = -\mu_0 \vec{J}
  \]
  \item \textbf{Solución integral del Potencial Vector:}
  \[
    \vec{A}(\vec{r}) = \frac{\mu_0}{4\pi} \int_V \frac{\vec{J}(\vec{r}')}{|\vec{r} - \vec{r}'|} dV'
  \]
\end{itemize}

\subsubsection*{14.10 Condiciones de Frontera Magnetostáticas}
\addcontentsline{toc}{subsubsection}{14.10 Condiciones de Frontera Magnetostáticas}

\begin{itemize}
  \item \textbf{Componente normal de la inducción magnética:}
  \[
    \hat{n} \cdot (\vec{B}_1 - \vec{B}_2) = 0 \implies B_{1n} = B_{2n}
  \]
  \item \textbf{Componente tangencial de la intensidad magnética:}
  \[
    \hat{n} \times (\vec{H}_1 - \vec{H}_2) = \vec{K}_f \implies H_{1t} - H_{2t} = K_f
  \]
\end{itemize}

\subsubsection*{14.11 Densidad y Energía Magnetostática Total}
\addcontentsline{toc}{subsubsection}{14.11 Densidad y Energía Magnetostática Total}

\begin{itemize}
  \item \textbf{Densidad volumétrica de energía magnética:}
  \[
    u_m = \frac{1}{2} \vec{B} \cdot \vec{H} = \frac{1}{2\mu_0} B^2
  \]
  \item \textbf{Energía magnética total:}
  \[
    U = \frac{1}{2} \int_V \vec{A} \cdot \vec{J} \, dV = \frac{1}{2} \int_{\text{todo el espacio}} \vec{B} \cdot \vec{H} \, dV
  \]
\end{itemize}

\section*{15. Inducción Electromagnética, Maxwell y Ondas}
\addcontentsline{toc}{section}{15. Inducción Electromagnética, Maxwell y Ondas}

\subsection*{Nivel Básico}
\addcontentsline{toc}{subsection}{Nivel Básico}

\subsubsection*{15.1 Ley de Faraday-Henry y Ley de Lenz}
\addcontentsline{toc}{subsubsection}{15.1 Ley de Faraday-Henry y Ley de Lenz}

\begin{itemize}
  \item \textbf{Fuerza Electromotriz Inducida ($\mathcal{E}$):}
  \[
    \mathcal{E} = -\frac{d\Phi_B}{dt} = -\frac{d}{dt}\left( \iint_S \vec{B} \cdot d\vec{A} \right)
  \]
  \item \textbf{FEM de movimiento en barra conductora (longitud $L$, velocidad $v \perp B$):}
  \[
    \mathcal{E} = B L v
  \]
\end{itemize}

\subsubsection*{15.2 Inductancia y Autoinducción}
\addcontentsline{toc}{subsubsection}{15.2 Inductancia y Autoinducción}

\begin{itemize}
  \item \textbf{Autoinductancia ($L$):}
  \[
    L = \frac{N \Phi_B}{I} \implies \mathcal{E}_L = -L \frac{dI}{dt}
  \]
  \item \textbf{Inductancia de un solenoide ideal (longitud $l$, área $A$, $N$ vueltas):}
  \[
    L = \mu_0 \frac{N^2 A}{l} = \mu_0 n^2 A l
  \]
  \item \textbf{Inductancia Mutua ($M$):}
  \[
    M_{12} = \frac{N_2 \Phi_{21}}{I_1}, \quad \mathcal{E}_2 = -M_{12} \frac{dI_1}{dt}
  \]
  \item \textbf{Energía acumulada en un inductor:}
  \[
    U = \frac{1}{2} L I^2
  \]
\end{itemize}

\subsection*{Nivel Intermedio}
\addcontentsline{toc}{subsection}{Nivel Intermedio}

\subsubsection*{15.3 Ley de Ampère-Maxwell y Corriente de Desplazamiento}
\addcontentsline{toc}{subsubsection}{15.3 Ley de Ampère-Maxwell y Corriente de Desplazamiento}

\begin{itemize}
  \item \textbf{Densidad de corriente de desplazamiento (Forma general macroscópica / Vacío):}
  \[
    \vec{J}_d = \frac{\partial \vec{D}}{\partial t} \quad (\text{General}), \quad \vec{J}_d = \varepsilon \frac{\partial \vec{E}}{\partial t} \quad (\text{Medio lineal}), \quad \vec{J}_d = \varepsilon_0 \frac{\partial \vec{E}}{\partial t} \quad (\text{Vacío})
  \]
  \item \textbf{Corriente de desplazamiento total:}
  \[
    I_d = \iint_S \vec{J}_d \cdot d\vec{A} = \frac{d\Phi_D}{dt} = \varepsilon_0 \frac{d\Phi_E}{dt} \quad (\text{en el vacío})
  \]
\end{itemize}

\subsubsection*{15.4 Ecuaciones de Maxwell en Forma Integral (Medios Lineales Generales)}
\addcontentsline{toc}{subsubsection}{15.4 Ecuaciones de Maxwell en Forma Integral (Medios Lineales Generales)}

\begin{itemize}
  \item \textbf{Ley de Gauss para el campo eléctrico:}
  \[
    \oiint_{\partial V} \vec{D} \cdot d\vec{A} = Q_{f,\text{enc}} \quad \left( \oiint_{\partial V} \vec{E} \cdot d\vec{A} = \frac{Q_{\text{enc}}}{\varepsilon_0} \right)
  \]
  \item \textbf{Ley de Gauss para el campo magnético:}
  \[
    \oiint_{\partial V} \vec{B} \cdot d\vec{A} = 0
  \]
  \item \textbf{Ley de Faraday de la Inducción:}
  \[
    \oint_C \vec{E} \cdot d\vec{\ell} = -\frac{d}{dt} \iint_S \vec{B} \cdot d\vec{A}
  \]
  \item \textbf{Ley de Ampère-Maxwell:}
  \[
    \oint_C \vec{H} \cdot d\vec{\ell} = I_{f,\text{enc}} + \frac{d}{dt} \iint_S \vec{D} \cdot d\vec{A}
  \]
\end{itemize}

\subsubsection*{15.5 Teorema de Poynting y Transporte de Energía}
\addcontentsline{toc}{subsubsection}{15.5 Teorema de Poynting y Transporte de Energía}

\begin{itemize}
  \item \textbf{Vector de Poynting ($\vec{S}$):}
  \[
    \vec{S} = \vec{E} \times \vec{H} \quad (\text{Medio material}), \quad \vec{S} = \frac{1}{\mu_0}(\vec{E} \times \vec{B}) \quad (\text{Vacío e isótropo})
  \]
  \item \textbf{Densidad volumétrica total de energía electromagnética:}
  \[
    u = \frac{1}{2} (\vec{D} \cdot \vec{E} + \vec{B} \cdot \vec{H}) \quad (\text{Medio lineal}), \quad u = \frac{1}{2} \left( \varepsilon_0 E^2 + \frac{1}{\mu_0} B^2 \right) \quad (\text{Vacío})
  \]
  \item \textbf{Teorema de Poynting Diferencial (Balance de Potencia):}
  \[
    \nabla \cdot \vec{S} + \frac{\partial u}{\partial t} = -\vec{J} \cdot \vec{E}
  \]
  \item \textbf{Intensidad de Onda Electromagnética ($\langle S \rangle$):}
  \[
    I = \langle |\vec{S}| \rangle = \frac{1}{2} \varepsilon_0 c E_0^2 = \frac{E_0^2}{2 \mu_0 c} = \frac{E_0 B_0}{2\mu_0}
  \]
\end{itemize}

\subsection*{Nivel Universitario / Avanzado}
\addcontentsline{toc}{subsection}{Nivel Universitario / Avanzado}

\subsubsection*{15.6 Ecuaciones de Maxwell en Forma Diferencial}
\addcontentsline{toc}{subsubsection}{15.6 Ecuaciones de Maxwell en Forma Diferencial}

\begin{itemize}
  \item \textbf{Forma Microscópica (en el vacío):}
  \[
    \nabla \cdot \vec{E} = \frac{\rho}{\varepsilon_0}
  \]
  \[
    \nabla \cdot \vec{B} = 0
  \]
  \[
    \nabla \times \vec{E} = -\frac{\partial \vec{B}}{\partial t}
  \]
  \[
    \nabla \times \vec{B} = \mu_0 \vec{J} + \mu_0 \varepsilon_0 \frac{\partial \vec{E}}{\partial t}
  \]
  \item \textbf{Forma Macroscópica (en medios materiales):}
  \[
    \nabla \cdot \vec{D} = \rho_f
  \]
  \[
    \nabla \cdot \vec{B} = 0
  \]
  \[
    \nabla \times \vec{E} = -\frac{\partial \vec{B}}{\partial t}
  \]
  \[
    \nabla \times \vec{H} = \vec{J}_f + \frac{\partial \vec{D}}{\partial t}
  \]
\end{itemize}

\subsubsection*{15.7 Potenciales Electrodinámicos y Transformaciones de Calibre}
\addcontentsline{toc}{subsubsection}{15.7 Potenciales Electrodinámicos y Transformaciones de Calibre}

\begin{itemize}
  \item \textbf{Definición de potenciales en términos dependientes del tiempo:}
  \[
    \vec{B} = \nabla \times \vec{A}, \quad \vec{E} = -\nabla V - \frac{\partial \vec{A}}{\partial t}
  \]
  \item \textbf{Transformaciones de Calibre (para cualquier función escalar $\Lambda(\vec{r}, t)$):}
  \[
    \vec{A}' = \vec{A} + \nabla \Lambda, \quad V' = V - \frac{\partial \Lambda}{\partial t}
  \]
  \item \textbf{Condición de Calibre de Lorenz:}
  \[
    \nabla \cdot \vec{A} + \frac{1}{c^2} \frac{\partial V}{\partial t} = 0 \implies \nabla \cdot \vec{A} + \mu_0 \varepsilon_0 \frac{\partial V}{\partial t} = 0
  \]
  \item \textbf{Ecuaciones de Onda Inhomogéneas para los Potenciales (en Calibre de Lorenz):}
  \[
    \nabla^2 V - \frac{1}{c^2} \frac{\partial^2 V}{\partial t^2} = -\frac{\rho}{\varepsilon_0} \iff \Box V = -\frac{\rho}{\varepsilon_0}
  \]
  \[
    \nabla^2 \vec{A} - \frac{1}{c^2} \frac{\partial^2 \vec{A}}{\partial t^2} = -\mu_0 \vec{J} \iff \Box \vec{A} = -\mu_0 \vec{J}
  \]
  \[
    (\Box = \nabla^2 - \frac{1}{c^2}\frac{\partial^2}{\partial t^2} = \text{Operador d'Alembertiano})
  \]
  \item \textbf{Potenciales Retardados de Liénard-Wiechert (Tiempo retardado $t_r = t - |\vec{r} - \vec{r}'|/c$):}
  \[
    V(\vec{r}, t) = \frac{1}{4\pi\varepsilon_0} \int \frac{\rho(\vec{r}', t_r)}{|\vec{r} - \vec{r}'|} dV'
  \]
  \[
    \vec{A}(\vec{r}, t) = \frac{\mu_0}{4\pi} \int \frac{\vec{J}(\vec{r}', t_r)}{|\vec{r} - \vec{r}'|} dV'
  \]
\end{itemize}

\subsubsection*{15.8 Ondas Electromagnéticas Planas en el Vacío}
\addcontentsline{toc}{subsubsection}{15.8 Ondas Electromagnéticas Planas en el Vacío}

\begin{itemize}
  \item \textbf{Ecuaciones de Onda Homogéneas:}
  \[
    \nabla^2 \vec{E} - \frac{1}{c^2} \frac{\partial^2 \vec{E}}{\partial t^2} = \vec{0}, \quad \nabla^2 \vec{B} - \frac{1}{c^2} \frac{\partial^2 \vec{B}}{\partial t^2} = \vec{0}
  \]
  \item \textbf{Velocidad de la luz en el vacío:}
  \[
    c = \frac{1}{\sqrt{\mu_0 \varepsilon_0}} \approx 2.9979 \times 10^8 \, \frac{\text{m}}{\text{s}}
  \]
  \item \textbf{Solución de Onda Plana Monocromática ($\vec{k} = \text{vector de onda}$):}
  \[
    \vec{E}(\vec{r}, t) = \vec{E}_0 e^{i(\vec{k} \cdot \vec{r} - \omega t)}, \quad \vec{B}(\vec{r}, t) = \vec{B}_0 e^{i(\vec{k} \cdot \vec{r} - \omega t)}
  \]
  \[
    \vec{B}_0 = \frac{1}{\omega} (\vec{k} \times \vec{E}_0) = \frac{1}{c} (\hat{k} \times \vec{E}_0), \quad \vec{k} \cdot \vec{E}_0 = 0, \quad \vec{k} \cdot \vec{B}_0 = 0
  \]
  \item \textbf{Impedancia Característica del Vacío ($\eta_0$):}
  \[
    \eta_0 = \sqrt{\frac{\mu_0}{\varepsilon_0}} \approx 120\pi \, \Omega \approx 376.73 \, \Omega
  \]
\end{itemize}

\subsubsection*{15.9 Tensor de Esfuerzos de Maxwell ($T_{ij}$) y Momento Electromagnético}
\addcontentsline{toc}{subsubsection}{15.9 Tensor de Esfuerzos de Maxwell ($T_{ij}$) y Momento Electromagnético}

\begin{itemize}
  \item \textbf{Tensor de Esfuerzos de Maxwell:}
  \[
    T_{ij} = \varepsilon_0 \left( E_i E_j - \frac{1}{2} \delta_{ij} E^2 \right) + \frac{1}{\mu_0} \left( B_i B_j - \frac{1}{2} \delta_{ij} B^2 \right)
  \]
  \item \textbf{Ley de Conservación del Momento Lineal:}
  \[
    \vec{f}_{\text{Lorentz}} + \frac{\partial \vec{g}}{\partial t} = \nabla \cdot \mathbf{T}, \quad \vec{g} = \varepsilon_0 (\vec{E} \times \vec{B}) = \frac{\vec{S}}{c^2} \quad (\text{Densidad de Momento})
  \]
  \item \textbf{Presión de Radiación ($P_{\text{rad}}$ para incidencia normal):}
  \[
    P_{\text{rad}} = \frac{I}{c} \quad (\text{Absorción total}), \quad P_{\text{rad}} = \frac{2I}{c} \quad (\text{Reflexión total})
  \]
\end{itemize}

\subsubsection*{15.10 Formulación Covariante / Cuadrivectorial de la Electrodinámica (Métrica $\eta_{\mu\nu} = \operatorname{diag}(+1, -1, -1, -1)$)}
\addcontentsline{toc}{subsubsection}{15.10 Formulación Covariante / Cuadrivectorial de la Electrodinámica (Métrica $\eta_{\mu\nu} = \operatorname{diag}(+1, -1, -1, -1)$)}

\begin{itemize}
  \item \textbf{Cuadrivector Posición, Gradiente y Corriente:}
  \[
    x^\mu = (ct, \vec{r}), \quad \partial_\mu = \left( \frac{1}{c}\frac{\partial}{\partial t}, \nabla \right), \quad \partial^\mu = \left( \frac{1}{c}\frac{\partial}{\partial t}, -\nabla \right), \quad J^\mu = (c\rho, \vec{J})
  \]
  \item \textbf{Cuadripotencial:}
  \[
    A^\mu = \left( \frac{V}{c}, \vec{A} \right), \quad A_\mu = \left( \frac{V}{c}, -\vec{A} \right)
  \]
  \item \textbf{Tensor de Campo Electromagnético de Faraday ($F^{\mu\nu} = \partial^\mu A^\nu - \partial^\nu A^\mu$):}
  \[
    F^{\mu\nu} = \begin{bmatrix} 0 & -E_x/c & -E_y/c & -E_z/c \\ E_x/c & 0 & -B_z & B_y \\ E_y/c & B_z & 0 & -B_x \\ E_z/c & -B_y & B_x & 0 \end{bmatrix}
  \]
  \item \textbf{Tensor Dual de Campo Electromagnético ($\tilde{F}^{\mu\nu} = \frac{1}{2} \epsilon^{\mu\nu\alpha\beta} F_{\alpha\beta}$):}
  \[
    \tilde{F}^{\mu\nu} = \begin{bmatrix} 0 & -B_x & -B_y & -B_z \\ B_x & 0 & E_z/c & -E_y/c \\ B_y & -E_z/c & 0 & E_x/c \\ B_z & E_y/c & -E_x/c & 0 \end{bmatrix}
  \]
  \item \textbf{Ecuaciones de Maxwell Covariantes:}
  \[
    \partial_\mu F^{\mu\nu} = \mu_0 J^\nu \quad (\text{Gauss y Ampère-Maxwell})
  \]
  \[
    \partial_\mu \tilde{F}^{\mu\nu} = 0 \quad (\text{Gauss Magnética y Faraday})
  \]
  \item \textbf{Invariantes de Lorentz del Campo Electromagnético:}
  \[
    I_1 = F^{\mu\nu} F_{\mu\nu} = 2 \left( B^2 - \frac{E^2}{c^2} \right) = \text{invariante}
  \]
  \[
    I_2 = \epsilon_{\mu\nu\alpha\beta} F^{\mu\nu} F^{\alpha\beta} = -\frac{4}{c} (\vec{E} \cdot \vec{B}) = \text{invariante}
  \]
\end{itemize}

\section*{Compendio V: Óptica y Luz}
\addcontentsline{toc}{section}{Compendio V: Óptica y Luz}

\textit{(Nivel Básico a Universitario)}

\section*{16. Óptica Geométrica y Sistemas Ópticos}
\addcontentsline{toc}{section}{16. Óptica Geométrica y Sistemas Ópticos}

\subsection*{Nivel Básico}
\addcontentsline{toc}{subsection}{Nivel Básico}

\subsubsection*{16.1 Propagación de la Luz y Leyes Fundamentales}
\addcontentsline{toc}{subsubsection}{16.1 Propagación de la Luz y Leyes Fundamentales}

\begin{itemize}
  \item \textbf{Índice de refracción absoluto ($c = \text{velocidad en el vacío}$):}
  \[
    n = \frac{c}{v}
  \]
  \item \textbf{Ley de Reflexión:}
  \[
    \theta_i = \theta_r
  \]
  \item \textbf{Ley de Refracción (Ley de Snell):}
  \[
    n_1 \sen(\theta_1) = n_2 \sen(\theta_2)
  \]
  \item \textbf{Ángulo Crítico y Reflexión Interna Total ($n_1 > n_2$):}
  \[
    \sen(\theta_c) = \frac{n_2}{n_1} \implies \theta_c = \arcsen\left( \frac{n_2}{n_1} \right)
  \]
\end{itemize}

\subsubsection*{16.2 Espejos Planos y Esféricos (Aproximación Paraxial)}
\addcontentsline{toc}{subsubsection}{16.2 Espejos Planos y Esféricos (Aproximación Paraxial)}

\begin{itemize}
  \item \textbf{Relación entre Radio de Curvatura y Distancia Focal:}
  \[
    f = \frac{R}{2}
  \]
  \item \textbf{Ecuación de Descartes para Espejos Esféricos ($s_o = \text{objeto}, \, s_i = \text{imagen}$):}
  \[
    \frac{1}{s_o} + \frac{1}{s_i} = \frac{1}{f} = \frac{2}{R}
  \]
  \item \textbf{Aumento Lateral Transversal ($m$):}
  \[
    m = \frac{y_i}{y_o} = -\frac{s_i}{s_o}
  \]
\end{itemize}

\subsection*{Nivel Intermedio}
\addcontentsline{toc}{subsection}{Nivel Intermedio}

\subsubsection*{16.3 Dioptrios Esféricos y Lentes Delgadas}
\addcontentsline{toc}{subsubsection}{16.3 Dioptrios Esféricos y Lentes Delgadas}

\begin{itemize}
  \item \textbf{Refracción en un Dioptrio Esférico:}
  \[
    \frac{n_1}{s_o} + \frac{n_2}{s_i} = \frac{n_2 - n_1}{R}
  \]
  \item \textbf{Fórmula del Fabricante de Lentes (Lente delgada en el aire $n_{\text{medio}} = 1$):}
  \[
    \frac{1}{f} = (n - 1)\left( \frac{1}{R_1} - \frac{1}{R_2} + \frac{(n - 1)d}{n R_1 R_2} \right) \xrightarrow{d \to 0} (n - 1)\left( \frac{1}{R_1} - \frac{1}{R_2} \right)
  \]
  \item \textbf{Ecuación de Gauss para Lentes Delgadas:}
  \[
    \frac{1}{s_o} + \frac{1}{s_i} = \frac{1}{f}
  \]
  \item \textbf{Forma Newtoniana de la Ecuación de Lentes ($x_o = s_o - f, \, x_i = s_i - f$):}
  \[
    x_o \cdot x_i = f^2
  \]
  \item \textbf{Potencia Óptica ($P$ en dioptrías $\text{m}^{-1}$):}
  \[
    P = \frac{1}{f}
  \]
  \item \textbf{Lentes Delgadas en Contacto:}
  \[
    P_{\text{eq}} = \sum_{i=1}^N P_i \implies \frac{1}{f_{\text{eq}}} = \frac{1}{f_1} + \frac{1}{f_2} + \dots + \frac{1}{f_N}
  \]
  \item \textbf{Lentes Separadas por una Distancia $d$:}
  \[
    \frac{1}{f_{\text{eq}}} = \frac{1}{f_1} + \frac{1}{f_2} - \frac{d}{f_1 f_2}
  \]
\end{itemize}

\subsubsection*{16.4 Prismas y Dispersión Cromática}
\addcontentsline{toc}{subsubsection}{16.4 Prismas y Dispersión Cromática}

\begin{itemize}
  \item \textbf{Desviación Angular en un Prisma de Ángulo de Ápice $A$:}
  \[
    \delta = \theta_1 + \theta_2' - A
  \]
  \item \textbf{Ángulo de Mínima Desviación ($\delta_{\min}$):}
  \[
    n = \frac{\sen\left( \frac{A + \delta_{\min}}{2} \right)}{\sen\left( \frac{A}{2} \right)}
  \]
  \item \textbf{Número de Abbe / Potencia Dispersiva ($V_d$):}
  \[
    V_d = \frac{n_d - 1}{n_F - n_C}
  \]
\end{itemize}

\subsection*{Nivel Universitario / Avanzado}
\addcontentsline{toc}{subsection}{Nivel Universitario / Avanzado}

\subsubsection*{16.5 Óptica Matricial (Matrices de Rayos ABCD)}
\addcontentsline{toc}{subsubsection}{16.5 Óptica Matricial (Matrices de Rayos ABCD)}

\begin{itemize}
  \item \textbf{Vector de Rayo Paraxial (Altura $y$, ángulo $\theta$ respecto al eje óptico):}
  \[
    \begin{bmatrix} y_2 \\ \theta_2 \end{bmatrix} = \begin{bmatrix} A & B \\ C & D \end{bmatrix} \begin{bmatrix} y_1 \\ \theta_1 \end{bmatrix}, \quad \det\left( \mathbf{M} \right) = AD - BC = \frac{n_1}{n_2}
  \]
  \item \textbf{Matriz de Traslación Libre en Medio Homogéneo (distancia $d$):}
  \[
    \mathbf{M}_{\text{tras}} = \begin{bmatrix} 1 & d \\ 0 & 1 \end{bmatrix}
  \]
  \item \textbf{Matriz de Refracción en Interfaz Plana:}
  \[
    \mathbf{M}_{\text{ref,plana}} = \begin{bmatrix} 1 & 0 \\ 0 & \frac{n_1}{n_2} \end{bmatrix}
  \]
  \item \textbf{Matriz de Refracción en Dioptrio Esférico:}
  \[
    \mathbf{M}_{\text{dioptrio}} = \begin{bmatrix} 1 & 0 \\ -\frac{n_2 - n_1}{n_2 R} & \frac{n_1}{n_2} \end{bmatrix}
  \]
  \item \textbf{Matriz de Lente Delgada:}
  \[
    \mathbf{M}_{\text{lente}} = \begin{bmatrix} 1 & 0 \\ -\frac{1}{f} & 1 \end{bmatrix}
  \]
  \item \textbf{Matriz de Espejo Esférico:}
  \[
    \mathbf{M}_{\text{espejo}} = \begin{bmatrix} 1 & 0 \\ -\frac{2}{R} & 1 \end{bmatrix}
  \]
\end{itemize}

\subsubsection*{16.6 Principio Variacional y Óptica Hamiltoniana}
\addcontentsline{toc}{subsubsection}{16.6 Principio Variacional y Óptica Hamiltoniana}

\begin{itemize}
  \item \textbf{Principio de Fermat (Camino Óptico Estacionario):}
  \[
    \delta \mathcal{S} = \delta \int_{P_1}^{P_2} n(\vec{r}) \, ds = 0
  \]
  \item \textbf{Longitud de Camino Óptico (OPL):}
  \[
    \text{OPL} = \int_C n(\vec{r}) \, ds
  \]
  \item \textbf{Ecuación Diferencial del Rayo:}
  \[
    \frac{d}{ds}\left( n \frac{d\vec{r}}{ds} \right) = \nabla n
  \]
  \item \textbf{Ecuación de la Eikonal (Límite de longitud de onda cero $\lambda \to 0$ para fase $S$):}
  \[
    |\nabla S|^2 = n^2(\vec{r})
  \]
\end{itemize}

\section*{17. Óptica Ondulatoria (Interferencia, Difracción y Polarización)}
\addcontentsline{toc}{section}{17. Óptica Ondulatoria (Interferencia, Difracción y Polarización)}

\subsection*{Nivel Básico}
\addcontentsline{toc}{subsection}{Nivel Básico}

\subsubsection*{17.1 Interferencia por División de Frente de Onda}
\addcontentsline{toc}{subsubsection}{17.1 Interferencia por División de Frente de Onda}

\begin{itemize}
  \item \textbf{Experimento de la Doble Rendija de Young (Separación $d$, distancia a pantalla $D \gg d$):}
  \item \textbf{Máximos de Interferencia (Constructiva):}
  \[
    d \sen\theta = m \lambda \implies y_{\text{máx}} \approx m \frac{\lambda D}{d}, \quad m \in \mathbb{Z}
  \]
  \item \textbf{Mínimos de Interferencia (Destructiva):}
  \[
    d \sen\theta = \left( m + \frac{1}{2} \right)\lambda \implies y_{\text{mín}} \approx \left( m + \frac{1}{2} \right)\frac{\lambda D}{d}, \quad m \in \mathbb{Z}
  \]
  \item \textbf{Separación entre Franjas Consecutivas ($\Delta y$):}
  \[
    \Delta y = \frac{\lambda D}{d}
  \]
\end{itemize}

\subsubsection*{17.2 Polarización de la Luz}
\addcontentsline{toc}{subsubsection}{17.2 Polarización de la Luz}

\begin{itemize}
  \item \textbf{Ley de Malus (Intensidad tras polarizador lineal con ángulo $\theta$):}
  \[
    I = I_0 \cos^2\theta
  \]
  \item \textbf{Ángulo de Polarización de Brewster:}
  \[
    \tan(\theta_B) = \frac{n_2}{n_1} \implies \theta_B + \theta_t = 90^\circ
  \]
\end{itemize}

\subsection*{Nivel Intermedio}
\addcontentsline{toc}{subsection}{Nivel Intermedio}

\subsubsection*{17.3 Distribución de Intensidad en Interferencia}
\addcontentsline{toc}{subsubsection}{17.3 Distribución de Intensidad en Interferencia}

\begin{itemize}
  \item \textbf{Superposición de Dos Ondas Armónicas Coherentes:}
  \[
    I = I_1 + I_2 + 2\sqrt{I_1 I_2} \cos(\delta)
  \]
  \item \textbf{Diferencia de Fase ($\delta$) por Diferencia de Camino Óptico ($\Delta r$):}
  \[
    \delta = k \Delta r + \Delta\phi_0 = \frac{2\pi}{\lambda} \Delta r + \Delta\phi_0
  \]
  \item \textbf{Visibilidad o Contraste de Franjas:}
  \[
    \mathcal{V} = \frac{I_{\text{máx}} - I_{\text{mín}}}{I_{\text{máx}} + I_{\text{mín}}} = \frac{2\sqrt{I_1 I_2}}{I_1 + I_2}
  \]
\end{itemize}

\subsubsection*{17.4 Interferencia en Películas Delgadas (Espesor $t$, incidencia casi normal)}
\addcontentsline{toc}{subsubsection}{17.4 Interferencia en Películas Delgadas (Espesor $t$, incidencia casi normal)}

\begin{itemize}
  \item \textbf{Condición de Interferencia Constructiva (Reflexión):}
  \[
    2 n t = \left( m + \frac{1}{2} \right)\lambda \quad (\text{con un solo cambio de fase de } \pi)
  \]
  \[
    2 n t = m \lambda \quad (\text{con } 0 \text{ o } 2 \text{ cambios de fase de } \pi)
  \]
  \item \textbf{Anillos de Newton (Radio del $m$-ésimo anillo brillante/oscuro por reflexión):}
  \[
    r_{\text{oscuro}} = \sqrt{m \lambda R}, \quad r_{\text{brillante}} = \sqrt{\left(m + \frac{1}{2}\right)\lambda R}
  \]
\end{itemize}

\subsubsection*{17.5 Difracción de Fraunhofer (Campo Lejano)}
\addcontentsline{toc}{subsubsection}{17.5 Difracción de Fraunhofer (Campo Lejano)}

\begin{itemize}
  \item \textbf{Rendija Simple de Ancho $a$:}
  \item \textbf{Mínimos de difracción:}
  \[
    a \sen\theta = m \lambda, \quad m = \pm 1, \pm 2, \dots
  \]
  \item \textbf{Distribución de intensidad:}
  \[
    I(\theta) = I_0 \left[ \frac{\sen(\beta)}{\beta} \right]^2 = I_0 \operatorname{sinc}^2(\beta), \quad \beta = \frac{\pi a}{\lambda}\sen\theta
  \]
  \item \textbf{Doble Rendija con Difracción e Interferencia Combinadas:}
  \[
    I(\theta) = I_0 \cos^2(\alpha) \left[ \frac{\sen(\beta)}{\beta} \right]^2, \quad \alpha = \frac{\pi d}{\lambda}\sen\theta, \quad \beta = \frac{\pi a}{\lambda}\sen\theta
  \]
  \item \textbf{Abertura Circular de Diámetro $D$ (Patrón de Airy):}
  \item \textbf{Primer anillo oscuro (Criterio de Rayleigh para resolución angular):}
  \[
    \sen\theta_{\min} \approx 1.22 \frac{\lambda}{D}
  \]
  \item \textbf{Redes de Difracción ($N$ rendijas, espaciado $d$):}
  \item \textbf{Ecuación de la red (Máximos principales):}
  \[
    d \sen\theta = m \lambda, \quad m \in \mathbb{Z}
  \]
  \item \textbf{Poder de Resolución Espectral ($R$):}
  \[
    \mathcal{R} = \frac{\lambda}{\Delta\lambda} = m N
  \]
  \item \textbf{Dispersión Angular ($D_\theta$):}
  \[
    D_\theta = \frac{d\theta}{d\lambda} = \frac{m}{d \cos\theta}
  \]
\end{itemize}

\subsubsection*{17.6 Ecuaciones de Fresnel para Interfaces Dieléctricas Planas}
\addcontentsline{toc}{subsubsection}{17.6 Ecuaciones de Fresnel para Interfaces Dieléctricas Planas}

\begin{itemize}
  \item \textbf{Coeficientes de Reflexión en Amplitud:}
  \[
    r_\perp = \frac{n_1 \cos\theta_i - n_2 \cos\theta_t}{n_1 \cos\theta_i + n_2 \cos\theta_t} = -\frac{\sen(\theta_i - \theta_t)}{\sen(\theta_i + \theta_t)}
  \]
  \[
    r_\parallel = \frac{n_2 \cos\theta_i - n_1 \cos\theta_t}{n_2 \cos\theta_i + n_1 \cos\theta_t} = \frac{\tan(\theta_i - \theta_t)}{\tan(\theta_i + \theta_t)}
  \]
  \item \textbf{Coeficientes de Transmisión en Amplitud:}
  \[
    t_\perp = \frac{2 n_1 \cos\theta_i}{n_1 \cos\theta_i + n_2 \cos\theta_t} = \frac{2 \sen\theta_t \cos\theta_i}{\sen(\theta_i + \theta_t)}
  \]
  \[
    t_\parallel = \frac{2 n_1 \cos\theta_i}{n_2 \cos\theta_i + n_1 \cos\theta_t} = \frac{2 \sen\theta_t \cos\theta_i}{\sen(\theta_i + \theta_t)\cos(\theta_i - \theta_t)}
  \]
  \item \textbf{Reflectancia ($R$) y Transmitancia ($T_w$) de Potencia:}
  \[
    R_\perp = |r_\perp|^2, \quad R_\parallel = |r_\parallel|^2, \quad R + T_w = 1
  \]
\end{itemize}

\subsection*{Nivel Universitario / Avanzado}
\addcontentsline{toc}{subsection}{Nivel Universitario / Avanzado}

\subsubsection*{17.7 Teoría Escalar de Difracción (Fresnel y Kirchhoff)}
\addcontentsline{toc}{subsubsection}{17.7 Teoría Escalar de Difracción (Fresnel y Kirchhoff)}

\begin{itemize}
  \item \textbf{Integral de Difracción de Fresnel-Kirchhoff:}
  \[
    U(P) = -\frac{i}{2\lambda} \iint_{\Sigma} U_0 \frac{e^{i k r}}{r} [\cos(\hat{n}, \vec{r}) - \cos(\hat{n}, \vec{r}_0)] dS
  \]
  \item \textbf{Integral de Difracción de Fresnel (Paraxial en distancia $z$):}
  \[
    U(x, y, z) = \frac{e^{i k z}}{i \lambda z} \iint_{-\infty}^{\infty} U(\xi, \eta, 0) \exp\left\{ \frac{i k}{2z}\left[ (x - \xi)^2 + (y - \eta)^2 \right] \right\} d\xi d\eta
  \]
  \item \textbf{Número de Fresnel ($N_F$):}
  \[
    N_F = \frac{a^2}{\lambda z} \quad (N_F \gg 1 \text{ Ópt. Geométrica}, \, N_F \sim 1 \text{ Fresnel}, \, N_F \ll 1 \text{ Fraunhofer})
  \]
  \item \textbf{Difracción de Fraunhofer como Transformada de Fourier:}
  \[
    U(x, y) \propto \mathcal{F}\{U(\xi, \eta)\}\Big|_{f_x = \frac{x}{\lambda z}, f_y = \frac{y}{\lambda z}}
  \]
\end{itemize}

\subsubsection*{17.8 Formalismo de Polarización (Jones y Stokes-Mueller)}
\addcontentsline{toc}{subsubsection}{17.8 Formalismo de Polarización (Jones y Stokes-Mueller)}

\begin{itemize}
  \item \textbf{Vector de Jones para Estados de Polarización Totalmente Coherentes:}
  \[
    \vec{J} = \begin{bmatrix} E_{0x} e^{i\phi_x} \\ E_{0y} e^{i\phi_y} \end{bmatrix}
  \]
  \item \textbf{Matrices de Jones de Elementos Ópticos Clásicos:}
  \[
    \mathbf{M}_{\text{pol,H}} = \begin{bmatrix} 1 & 0 \\ 0 & 0 \end{bmatrix}, \quad \mathbf{M}_{\text{retardador}}(\Gamma) = \begin{bmatrix} e^{i\Gamma/2} & 0 \\ 0 & e^{-i\Gamma/2} \end{bmatrix}
  \]
  \item \textbf{Parámetros de Stokes (Para luz arbitraria o parcialmente polarizada):}
  \[
    S_0 = I = I_H + I_V, \quad S_1 = Q = I_H - I_V, \quad S_2 = U = I_{+45^\circ} - I_{-45^\circ}, \quad S_3 = V = I_R - I_L
  \]
  \item \textbf{Grado de Polarización ($DOP$):}
  \[
    \text{DOP} = \frac{\sqrt{S_1^2 + S_2^2 + S_3^2}}{S_0} \le 1
  \]
  \item \textbf{Ecuación de Transformación de Mueller:}
  \[
    \vec{S}_{\text{sal}} = \mathbf{M}_{\text{Mueller}} \vec{S}_{\text{ent}}
  \]
\end{itemize}

\subsubsection*{17.9 Birrefringencia y Óptica de Cristales Anisótropos}
\addcontentsline{toc}{subsubsection}{17.9 Birrefringencia y Óptica de Cristales Anisótropos}

\begin{itemize}
  \item \textbf{Elipsoide de Índices Ópticos (Indicatriz Óptica):}
  \[
    \frac{x^2}{n_x^2} + \frac{y^2}{n_y^2} + \frac{z^2}{n_z^2} = 1
  \]
  \item \textbf{Retardo de Fase en Lámina Birrefringente de Espesor $d$:}
  \[
    \Delta\phi = \Gamma = \frac{2\pi}{\lambda} |n_e - n_o| d
  \]
  \[
    \text{Lámina de Cuarto de Onda: } d = \frac{\lambda}{4|n_e - n_o|}, \quad \text{Lámina de Media Onda: } d = \frac{\lambda}{2|n_e - n_o|}
  \]
\end{itemize}

\section*{18. Óptica Electromagnética, Dispersión y Guías de Ondas}
\addcontentsline{toc}{section}{18. Óptica Electromagnética, Dispersión y Guías de Ondas}

\subsection*{Nivel Intermedio}
\addcontentsline{toc}{subsection}{Nivel Intermedio}

\subsubsection*{18.1 Dispersión y Propagación en Medios Materiales}
\addcontentsline{toc}{subsubsection}{18.1 Dispersión y Propagación en Medios Materiales}

\begin{itemize}
  \item \textbf{Relación entre Índice de Refracción y Constantes Dieléctricas:}
  \[
    n = \sqrt{\varepsilon_r \mu_r} \approx \sqrt{\varepsilon_r}
  \]
  \item \textbf{Ecuación Empírica de Dispersión de Cauchy:}
  \[
    n(\lambda) = A + \frac{B}{\lambda^2} + \frac{C}{\lambda^4} + \dots
  \]
  \item \textbf{Ecuación de Sellmeier:}
  \[
    n^2(\lambda) = 1 + \sum_{i} \frac{B_i \lambda^2}{\lambda^2 - C_i}
  \]
  \item \textbf{Velocidad de Fase y Velocidad de Grupo:}
  \[
    v_p = \frac{\omega}{k} = \frac{c}{n}, \quad v_g = \frac{d\omega}{dk} = \frac{c}{n_g}, \quad n_g = n - \lambda \frac{dn}{d\lambda}
  \]
  \item \textbf{Dispersión por Retardo de Grupo (GVD / Parámetro $D$):}
  \[
    D = -\frac{\lambda}{c} \frac{d^2 n}{d\lambda^2} = -\frac{2\pi c}{\lambda^2} \beta_2, \quad \beta_2 = \frac{d^2 k}{d\omega^2}
  \]
\end{itemize}

\subsubsection*{18.2 Atenuación y Absorción}
\addcontentsline{toc}{subsubsection}{18.2 Atenuación y Absorción}

\begin{itemize}
  \item \textbf{Índice de Refracción Complejo:}
  \[
    \tilde{n} = n + i \kappa
  \]
  \item \textbf{Ley de Beer-Lambert (Coeficiente de absorción $\alpha$):}
  \[
    I(z) = I_0 e^{-\alpha z}, \quad \alpha = \frac{2\omega\kappa}{c} = \frac{4\pi\kappa}{\lambda_0}
  \]
  \item \textbf{Profundidad de Penetración / Efecto Piel Óptico ($\delta_s$):}
  \[
    \delta_s = \frac{1}{\alpha} = \frac{\lambda_0}{4\pi\kappa}
  \]
\end{itemize}

\subsection*{Nivel Universitario / Avanzado}
\addcontentsline{toc}{subsection}{Nivel Universitario / Avanzado}

\subsubsection*{18.3 Modelo Microeléctrico de Lorentz-Drude y Relaciones Dispersivas}
\addcontentsline{toc}{subsubsection}{18.3 Modelo Microeléctrico de Lorentz-Drude y Relaciones Dispersivas}

\begin{itemize}
  \item \textbf{Permitividad Relativa Compleja del Modelo de Oscilador de Lorentz:}
  \[
    \tilde{\varepsilon}_r(\omega) = 1 + \frac{N q^2}{\varepsilon_0 m_e} \sum_j \frac{f_j}{\omega_{0j}^2 - \omega^2 - i\gamma_j \omega}
  \]
  \item \textbf{Frecuencia de Plasma (Metales / Modelo de Drude $\omega_{0} = 0$):}
  \[
    \omega_p^2 = \frac{N e^2}{\varepsilon_0 m_e} \implies \varepsilon(\omega) = 1 - \frac{\omega_p^2}{\omega^2 + i\gamma\omega}
  \]
  \item \textbf{Relaciones de Kramers-Kronig (Causalidad en la respuesta dieléctrica):}
  \[
    \text{Re}[\chi(\omega)] = \frac{2}{\pi} \mathcal{P} \int_0^\infty \frac{\Omega \, \text{Im}[\chi(\Omega)]}{\Omega^2 - \omega^2} d\Omega
  \]
  \[
    \text{Im}[\chi(\omega)] = -\frac{2\omega}{\pi} \mathcal{P} \int_0^\infty \frac{\text{Re}[\chi(\Omega)]}{\Omega^2 - \omega^2} d\Omega
  \]
\end{itemize}

\subsubsection*{18.4 Fibras Ópticas y Guías de Ondas Dieléctricas}
\addcontentsline{toc}{subsubsection}{18.4 Fibras Ópticas y Guías de Ondas Dieléctricas}

\begin{itemize}
  \item \textbf{Apertura Numérica ($\text{NA}$):}
  \[
    \text{NA} = \sen(\theta_{\text{máx}}) = \sqrt{n_{\text{núcleo}}^2 - n_{\text{revestimiento}}^2} \approx n_1 \sqrt{2\Delta}, \quad \Delta = \frac{n_1 - n_2}{n_1}
  \]
  \item \textbf{Parámetro de Frecuencia Normalizada (Número $V$ para fibra de radio $a$):}
  \[
    V = \frac{2\pi a}{\lambda_0} \sqrt{n_1^2 - n_2^2} = \frac{2\pi a}{\lambda_0} \text{NA}
  \]
  \item \textbf{Condición de Monomodo en Fibra de Índice Escalonado:}
  \[
    V < 2.4048 \quad (\text{Primer cero de la función de Bessel } J_0)
  \]
  \item \textbf{Número Aproximado de Modos en Fibra Multimodo de Índice Escalonado ($V \gg 1$):}
  \[
    M \approx \frac{V^2}{2}
  \]
\end{itemize}

\subsubsection*{18.5 Ondas Evanescentes en Reflexión Total Interna}
\addcontentsline{toc}{subsubsection}{18.5 Ondas Evanescentes en Reflexión Total Interna}

\begin{itemize}
  \item \textbf{Vector de onda transversal transmitido imaginario:}
  \[
    k_{tz} = i \alpha = i k_0 \sqrt{n_1^2 \sin^2(\theta_i) - n_2^2}
  \]
  \item \textbf{Decaimiento de la Amplitud del Campo Evanescente:}
  \[
    \vec{E}_t(z) = \vec{E}_{t0} e^{-\alpha z} e^{i(k_{tx} x - \omega t)}
  \]
  \item \textbf{Desplazamiento Lateral de Goos-Hänchen (Polarización s / TE cerca del ángulo crítico):}
  \[
    \Delta x_{\text{GH}} = \frac{2}{k_0 \sqrt{n_1^2 \sin^2(\theta_i) - n_2^2}} \quad \left( \text{Formulación general: } \Delta x_{\text{GH}} = -\frac{d\phi}{dk_x} \right)
  \]
\end{itemize}

\section*{19. Óptica Cuántica, Radiación, Haces Láser y No Lineal}
\addcontentsline{toc}{section}{19. Óptica Cuántica, Radiación, Haces Láser y No Lineal}

\subsection*{Nivel Básico}
\addcontentsline{toc}{subsection}{Nivel Básico}

\subsubsection*{19.1 Fotones y Naturaleza Cuántica de la Luz}
\addcontentsline{toc}{subsubsection}{19.1 Fotones y Naturaleza Cuántica de la Luz}

\begin{itemize}
  \item \textbf{Energía del Fotón (Relación de Planck-Einstein):}
  \[
    E = h f = \hbar \omega, \quad \hbar = \frac{h}{2\pi}
  \]
  \item \textbf{Momento Lineal del Fotón (Relación de De Broglie):}
  \[
    p = \frac{h}{\lambda} = \hbar k = \frac{E}{c}
  \]
  \item \textbf{Ecuación del Efecto Fotoeléctrico de Einstein:}
  \[
    E_{k,\text{máx}} = e V_{\text{corte}} = h f - \Phi
  \]
  \item \textbf{Presión de Radiación ($P_{\text{rad}}$ con densidad de flujo radiante / irradiancia $I$):}
  \[
    P_{\text{rad}} = \frac{I}{c} \quad (\text{Absorción pura}), \quad P_{\text{rad}} = \frac{2I}{c} \quad (\text{Reflexión perfecta})
  \]
\end{itemize}

\subsection*{Nivel Intermedio}
\addcontentsline{toc}{subsection}{Nivel Intermedio}

\subsubsection*{19.2 Leyes de Radiación Térmica (Cuerpo Negro)}
\addcontentsline{toc}{subsubsection}{19.2 Leyes de Radiación Térmica (Cuerpo Negro)}

\begin{itemize}
  \item \textbf{Ley de Distribución Espectral de Planck:}
  \[
    u_\lambda(\lambda, T) = \frac{8\pi h c}{\lambda^5} \frac{1}{e^{\frac{h c}{\lambda k_B T}} - 1}, \quad I_\nu(\nu, T) = \frac{2h\nu^3}{c^2}\frac{1}{e^{\frac{h\nu}{k_B T}} - 1}
  \]
  \item \textbf{Ley de Desplazamiento de Wien:}
  \[
    \lambda_{\text{máx}} T = b \approx 2.8978 \times 10^{-3} \, \text{m}\cdot\text{K}
  \]
  \item \textbf{Ley de Stefan-Boltzmann:}
  \[
    j^* = \sigma T^4 = \varepsilon \sigma T^4, \quad \sigma = \frac{2\pi^5 k_B^4}{15 c^2 h^3} \approx 5.6704 \times 10^{-8} \, \frac{\text{W}}{\text{m}^2\cdot\text{K}^4}
  \]
\end{itemize}

\subsubsection*{19.3 Radiometría y Fotometría}
\addcontentsline{toc}{subsubsection}{19.3 Radiometría y Fotometría}

\begin{itemize}
  \item \textbf{Flujo Radiante ($\Phi_e$ en W) y Flujo Luminoso ($\Phi_v$ en lúmenes lm):}
  \[
    \Phi_v = K_m \int_0^\infty \Phi_{e,\lambda}(\lambda) V(\lambda) d\lambda, \quad K_m \approx 683 \, \frac{\text{lm}}{\text{W}}
  \]
  \item \textbf{Irradiancia ($E_e$) e Iluminancia ($E_v$ en lux):}
  \[
    E_e = \frac{d\Phi_e}{dA}, \quad E_v = \frac{d\Phi_v}{dA}
  \]
  \item \textbf{Intensidad Radiante ($I_e$) e Intensidad Luminosa ($I_v$ en candelas cd):}
  \[
    I = \frac{d\Phi}{d\Omega}
  \]
  \item \textbf{Ley de Lambert (Superficie Difusa Emisora / Reflectora):}
  \[
    I(\theta) = I_0 \cos\theta
  \]
\end{itemize}

\subsubsection*{19.4 Física del Láser y Coeficientes de Einstein}
\addcontentsline{toc}{subsubsection}{19.4 Física del Láser y Coeficientes de Einstein}

\begin{itemize}
  \item \textbf{Relación entre Emisión Espontánea ($A_{21}$), Estimulada ($B_{21}$) y Absorción ($B_{12}$):}
  \[
    g_1 B_{12} = g_2 B_{21}, \quad A_{21} = \frac{8\pi h \nu^3}{c^3} B_{21}
  \]
  \item \textbf{Condición de Inversión de Población para Ganancia Óptica:}
  \[
    N_2 - \frac{g_2}{g_1} N_1 > 0
  \]
  \item \textbf{Ganancia Óptica en Medio Activo (Sección eficaz estimulada $\sigma_{21}$):}
  \[
    I(z) = I_0 e^{\gamma(\nu) z}, \quad \gamma(\nu) = \sigma_{21}(\nu)\left[ N_2 - \frac{g_2}{g_1} N_1 \right]
  \]
  \item \textbf{Condición de Umbral Láser en Cavidad de Longitud $L$ (Reflectancias $R_1, R_2$):}
  \[
    \gamma_{\text{th}} = \alpha_{\text{pérdidas}} + \frac{1}{2L}\ln\left( \frac{1}{R_1 R_2} \right)
  \]
\end{itemize}

\subsection*{Nivel Universitario / Avanzado}
\addcontentsline{toc}{subsection}{Nivel Universitario / Avanzado}

\subsubsection*{19.5 Propagación de Haces Gaussianos ($\text{TEM}_{00}$)}
\addcontentsline{toc}{subsubsection}{19.5 Propagación de Haces Gaussianos ($\text{TEM}_{00}$)}

\begin{itemize}
  \item \textbf{Parámetro de Rayleigh / Distancia Confocal ($z_R$):}
  \[
    z_R = \frac{\pi w_0^2}{\lambda_0}
  \]
  \item \textbf{Cintura y Perfil del Radio del Haz ($w(z)$ con cintura mínima $w_0$ en $z=0$):}
  \[
    w(z) = w_0 \sqrt{1 + \left( \frac{z}{z_R} \right)^2}
  \]
  \item \textbf{Radio de Curvatura del Frente de Onda ($R(z)$):}
  \[
    R(z) = z \left[ 1 + \left( \frac{z_R}{z} \right)^2 \right]
  \]
  \item \textbf{Divergencia Angular de Campo Lejano ($\theta_{\text{div}}$):}
  \[
    \theta_{\text{div}} = \lim_{z \to \infty} \frac{w(z)}{z} = \frac{\lambda_0}{\pi w_0}
  \]
  \item \textbf{Fase de Gouy:}
  \[
    \zeta(z) = \arctan\left( \frac{z}{z_R} \right)
  \]
  \item \textbf{Parámetro Complejo del Haz $q(z)$ y Transformación $ABCD$:}
  \[
    \frac{1}{q(z)} = \frac{1}{R(z)} - i \frac{\lambda_0}{\pi n w^2(z)} \implies q_2 = \frac{A q_1 + B}{C q_1 + D}
  \]
\end{itemize}

\subsubsection*{19.6 Óptica No Lineal}
\addcontentsline{toc}{subsubsection}{19.6 Óptica No Lineal}

\begin{itemize}
  \item \textbf{Expansión de la Polarización No Lineal:}
  \[
    \vec{P} = \varepsilon_0 \left( \chi^{(1)} \cdot \vec{E} + \chi^{(2)} : \vec{E} \vec{E} + \chi^{(3)} \vdots \vec{E} \vec{E} \vec{E} + \dots \right) = \vec{P}_L + \vec{P}_{\text{NL}}
  \]
  \item \textbf{Generación de Segundo Armónico (SHG):}
  \[
    P^{(2)}(2\omega) = \varepsilon_0 \chi^{(2)} E^2(\omega)
  \]
  \item \textbf{Condición de Ajuste de Fase (Phase Matching):}
  \[
    \Delta k = k(2\omega) - 2k(\omega) = 0 \implies n(2\omega) = n(\omega)
  \]
  \item \textbf{Efecto Kerr Óptico (Índice de refracción dependiente de la intensidad):}
  \[
    n(I) = n_0 + n_2 I, \quad n_2 = \frac{3 \chi^{(3)}}{4 \varepsilon_0 c n_0^2}
  \]
  \item \textbf{Efecto Pockels (Lineal electro-óptico):}
  \[
    \Delta\left(\frac{1}{n^2}\right) = r_{ijk} E_k
  \]
  \item \textbf{Efecto Kerr Electro-óptico (Cuadrático electro-óptico):}
  \[
    \Delta n = \lambda K_{\text{Kerr}} E^2
  \]
\end{itemize}

\section*{Compendio VI: Física Moderna}
\addcontentsline{toc}{section}{Compendio VI: Física Moderna}

\textit{(Nivel Básico a Universitario)}

\section*{20. Teoría de la Relatividad (Especial y General)}
\addcontentsline{toc}{section}{20. Teoría de la Relatividad (Especial y General)}

\subsection*{Nivel Básico}
\addcontentsline{toc}{subsection}{Nivel Básico}

\subsubsection*{20.1 Cinemática Relativista Unidimensional}
\addcontentsline{toc}{subsubsection}{20.1 Cinemática Relativista Unidimensional}

\begin{itemize}
  \item \textbf{Factor de Lorentz ($\gamma$ con $\beta = v/c$):}
  \[
    \gamma = \frac{1}{\sqrt{1 - \frac{v^2}{c^2}}} = \frac{1}{\sqrt{1 - \beta^2}}
  \]
  \item \textbf{Dilatación Temporal ($\Delta t_0 = \text{tiempo propio}$):}
  \[
    \Delta t = \gamma \Delta t_0 = \frac{\Delta t_0}{\sqrt{1 - \frac{v^2}{c^2}}}
  \]
  \item \textbf{Contracción de la Longitud ($L_0 = \text{longitud propia}$ en dirección del movimiento):}
  \[
    L = \frac{L_0}{\gamma} = L_0 \sqrt{1 - \frac{v^2}{c^2}}
  \]
  \item \textbf{Adición Relativista de Velocidades (Movimiento en el eje $x$):}
  \[
    u_x' = \frac{u_x - v}{1 - \frac{u_x v}{c^2}} \iff u_x = \frac{u_x' + v}{1 + \frac{u_x' v}{c^2}}
  \]
\end{itemize}

\subsubsection*{20.2 Dinámica y Equivalencia Masa-Energía}
\addcontentsline{toc}{subsubsection}{20.2 Dinámica y Equivalencia Masa-Energía}

\begin{itemize}
  \item \textbf{Momento Lineal Relativista:}
  \[
    \vec{p} = \gamma m_0 \vec{v} = \frac{m_0 \vec{v}}{\sqrt{1 - \frac{v^2}{c^2}}}
  \]
  \item \textbf{Energía en Reposo:}
  \[
    E_0 = m_0 c^2
  \]
  \item \textbf{Energía Total de una Partícula Libre:}
  \[
    E = \gamma m_0 c^2 = E_k + m_0 c^2
  \]
  \item \textbf{Energía Cinética Relativista:}
  \[
    E_k = (\gamma - 1) m_0 c^2 = m_0 c^2 \left( \frac{1}{\sqrt{1 - \frac{v^2}{c^2}}} - 1 \right)
  \]
  \item \textbf{Relación Fundamental Energía-Momento:}
  \[
    E^2 = (p c)^2 + (m_0 c^2)^2 \implies E = \sqrt{p^2 c^2 + m_0^2 c^4}
  \]
  \[
    (\text{Para fotones y partículas sin masa } m_0 = 0 \implies E = p c)
  \]
\end{itemize}

\subsubsection*{20.3 Efecto Doppler Relativista}
\addcontentsline{toc}{subsubsection}{20.3 Efecto Doppler Relativista}

\begin{itemize}
  \item \textbf{Fuente y Observador en Dirección Longitudinal:}
  \[
    f_o = f_s \sqrt{\frac{1 - \beta}{1 + \beta}} \quad (\text{Alejamiento}) \quad | \quad f_o = f_s \sqrt{\frac{1 + \beta}{1 - \beta}} \quad (\text{Acercamiento})
  \]
  \item \textbf{Efecto Doppler Transversal ($\theta = 90^\circ$):}
  \[
    f_o = \frac{f_s}{\gamma} = f_s \sqrt{1 - \beta^2}
  \]
\end{itemize}

\subsection*{Nivel Intermedio}
\addcontentsline{toc}{subsection}{Nivel Intermedio}

\subsubsection*{20.4 Transformaciones de Lorentz y Espacio-Tiempo de Minkowski}
\addcontentsline{toc}{subsubsection}{20.4 Transformaciones de Lorentz y Espacio-Tiempo de Minkowski}

\begin{itemize}
  \item \textbf{Transformación de Coordenadas de Lorentz (Movimiento estándar a lo largo de $x$):}
  \[
    x' = \gamma (x - v t), \quad y' = y, \quad z' = z, \quad t' = \gamma \left( t - \frac{v x}{c^2} \right)
  \]
  \item \textbf{Transformación Inversa:}
  \[
    x = \gamma (x' + v t'), \quad y = y', \quad z = z', \quad t = \gamma \left( t' + \frac{v x'}{c^2} \right)
  \]
  \item \textbf{Intervalo Espaciotemporal Invariante ($s^2$):}
  \[
    \Delta s^2 = c^2 \Delta t^2 - (\Delta x^2 + \Delta y^2 + \Delta z^2) = c^2 \Delta t^2 - |\Delta\vec{r}|^2 = \text{invariante}
  \]
  \[
    \Delta s^2 > 0 \quad (\text{Tipo Tiempo}), \quad \Delta s^2 = 0 \quad (\text{Tipo Luz / Nulo}), \quad \Delta s^2 < 0 \quad (\text{Tipo Espacio})
  \]
  \item \textbf{Rapidez Relativista ($\theta$ / Parametrización Hiperbólica):}
  \[
    \beta = \tanh(\theta), \quad \gamma = \cosh(\theta), \quad \gamma\beta = \sinh(\theta) \implies \theta_{\text{total}} = \theta_1 + \theta_2
  \]
\end{itemize}

\subsubsection*{20.5 Formalismo Cuadrivectorial}
\addcontentsline{toc}{subsubsection}{20.5 Formalismo Cuadrivectorial}

\begin{itemize}
  \item \textbf{Métrica de Minkowski $\eta_{\mu\nu}$ (Convención $+,-,-,-$):}
  \[
    \eta_{\mu\nu} = \operatorname{diag}(1, -1, -1, -1)
  \]
  \item \textbf{Cuadrivector Posición:}
  \[
    X^\mu = (ct, x, y, z) = (ct, \vec{r})
  \]
  \item \textbf{Cuadrivector Velocidad (Tiempo propio $\tau$ con $d\tau = dt/\gamma$):}
  \[
    U^\mu = \frac{dX^\mu}{d\tau} = \gamma(c, \vec{v}), \quad U^\mu U_\mu = \eta_{\mu\nu} U^\mu U^\nu = c^2
  \]
  \item \textbf{Cuadrivector Momento / Cuadrimomento:}
  \[
    P^\mu = m_0 U^\mu = \left( \frac{E}{c}, \vec{p} \right), \quad P^\mu P_\mu = \frac{E^2}{c^2} - |\vec{p}|^2 = m_0^2 c^2
  \]
  \item \textbf{Cuadrivector Fuerza de Minkowski:}
  \[
    K^\mu = \frac{dP^\mu}{d\tau} = \gamma \left( \frac{\vec{F} \cdot \vec{v}}{c}, \vec{F} \right)
  \]
  \item \textbf{Cuadrivector Onda de De Broglie:}
  \[
    K^\mu = \left( \frac{\omega}{c}, \vec{k} \right), \quad P^\mu = \hbar K^\mu
  \]
\end{itemize}

\subsection*{Nivel Universitario / Avanzado}
\addcontentsline{toc}{subsection}{Nivel Universitario / Avanzado}

\subsubsection*{20.6 Elementos de Relatividad General}
\addcontentsline{toc}{subsubsection}{20.6 Elementos de Relatividad General}

\begin{itemize}
  \item \textbf{Principio de Equivalencia y Desplazamiento al Rojo Gravitacional:}
  \[
    \frac{\Delta f}{f_0} = \frac{\Delta\Phi_g}{c^2} \implies f(r) = f_\infty \sqrt{1 - \frac{2GM}{c^2 r}}
  \]
  \item \textbf{Métrica General del Espacio-Tiempo Curvo:}
  \[
    ds^2 = g_{\mu\nu} dx^\mu dx^\nu
  \]
  \item \textbf{Símbolos de Christoffel (Conexión Levi-Civita):}
  \[
    \Gamma^\mu_{\alpha\beta} = \frac{1}{2} g^{\mu\sigma} \left( \frac{\partial g_{\sigma\alpha}}{\partial x^\beta} + \frac{\partial g_{\sigma\beta}}{\partial x^\alpha} - \frac{\partial g_{\alpha\beta}}{\partial x^\sigma} \right)
  \]
  \item \textbf{Ecuación de las Geodésicas:}
  \[
    \frac{d^2 x^\mu}{d\lambda^2} + \Gamma^\mu_{\alpha\beta} \frac{dx^\alpha}{d\lambda} \frac{dx^\beta}{d\lambda} = 0
  \]
  \item \textbf{Tensor de Curvatura de Riemann:}
  \[
    R^\rho_{\ \sigma\mu\nu} = \partial_\mu \Gamma^\rho_{\nu\sigma} - \partial_\nu \Gamma^\rho_{\mu\sigma} + \Gamma^\rho_{\mu\lambda} \Gamma^\lambda_{\nu\sigma} - \Gamma^\rho_{\nu\lambda} \Gamma^\lambda_{\mu\sigma}
  \]
  \item \textbf{Tensor de Ricci y Escalar de Curvatura:}
  \[
    R_{\mu\nu} = R^\lambda_{\ \mu\lambda\nu}, \quad R = g^{\mu\nu} R_{\mu\nu}
  \]
  \item \textbf{Ecuaciones de Campo de Einstein:}
  \[
    G_{\mu\nu} + \Lambda g_{\mu\nu} = \frac{8\pi G}{c^4} T_{\mu\nu}, \quad G_{\mu\nu} = R_{\mu\nu} - \frac{1}{2} R g_{\mu\nu}
  \]
  \item \textbf{Métrica de Schwarzschild (Masa esférica estática en el vacío):}
  \[
    ds^2 = -\left( 1 - \frac{r_s}{r} \right) c^2 dt^2 + \left( 1 - \frac{r_s}{r} \right)^{-1} dr^2 + r^2 (d\theta^2 + \sen^2\theta \, d\phi^2)
  \]
  \item \textbf{Radio de Schwarzschild (Horizonte de sucesos):}
  \[
    r_s = \frac{2GM}{c^2}
  \]
\end{itemize}

\section*{21. Mecánica Cuántica y Física Atómica}
\addcontentsline{toc}{section}{21. Mecánica Cuántica y Física Atómica}

\subsection*{Nivel Básico}
\addcontentsline{toc}{subsection}{Nivel Básico}

\subsubsection*{21.1 Fenomenología Cuántica Temprana}
\addcontentsline{toc}{subsubsection}{21.1 Fenomenología Cuántica Temprana}

\begin{itemize}
  \item \textbf{Hipótesis de De Broglie (Dualidad onda-partícula):}
  \[
    \lambda = \frac{h}{p} = \frac{h}{\gamma m v}, \quad p = \hbar k = \frac{h}{\lambda}
  \]
  \item \textbf{Efecto Fotoeléctrico (Einstein):}
  \[
    E_k^{\text{máx}} = e V_s = h f - \Phi_0 = \hbar \omega - \Phi_0, \quad f_0 = \frac{\Phi_0}{h}
  \]
  \item \textbf{Dispersión Compton (Corrimiento en longitud de onda):}
  \[
    \Delta\lambda = \lambda' - \lambda = \frac{h}{m_e c}(1 - \cos\theta) = \lambda_C (1 - \cos\theta)
  \]
  \[
    \lambda_C = \frac{h}{m_e c} \approx 2.4263 \times 10^{-12} \, \text{m} \quad (\text{Longitud de onda Compton})
  \]
  \item \textbf{Principio de Incertidumbre de Heisenberg:}
  \[
    \Delta x \cdot \Delta p_x \ge \frac{\hbar}{2}, \quad \Delta E \cdot \Delta t \ge \frac{\hbar}{2}
  \]
\end{itemize}

\subsubsection*{21.2 Modelo Atómico de Bohr (Átomos Hidrogenoides de número atómico $Z$)}
\addcontentsline{toc}{subsubsection}{21.2 Modelo Atómico de Bohr (Átomos Hidrogenoides de número atómico $Z$)}

\begin{itemize}
  \item \textbf{Cuantización del Momento Angular Orbital:}
  \[
    L = m_e v r = n \hbar = n \frac{h}{2\pi}, \quad n = 1, 2, 3, \dots
  \]
  \item \textbf{Radios de Órbitas Permitidas ($a_0 = \text{Radio de Bohr}$):}
  \[
    r_n = \frac{4\pi\varepsilon_0 \hbar^2}{m_e e^2} \frac{n^2}{Z} = a_0 \frac{n^2}{Z}, \quad a_0 = \frac{4\pi\varepsilon_0 \hbar^2}{m_e e^2} \approx 0.529 \, \text{Å}
  \]
  \item \textbf{Niveles de Energía Cuantizados:}
  \[
    E_n = -\frac{m_e Z^2 e^4}{32 \pi^2 \varepsilon_0^2 \hbar^2} \frac{1}{n^2} = -E_R \frac{Z^2}{n^2} = -13.6 \, \text{eV} \cdot \frac{Z^2}{n^2}
  \]
  \item \textbf{Fórmula de Rydberg para Transiciones Espectrales:}
  \[
    \frac{1}{\lambda} = R_\infty Z^2 \left( \frac{1}{n_1^2} - \frac{1}{n_2^2} \right), \quad R_\infty = \frac{m_e e^4}{8 \varepsilon_0^2 h^3 c} \approx 1.09737 \times 10^7 \, \text{m}^{-1}
  \]
\end{itemize}

\subsection*{Nivel Intermedio}
\addcontentsline{toc}{subsection}{Nivel Intermedio}

\subsubsection*{21.3 Ecuación de Schrödinger No Relativista}
\addcontentsline{toc}{subsubsection}{21.3 Ecuación de Schrödinger No Relativista}

\begin{itemize}
  \item \textbf{Dependiente del Tiempo (1D y 3D):}
  \[
    i\hbar \frac{\partial \Psi(\vec{r}, t)}{\partial t} = \hat{H}\Psi(\vec{r}, t) = \left[ -\frac{\hbar^2}{2m}\nabla^2 + V(\vec{r}, t) \right] \Psi(\vec{r}, t)
  \]
  \item \textbf{Independiente del Tiempo (Estados Estacionarios $\Psi(\vec{r}, t) = \psi(\vec{r})e^{-iEt/\hbar}$):}
  \[
    \hat{H}\psi(\vec{r}) = E\psi(\vec{r}) \implies -\frac{\hbar^2}{2m}\nabla^2\psi(\vec{r}) + V(\vec{r})\psi(\vec{r}) = E\psi(\vec{r})
  \]
  \item \textbf{Densidad de Probabilidad ($P$) y Densidad de Corriente de Probabilidad ($\vec{J}$):}
  \[
    P(\vec{r}, t) = |\Psi(\vec{r}, t)|^2 = \Psi^* \Psi, \quad \int_{\text{todo el espacio}} |\Psi|^2 dV = 1
  \]
  \[
    \vec{J} = \frac{\hbar}{2mi}(\Psi^* \nabla\Psi - \Psi \nabla\Psi^*) \implies \frac{\partial P}{\partial t} + \nabla \cdot \vec{J} = 0
  \]
\end{itemize}

\subsubsection*{21.4 Problemas Unidimensionales Notables}
\addcontentsline{toc}{subsubsection}{21.4 Problemas Unidimensionales Notables}

\begin{itemize}
  \item \textbf{Pozo de Potencial Infinito (Caja 1D de longitud $L$ en $[0, L]$):}
  \[
    \psi_n(x) = \sqrt{\frac{2}{L}}\sen\left( \frac{n\pi x}{L} \right), \quad E_n = \frac{n^2 \pi^2 \hbar^2}{2m L^2} = \frac{n^2 h^2}{8m L^2}, \quad n = 1, 2, 3, \dots
  \]
  \item \textbf{Oscilador Armónico Cuántico Simple ($\omega = \sqrt{k/m}$):}
  \[
    E_n = \hbar\omega \left( n + \frac{1}{2} \right), \quad n = 0, 1, 2, \dots
  \]
  \[
    \text{Operadores de Aniquilación/Creación (Escalera): } \hat{a} = \sqrt{\frac{m\omega}{2\hbar}}\left(\hat{x} + \frac{i\hat{p}}{m\omega}\right), \quad \hat{a}^\dagger = \sqrt{\frac{m\omega}{2\hbar}}\left(\hat{x} - \frac{i\hat{p}}{m\omega}\right)
  \]
  \[
    [\hat{a}, \hat{a}^\dagger] = 1, \quad \hat{H} = \hbar\omega\left(\hat{a}^\dagger \hat{a} + \frac{1}{2}\right) = \hbar\omega\left(\hat{N} + \frac{1}{2}\right)
  \]
  \item \textbf{Efecto Túnel a través de Barrera Rectangular ($V_0 > E$, ancho $a$):}
  \[
    T \approx \exp(-2\kappa a), \quad \kappa = \frac{\sqrt{2m(V_0 - E)}}{\hbar} \quad (\text{para } \kappa a \gg 1)
  \]
\end{itemize}

\subsubsection*{21.5 Momento Angular Cuántico y Átomo de Hidrógeno}
\addcontentsline{toc}{subsubsection}{21.5 Momento Angular Cuántico y Átomo de Hidrógeno}

\begin{itemize}
  \item \textbf{Operadores de Momento Angular Orbital ($\hat{\vec{L}} = \hat{\vec{r}} \times \hat{\vec{p}}$):}
  \[
    [\hat{L}_x, \hat{L}_y] = i\hbar \hat{L}_z, \quad [\hat{L}_y, \hat{L}_z] = i\hbar \hat{L}_x, \quad [\hat{L}_z, \hat{L}_x] = i\hbar \hat{L}_y, \quad [\hat{L}^2, \hat{L}_z] = 0
  \]
  \item \textbf{Autovalores de Momento Angular y Proyección:}
  \[
    \hat{L}^2 |l, m_l\rangle = \hbar^2 l(l + 1)|l, m_l\rangle, \quad l = 0, 1, 2, \dots
  \]
  \[
    \hat{L}_z |l, m_l\rangle = \hbar m_l |l, m_l\rangle, \quad m_l = -l, -l+1, \dots, l
  \]
  \item \textbf{Función de Onda del Átomo de Hidrógeno:}
  \[
    \psi_{n, l, m_l}(r, \theta, \phi) = R_{nl}(r) Y_l^{m_l}(\theta, \phi)
  \]
  \[
    (R_{nl} = \text{Polinomios asociados de Laguerre}, \, Y_l^{m_l} = \text{Armónicos esféricos})
  \]
\end{itemize}

\subsection*{Nivel Universitario / Avanzado}
\addcontentsline{toc}{subsection}{Nivel Universitario / Avanzado}

\subsubsection*{21.6 Formalismo de Dirac (Bra-Ket), Postulados y Espacio de Hilbert}
\addcontentsline{toc}{subsubsection}{21.6 Formalismo de Dirac (Bra-Ket), Postulados y Espacio de Hilbert}

\begin{itemize}
  \item \textbf{Valor Esperado de un Observable $\hat{A}$ en el Estado $|\psi\rangle$:}
  \[
    \langle A \rangle = \frac{\langle\psi|\hat{A}|\psi\rangle}{\langle\psi|\psi\rangle}
  \]
  \item \textbf{Relación Generalizada de Incertidumbre de Robertson-Schrödinger:}
  \[
    \sigma_A \sigma_B \ge \frac{1}{2} |\langle [\hat{A}, \hat{B}] \rangle|
  \]
  \item \textbf{Relación de Conmutación Canónica:}
  \[
    [\hat{x}_j, \hat{p}_k] = i\hbar \delta_{jk}
  \]
  \item \textbf{Ecuación de Movimiento de Heisenberg para un Operador $\hat{A}$:}
  \[
    \frac{d\hat{A}_H}{dt} = \frac{i}{\hbar}[\hat{H}, \hat{A}_H] + \left( \frac{\partial\hat{A}}{\partial t} \right)_H
  \]
\end{itemize}

\subsubsection*{21.7 Espín $1/2$ y Matrices de Pauli ($\hat{\vec{S}} = \frac{\hbar}{2}\vec{\sigma}$)}
\addcontentsline{toc}{subsubsection}{21.7 Espín $1/2$ y Matrices de Pauli ($\hat{\vec{S}} = \frac{\hbar}{2}\vec{\sigma}$)}

\begin{itemize}
  \item \textbf{Matrices de Pauli:}
  \[
    \sigma_x = \begin{bmatrix} 0 & 1 \\ 1 & 0 \end{bmatrix}, \quad \sigma_y = \begin{bmatrix} 0 & -i \\ i & 0 \end{bmatrix}, \quad \sigma_z = \begin{bmatrix} 1 & 0 \\ 0 & -1 \end{bmatrix}
  \]
  \item \textbf{Álgebra de Pauli:}
  \[
    \sigma_i \sigma_j = \delta_{ij} I + i \sum_k \epsilon_{ijk} \sigma_k, \quad [\sigma_i, \sigma_j] = 2i \epsilon_{ijk} \sigma_k, \quad \{\sigma_i, \sigma_j\} = 2\delta_{ij} I
  \]
  \item \textbf{Momento Magnético de Espín:}
  \[
    \vec{\mu}_s = -g_s \frac{e}{2m_e}\vec{S} \approx -\frac{e}{m_e}\vec{S} = -\mu_B \vec{\sigma}, \quad \mu_B = \frac{e\hbar}{2m_e} \quad (\text{Magnetón de Bohr})
  \]
\end{itemize}

\subsubsection*{21.8 Teoría de Perturbaciones y Métodos de Aproximación}
\addcontentsline{toc}{subsubsection}{21.8 Teoría de Perturbaciones y Métodos de Aproximación}

\begin{itemize}
  \item \textbf{Perturbaciones Estacionarias No Degeneradas (1er y 2do orden):}
  \[
    E_n^{(1)} = \langle n^{(0)} | \hat{H}' | n^{(0)} \rangle
  \]
  \[
    E_n^{(2)} = \sum_{k \neq n} \frac{|\langle k^{(0)} | \hat{H}' | n^{(0)} \rangle|^2}{E_n^{(0)} - E_k^{(0)}}
  \]
  \[
    |n^{(1)}\rangle = \sum_{k \neq n} \frac{\langle k^{(0)} | \hat{H}' | n^{(0)} \rangle}{E_n^{(0)} - E_k^{(0)}} |k^{(0)}\rangle
  \]
  \item \textbf{Regla de Oro de Fermi (Tasa de transición dependiente del tiempo):}
  \[
    \Gamma_{i \to f} = W_{i \to f} = \frac{2\pi}{\hbar} |\langle f | \hat{H}' | i \rangle|^2 \rho(E_f)
  \]
\end{itemize}

\subsubsection*{21.9 Mecánica Cuántica Relativista}
\addcontentsline{toc}{subsubsection}{21.9 Mecánica Cuántica Relativista}

\begin{itemize}
  \item \textbf{Ecuación de Klein-Gordon (Partículas escalares de espín 0 con $\Box = \nabla^2 - \frac{1}{c^2}\frac{\partial^2}{\partial t^2}$):}
  \[
    \left( \Box - \frac{m^2 c^2}{\hbar^2} \right) \phi = 0 \iff \nabla^2\phi - \frac{1}{c^2}\frac{\partial^2\phi}{\partial t^2} = \left(\frac{mc}{\hbar}\right)^2 \phi
  \]
  \[
    \left( \text{O con firma } \Box = \frac{1}{c^2}\frac{\partial^2}{\partial t^2} - \nabla^2 = \partial_\mu \partial^\mu: \quad \left(\Box + \frac{m^2 c^2}{\hbar^2}\right)\phi = 0 \right)
  \]
  \item \textbf{Ecuación de Dirac (Fermiones de espín 1/2 en representación covariante):}
  \[
    (i\gamma^\mu \partial_\mu - \frac{mc}{\hbar})\psi = 0 \iff (i\hbar \gamma^\mu \partial_\mu - mc)\psi = 0
  \]
  \item \textbf{Álgebra de Clifford para las Matrices Gamma de Dirac ($\{\gamma^\mu, \gamma^\nu\} = 2\eta^{\mu\nu} I_{4\times4}$):}
  \[
    \gamma^0 = \begin{bmatrix} I & 0 \\ 0 & -I \end{bmatrix}, \quad \gamma^k = \begin{bmatrix} 0 & \sigma_k \\ -\sigma_k & 0 \end{bmatrix}, \quad \gamma^5 = i\gamma^0 \gamma^1 \gamma^2 \gamma^3 = \begin{bmatrix} 0 & I \\ I & 0 \end{bmatrix}
  \]
\end{itemize}

\section*{22. Física Nuclear, Radiactividad y Partículas Elementales}
\addcontentsline{toc}{section}{22. Física Nuclear, Radiactividad y Partículas Elementales}

\subsection*{Nivel Básico}
\addcontentsline{toc}{subsection}{Nivel Básico}

\subsubsection*{22.1 Estructura Nuclear y Defecto de Masa}
\addcontentsline{toc}{subsubsection}{22.1 Estructura Nuclear y Defecto de Masa}

\begin{itemize}
  \item \textbf{Radio Nuclear Empírico ($A = \text{número másico}$):}
  \[
    R \approx R_0 A^{1/3}, \quad R_0 \approx 1.2 \, \text{fm} = 1.2 \times 10^{-15} \, \text{m}
  \]
  \item \textbf{Defecto de Masa ($\Delta m$ para un núcleo $_Z^A X_N$ con $N = A - Z$):}
  \[
    \Delta m = Z m_p + N m_n - M_{\text{núcleo}}(A, Z)
  \]
  \item \textbf{Energía de Enlace Nuclear Total ($B$ o $E_b$):}
  \[
    B(A, Z) = \Delta m \cdot c^2 = [Z m_p + (A - Z) m_n - M_{\text{núcleo}}] c^2
  \]
  \item \textbf{Energía de Enlace por Nucleón ($B/A$):}
  \[
    \frac{B}{A} = \frac{B(A, Z)}{A} \quad (\approx 8.8 \, \text{MeV/nucleón para el } ^{56}\text{Fe})
  \]
\end{itemize}

\subsubsection*{22.2 Cinética de la Desintegración Radiactiva}
\addcontentsline{toc}{subsubsection}{22.2 Cinética de la Desintegración Radiactiva}

\begin{itemize}
  \item \textbf{Ley de Desintegración Radiactiva:}
  \[
    N(t) = N_0 e^{-\lambda t}
  \]
  \item \textbf{Actividad Radiactiva ($A(t)$ en Becquerels $\text{Bq} = \text{des/s}$ o Curie $\text{Ci}$):}
  \[
    A(t) = -\frac{dN}{dt} = \lambda N(t) = A_0 e^{-\lambda t}, \quad A_0 = \lambda N_0
  \]
  \item \textbf{Periodo de Semidesintegración / Vida Media ($t_{1/2}$):}
  \[
    t_{1/2} = \frac{\ln(2)}{\lambda} \approx \frac{0.693}{\lambda}
  \]
  \item \textbf{Tiempo de Vida Medio ($\tau$):}
  \[
    \tau = \frac{1}{\lambda} = \frac{t_{1/2}}{\ln(2)} \approx 1.443 \, t_{1/2}
  \]
\end{itemize}

\subsubsection*{22.3 Tipos Principales de Desintegración Radiactiva}
\addcontentsline{toc}{subsubsection}{22.3 Tipos Principales de Desintegración Radiactiva}

\begin{itemize}
  \item \textbf{Desintegración Alfa ($\alpha = \,_2^4\text{He}$):}
  \[
    _Z^A X \longrightarrow \,_{Z-2}^{A-4} Y + \,_2^4\text{He} + Q_\alpha
  \]
  \item \textbf{Desintegración Beta Negativa ($\beta^-$ con emisión de antineutrino electrónico):}
  \[
    _Z^A X \longrightarrow \,_{Z+1}^A Y + e^- + \bar{\nu}_e + Q_{\beta^-} \quad (n \to p + e^- + \bar{\nu}_e)
  \]
  \item \textbf{Desintegración Beta Positiva ($\beta^+$ con emisión de neutrino electrónico):}
  \[
    _Z^A X \longrightarrow \,_{Z-1}^A Y + e^+ + \nu_e + Q_{\beta^+} \quad (p \to n + e^+ + \nu_e)
  \]
  \item \textbf{Captura Electrónica ($\text{CE}$):}
  \[
    _Z^A X + e^- \longrightarrow \,_{Z-1}^A Y + \nu_e + Q_{\text{CE}}
  \]
  \item \textbf{Emisión Gamma ($\gamma = \text{fotón nuclear}$):}
  \[
    _Z^A X^* \longrightarrow \,_Z^A X + \gamma
  \]
\end{itemize}

\subsection*{Nivel Intermedio}
\addcontentsline{toc}{subsection}{Nivel Intermedio}

\subsubsection*{22.4 Balances Energéticos y Modelos Nucleares}
\addcontentsline{toc}{subsubsection}{22.4 Balances Energéticos y Modelos Nucleares}

\begin{itemize}
  \item \textbf{Valor $Q$ de una Reacción Nuclear ($a + X \to Y + b$):}
  \[
    Q = (m_a + m_X - m_Y - m_b) c^2 = E_{k,Y} + E_{k,b} - E_{k,a}
  \]
  \[
    Q > 0 \quad (\text{Exoenergética / Exotérmica}), \quad Q < 0 \quad (\text{Endoenergética / Endotérmica})
  \]
  \item \textbf{Energía Umbral para Reacciones Endoenergéticas ($Q < 0$ sobre blanco estacionario):}
  \[
    E_{\text{umbral}} = |Q|\left( 1 + \frac{m_a}{m_X} \right)
  \]
  \item \textbf{Fórmula Semiempírica de Masas de Bethe-Weizsäcker (Modelo de la Gota Líquida):}
  \[
    B(A, Z) = a_v A - a_s A^{2/3} - a_c \frac{Z(Z - 1)}{A^{1/3}} - a_a \frac{(A - 2Z)^2}{A} + \delta(A, Z)
  \]
  \[
    \delta(A, Z) = \begin{cases} +a_p A^{-1/2} & \text{para } Z \text{ par, } N \text{ par} \\ 0 & \text{para } A \text{ impar (par-impar / impar-par)} \\ -a_p A^{-1/2} & \text{para } Z \text{ impar, } N \text{ impar} \end{cases} \quad (\text{o con término } \delta \propto \pm a_p A^{-3/4})
  \]
  \item \textbf{Ley de Geiger-Nuttall para Desintegración Alfa ($E_\alpha = \text{energía de la partícula } \alpha$):}
  \[
    \log_{10}(\lambda) = A_G + B_G \frac{Z}{\sqrt{E_\alpha}}
  \]
\end{itemize}

\subsubsection*{22.5 Secciones Eficaces y Dosimetría}
\addcontentsline{toc}{subsubsection}{22.5 Secciones Eficaces y Dosimetría}

\begin{itemize}
  \item \textbf{Tasa de Reacción Nuclear (Flujo $\Phi$, densidad de núcleos blanco $n_t$, sección eficaz $\sigma$):}
  \[
    R = \Phi \sigma N_t = I n_t x \sigma, \quad 1 \, \text{barn} (\text{b}) = 10^{-28} \, \text{m}^2 = 100 \, \text{fm}^2
  \]
  \item \textbf{Dosis Absorbida ($D$ en Grays $\text{Gy} = \text{J/kg}$) y Dosis Equivalente ($H$ en Sieverts $\text{Sv}$ con factor de calidad $w_R$):}
  \[
    D = \frac{dE_{\text{dep}}}{dm}, \quad H = w_R \cdot D
  \]
\end{itemize}

\subsection*{Nivel Universitario / Avanzado}
\addcontentsline{toc}{subsection}{Nivel Universitario / Avanzado}

\subsubsection*{22.6 Física de Partículas y Modelo Estándar}
\addcontentsline{toc}{subsubsection}{22.6 Física de Partículas y Modelo Estándar}

\begin{itemize}
  \item \textbf{Potencial Nuclear Fuerte de Yukawa (Mediado por mesones de masa $m_\pi$):}
  \[
    V(r) = -g^2 \frac{e^{-\mu r}}{r} = -g^2 \frac{e^{-r / \lambda_C}}{r}, \quad \mu = \frac{m_\pi c}{\hbar}
  \]
  \item \textbf{Fórmula de Gell-Mann-Nishijima (Carga, Isoespín, Hipercarga y Número Bariónico):}
  \[
    Q = I_3 + \frac{Y}{2} = I_3 + \frac{B + S + C + B' + T}{2}
  \]
  \[
    (\text{donde } B = \text{bariónico}, \, S = \text{extrañeza}, \, C = \text{encanto}, \, B' = \text{bottomness/belleza}, \, T = \text{topness/verdad})
  \]
  \item \textbf{Matriz CKM (Cabibbo-Kobayashi-Maskawa para mezcla de quarks):}
  \[
    \begin{bmatrix} d' \\ s' \\ b' \end{bmatrix} = \mathbf{V}_{\text{CKM}} \begin{bmatrix} d \\ s \\ b \end{bmatrix} = \begin{bmatrix} V_{ud} & V_{us} & V_{ub} \\ V_{cd} & V_{cs} & V_{cb} \\ V_{td} & V_{ts} & V_{tb} \end{bmatrix} \begin{bmatrix} d \\ s \\ b \end{bmatrix}
  \]
  \item \textbf{Relación de Dispersión de Breit-Wigner (Resonancia de desintegración con anchura $\Gamma$):}
  \[
    \sigma(E) = \frac{\pi}{k^2} \frac{g \Gamma_{\text{in}} \Gamma_{\text{out}}}{(E - E_0)^2 + (\Gamma/2)^2}, \quad \tau = \frac{\hbar}{\Gamma}
  \]
\end{itemize}

\section*{23. Física del Estado Sólido y Materia Condensada}
\addcontentsline{toc}{section}{23. Física del Estado Sólido y Materia Condensada}

\subsection*{Nivel Básico / Intermedio}
\addcontentsline{toc}{subsection}{Nivel Básico / Intermedio}

\subsubsection*{23.1 Cristalografía y Difracción}
\addcontentsline{toc}{subsubsection}{23.1 Cristalografía y Difracción}

\begin{itemize}
  \item \textbf{Ley de Bragg para Difracción de Rayos X ($d_{hkl} = \text{espaciado interplanar}$):}
  \[
    2 d_{hkl} \sen\theta = n \lambda, \quad n \in \mathbb{N}
  \]
  \item \textbf{Espaciado Interplanar para Red Cúbica Simple (Parámetro de red $a$):}
  \[
    d_{hkl} = \frac{a}{\sqrt{h^2 + k^2 + l^2}}
  \]
  \item \textbf{Vectores Primitivos de la Red Recíproca:}
  \[
    \vec{b}_1 = 2\pi \frac{\vec{a}_2 \times \vec{a}_3}{\vec{a}_1 \cdot (\vec{a}_2 \times \vec{a}_3)}, \quad \vec{b}_2 = 2\pi \frac{\vec{a}_3 \times \vec{a}_1}{\vec{a}_1 \cdot (\vec{a}_2 \times \vec{a}_3)}, \quad \vec{b}_3 = 2\pi \frac{\vec{a}_1 \times \vec{a}_2}{\vec{a}_1 \cdot (\vec{a}_2 \times \vec{a}_3)}
  \]
  \item \textbf{Condición de Difracción de Laue ($\vec{G} = \vec{G}_{hkl} = \text{vector de red recíproca}$):}
  \[
    \Delta\vec{k} = \vec{k}' - \vec{k} = \vec{G} \iff 2\vec{k} \cdot \vec{G} + |\vec{G}|^2 = 0 \quad (\text{para dispersión elástica } |\vec{k}'| = |\vec{k}|)
  \]
  \[
    \left( \text{O definiendo } \Delta\vec{k} = \vec{k} - \vec{k}' = \vec{G}: \quad 2\vec{k} \cdot \vec{G} = |\vec{G}|^2 \right)
  \]
\end{itemize}

\subsubsection*{23.2 Gas de Electrones Libres de Fermi}
\addcontentsline{toc}{subsubsection}{23.2 Gas de Electrones Libres de Fermi}

\begin{itemize}
  \item \textbf{Vector de Onda de Fermi y Momento de Fermi:}
  \[
    k_F = (3\pi^2 n)^{1/3}, \quad p_F = \hbar k_F = \hbar (3\pi^2 n)^{1/3}, \quad n = \frac{N}{V}
  \]
  \item \textbf{Energía de Fermi a $T = 0\text{ K}$:}
  \[
    E_F = \frac{\hbar^2 k_F^2}{2m} = \frac{\hbar^2}{2m}(3\pi^2 n)^{2/3}
  \]
  \item \textbf{Temperatura de Fermi y Velocidad de Fermi:}
  \[
    T_F = \frac{E_F}{k_B}, \quad v_F = \frac{\hbar k_F}{m} = \sqrt{\frac{2E_F}{m}}
  \]
  \item \textbf{Densidad de Estados Tridimensional ($g(E)$ o $D(E)$):}
  \[
    g(E) = \frac{V}{2\pi^2}\left( \frac{2m}{\hbar^2} \right)^{3/2} \sqrt{E} = \frac{3N}{2E_F^{3/2}}\sqrt{E}
  \]
  \item \textbf{Capacidad Calorífica Electrónica Lineal ($T \ll T_F$):}
  \[
    C_{v,e} = \gamma T = \frac{\pi^2}{2}\left( \frac{k_B T}{E_F} \right) N k_B = \frac{\pi^2}{3} g(E_F) k_B^2 T
  \]
\end{itemize}

\subsubsection*{23.3 Transporte Electrónico y Térmico}
\addcontentsline{toc}{subsubsection}{23.3 Transporte Electrónico y Térmico}

\begin{itemize}
  \item \textbf{Conductividad Eléctrica de Drude-Sommerfeld:}
  \[
    \sigma = \frac{n e^2 \tau_c}{m}
  \]
  \item \textbf{Ley de Wiedemann-Franz (Número de Lorenz $L$):}
  \[
    \frac{K}{\sigma} = L T, \quad L = \frac{\pi^2}{3}\left( \frac{k_B}{e} \right)^2 \approx 2.44 \times 10^{-8} \, \frac{\text{W}\cdot\Omega}{\text{K}^2}
  \]
\end{itemize}

\subsection*{Nivel Universitario / Avanzado}
\addcontentsline{toc}{subsection}{Nivel Universitario / Avanzado}

\subsubsection*{23.4 Teoría de Bandas y Semiconductores}
\addcontentsline{toc}{subsubsection}{23.4 Teoría de Bandas y Semiconductores}

\begin{itemize}
  \item \textbf{Teorema de Bloch para Electrones en un Potencial Periódico ($V(\vec{r} + \vec{R}) = V(\vec{r})$):}
  \[
    \psi_{\vec{k}}(\vec{r}) = e^{i\vec{k} \cdot \vec{r}} u_{\vec{k}}(\vec{r}), \quad u_{\vec{k}}(\vec{r} + \vec{R}) = u_{\vec{k}}(\vec{r})
  \]
  \item \textbf{Tensor de Masa Efectiva del Portador:}
  \[
    \left( \frac{1}{m^*} \right)_{ij} = \frac{1}{\hbar^2} \frac{\partial^2 E(\vec{k})}{\partial k_i \partial k_j}
  \]
  \item \textbf{Concentración Intrínseca de Portadores en Semiconductores ($E_g = \text{Banda prohibida / Bandgap}$):}
  \[
    n_i = \sqrt{N_c N_v} \exp\left( -\frac{E_g}{2 k_B T} \right)
  \]
  \[
    N_c = 2\left( \frac{m_e^* k_B T}{2\pi\hbar^2} \right)^{3/2}, \quad N_v = 2\left( \frac{m_h^* k_B T}{2\pi\hbar^2} \right)^{3/2}
  \]
  \item \textbf{Ley de Acción de Masas en Semiconductores:}
  \[
    n \cdot p = n_i^2
  \]
  \item \textbf{Nivel de Fermi Intrínseco:}
  \[
    E_{Fi} = \frac{E_c + E_v}{2} + \frac{3}{4} k_B T \ln\left( \frac{m_h^*}{m_e^*} \right)
  \]
\end{itemize}

\subsubsection*{23.5 Fonones y Propiedades Térmicas de la Red}
\addcontentsline{toc}{subsubsection}{23.5 Fonones y Propiedades Térmicas de la Red}

\begin{itemize}
  \item \textbf{Modelo de Debye para la Capacidad Calorífica de la Red ($T \ll \Theta_D$):}
  \[
    C_v = \frac{12\pi^4}{5} N k_B \left( \frac{T}{\Theta_D} \right)^3 \propto T^3, \quad \Theta_D = \frac{\hbar v_s}{k_B}(6\pi^2 n)^{1/3}
  \]
  \item \textbf{Modelo de Einstein para la Capacidad Calorífica:}
  \[
    C_v = 3 N k_B \left( \frac{\Theta_E}{T} \right)^2 \frac{e^{\Theta_E / T}}{(e^{\Theta_E / T} - 1)^2}, \quad \Theta_E = \frac{\hbar\omega_E}{k_B}
  \]
\end{itemize}

\subsubsection*{23.6 Superconductividad (Teoría Clásica de London y BCS)}
\addcontentsline{toc}{subsubsection}{23.6 Superconductividad (Teoría Clásica de London y BCS)}

\begin{itemize}
  \item \textbf{Ecuaciones de London ($\lambda_L = \text{longitud de penetración de London}$):}
  \[
    \frac{\partial \vec{J}_s}{\partial t} = \frac{n_s e^2}{m} \vec{E}, \quad \nabla \times \vec{J}_s = -\frac{n_s e^2}{m}\vec{B} \implies \nabla^2\vec{B} = \frac{1}{\lambda_L^2}\vec{B}
  \]
  \item \textbf{Efecto Meissner (Atenuación exponencial del campo magnético en el superconductor):}
  \[
    B(x) = B_0 \exp\left( -\frac{x}{\lambda_L} \right), \quad \lambda_L = \sqrt{\frac{m}{\mu_0 n_s e^2}}
  \]
  \item \textbf{Campo Magnético Crítico Termodinámico:}
  \[
    B_c(T) = B_c(0)\left[ 1 - \left( \frac{T}{T_c} \right)^2 \right]
  \]
  \item \textbf{Brecha de Energía Superconductora BCS a $T = 0\text{ K}$:}
  \[
    \Delta(0) \approx 1.764 \, k_B T_c
  \]
  \item \textbf{Longitud de Coherencia de Cooper / BCS ($\xi_0$):}
  \[
    \xi_0 = \frac{\hbar v_F}{\pi \Delta(0)}
  \]
\end{itemize}


\end{document}
